% Dipôles commandés, capteurs et chaînes électroniques — Physique Tle D — Cours signé OnBuch+
% Programme officiel MINESEC. Compiler avec : tectonic P13-dipoles-commandes-capteurs-chaines-electroniques.tex
\newcommand{\DOCMATIERE}{Physique}
\newcommand{\DOCNIVEAU}{Tle D}
\newcommand{\DOCMODULE}{Module 3 : Électricité — évolutions temporelles des circuits électriques et électroniques}
\newcommand{\DOCLECON}{P13}
\newcommand{\DOCTITRE}{Dipôles commandés, capteurs et chaînes électroniques}
% OnBuch+ — préambule « cours signés » ENRICHI (compilé avec tectonic / XeLaTeX)
% Système visuel « Pop » d'OnBuch+ : fond crème (jamais blanc), bordures encre
% épaisses, ombres « dures » décalées, titres Archivo Black en capitales.
% Version MATHÉMATIQUES : graphes (pgfplots/tikz), boîtes pédagogiques
% (définition, propriété, méthode, attention, le savais-tu, exercice résolu,
% exercice, TP, bilan), page de garde et signature OnBuch+ ancrée en bas.
%
% Macros attendues dans main.tex AVANT \input{preamble} :
%   \DOCMATIERE  \DOCNIVEAU  \DOCMODULE  \DOCLECON (numéro)  \DOCTITRE
\documentclass[11pt]{article}

\usepackage{fontspec}
\usepackage[french]{babel}
\usepackage[a4paper,margin=2cm,top=3.4cm,bottom=2.8cm,headheight=1.4cm,footskip=1.3cm]{geometry}
\usepackage{amsmath,amssymb,amsfonts}
\usepackage{mathtools} % \xmapsto, \coloneqq... souvent utilisés par les modèles
\usepackage{cancel} % simplifications barrées (bilans de forces)
% \abs{z} (module, valeur absolue) et \norm{u} (norme), utilisés par les modèles
\providecommand{\abs}{}\let\abs\relax\DeclarePairedDelimiter{\abs}{\lvert}{\rvert}
\providecommand{\norm}{}\let\norm\relax\DeclarePairedDelimiter{\norm}{\lVert}{\rVert}
\usepackage[table]{xcolor} % [table] : fournit aussi \rowcolors (alternance de lignes)
\usepackage{pagecolor}
\usepackage{tikz}
\usetikzlibrary{arrows.meta,positioning,calc,shapes.geometric,shapes.misc,shapes.arrows,decorations.pathmorphing,decorations.markings,patterns,shadows,mindmap,trees,fit,backgrounds,matrix,intersections,angles,quotes}
\usepackage{pgfplots}
\usepackage[version=4]{mhchem} % équations nucléaires \ce{^{235}_{92}U}
\usepackage[european,siunitx]{circuitikz} % schémas électriques
\pgfplotsset{compat=1.18}
\usepgfplotslibrary{fillbetween} % aires entre courbes (calcul intégral)
\usepackage{enumitem}
\usepackage{fancyhdr}
% [most] charge toutes les bibliothèques tcolorbox (listings, minted, raster,
% documentation...) et gonfle inutilement la mémoire TeX sur un document long
% (~300 boîtes) : on ne charge que ce qui est réellement utilisé.
\usepackage[skins,breakable]{tcolorbox}
\usepackage{graphicx}
\usepackage{colortbl}
\usepackage{booktabs}
\usepackage{tabularx}
\usepackage{diagbox} % case d'en-tête diagonale des tableaux à double entrée
% Colonnes extensibles centrée / gauche / droite (C, L, R), souvent utilisées par les modèles
\newcolumntype{C}{>{\centering\arraybackslash}X}
\newcolumntype{L}{>{\raggedright\arraybackslash}X}
\newcolumntype{R}{>{\raggedleft\arraybackslash}X}
\usepackage{array}
\usepackage{multicol}
\usepackage{siunitx}
% parse-numbers=false : accepte aussi \SI{2,5 \times 10^{-3}}{...}
\sisetup{locale=FR,output-decimal-marker={,},group-minimum-digits=4,parse-numbers=false}
\DeclareSIUnit\ohm{\ensuremath{\symup{\Omega}}} % U+2126 absent de la police
\DeclareSIUnit\division{div} % réglages d'oscilloscope (ms/div, V/div)
\usepackage{qrcode}
% \degree : commande fréquemment utilisée par les modèles pour un angle ou une
% température en dehors de \SI{}{} (non fournie par les paquets chargés).
\providecommand{\degree}{^{\circ}}
% \qed (fin de démonstration), utilisé par les modèles sans amsthm
\providecommand{\qed}{\hfill$\square$}
% \barre : double barre des valeurs interdites dans un tableau de variations (tkz-tab)
\providecommand{\barre}{\ensuremath{\|}}
% environnement proof (amsthm non chargé) utilisé par les modèles
\ifdefined\proof\else\newenvironment{proof}[1][Démonstration]{\par\noindent\textit{#1.}\ }{\qed\par}\fi
\usepackage{lastpage}
\usepackage{hyperref}
\hypersetup{hidelinks}

% ============================================================
% Palette « Pop » OnBuch+ — mêmes hex que la classe P (plus_ui.dart)
% ============================================================
\definecolor{poporange}{HTML}{F59321}
\definecolor{popdark}{HTML}{E07A0C}
\definecolor{popcream}{HTML}{FFF6E8}
\definecolor{popink}{HTML}{1C1714}
\definecolor{poppurple}{HTML}{7A5AE0}
\definecolor{popgreen}{HTML}{2E9E6A}
\definecolor{poppink}{HTML}{E0457B}
\definecolor{popblue}{HTML}{2F7DE1}
\definecolor{popgold}{HTML}{C98A12}
\definecolor{popmuted}{HTML}{8A7F76}
\definecolor{popline}{HTML}{EADFD2}
\definecolor{popgray}{HTML}{8A7F76}
\colorlet{popgrey}{popgray}
% Alias tolérants (le modèle nomme parfois les couleurs différemment)
\colorlet{popwhite}{white}
\colorlet{popblack}{popink}
\colorlet{popred}{poppink}
\colorlet{popyellow}{popgold}
% Teintes claires pour fonds de tableaux / zones
\colorlet{popblueL}{popblue!10!white}
\colorlet{poporangeL}{poporange!14!white}
\colorlet{popgreenL}{popgreen!12!white}
\colorlet{poppurpleL}{poppurple!12!white}
\colorlet{poppinkL}{poppink!12!white}
\colorlet{popgoldL}{popgold!14!white}

\pagecolor{popcream}

% ============================================================
% Typographie
% ============================================================
\defaultfontfeatures{Ligatures=TeX}
\setmainfont{PlusJakartaSans-Medium.ttf}[Path=../../../fonts/,BoldFont=PlusJakartaSans-Bold.ttf]
% Maths en sans-serif (Fira Math) avec chiffres et lettres droites de Plus
% Jakarta, pour que les formules se fondent dans le texte.
\usepackage[math-style=ISO,bold-style=ISO]{unicode-math}
\setmathfont{FiraMath-Regular.otf}[Path=../../../fonts/]
\setmathfont{PlusJakartaSans-Medium.ttf}[Path=../../../fonts/,range={up/num,up/{Latin,latin},\mathup}]
\usepackage{icomma} % virgule décimale sans espace en mode math
% Caractères mathématiques Unicode absents des polices de texte (Plus Jakarta,
% Archivo Black) : rendus par leur équivalent mathématique, en texte comme en
% formule — sinon ils s'affichent en carré vide (ex. « Arithmétique dans ℤ »).
\usepackage{newunicodechar}
\newunicodechar{ℝ}{\ensuremath{\mathbb{R}}}
\newunicodechar{ℤ}{\ensuremath{\mathbb{Z}}}
\newunicodechar{ℂ}{\ensuremath{\mathbb{C}}}
\newunicodechar{ℕ}{\ensuremath{\mathbb{N}}}
\newunicodechar{ℚ}{\ensuremath{\mathbb{Q}}}
\newunicodechar{𝔻}{\ensuremath{\mathbb{D}}}
\newunicodechar{α}{\ensuremath{\alpha}}
\newunicodechar{β}{\ensuremath{\beta}}
\newunicodechar{γ}{\ensuremath{\gamma}}
\newunicodechar{δ}{\ensuremath{\delta}}
\newunicodechar{ε}{\ensuremath{\varepsilon}}
\newunicodechar{θ}{\ensuremath{\theta}}
\newunicodechar{λ}{\ensuremath{\lambda}}
\newunicodechar{μ}{\ensuremath{\mu}}
\newunicodechar{σ}{\ensuremath{\sigma}}
\newunicodechar{τ}{\ensuremath{\tau}}
\newunicodechar{φ}{\ensuremath{\varphi}}
\newunicodechar{ψ}{\ensuremath{\psi}}
\newunicodechar{ω}{\ensuremath{\omega}}
\newunicodechar{Σ}{\ensuremath{\Sigma}}
% majuscules produites par \MakeUppercase dans les titres (α-aminés -> Α)
\newunicodechar{Α}{\ensuremath{\alpha}}
\newunicodechar{Β}{\ensuremath{\beta}}
\newunicodechar{Γ}{\ensuremath{\Gamma}}
\newunicodechar{Δ}{\ensuremath{\Delta}}
\newunicodechar{Ε}{\ensuremath{\varepsilon}}
\newunicodechar{Θ}{\ensuremath{\Theta}}
\newunicodechar{Λ}{\ensuremath{\Lambda}}
\newunicodechar{Μ}{\ensuremath{\mu}}
\newunicodechar{Τ}{\ensuremath{\tau}}
\newunicodechar{Φ}{\ensuremath{\Phi}}
\newunicodechar{Ψ}{\ensuremath{\Psi}}
\newunicodechar{Ω}{\ensuremath{\Omega}}
\newunicodechar{↦}{\ensuremath{\mapsto}}
\newunicodechar{∈}{\ensuremath{\in}}
\newunicodechar{∉}{\ensuremath{\notin}}
\newunicodechar{∧}{\ensuremath{\wedge}}
\newunicodechar{⇔}{\ensuremath{\Leftrightarrow}}
\newunicodechar{⟺}{\ensuremath{\Longleftrightarrow}}
\newunicodechar{⇒}{\ensuremath{\Rightarrow}}
\newunicodechar{⟹}{\ensuremath{\Longrightarrow}}
\newunicodechar{∘}{\ensuremath{\circ}}
\newunicodechar{∪}{\ensuremath{\cup}}
\newunicodechar{∩}{\ensuremath{\cap}}
\newunicodechar{≡}{\ensuremath{\equiv}}
\newunicodechar{⊂}{\ensuremath{\subset}}
\newunicodechar{⊥}{\ensuremath{\perp}}
\newunicodechar{∃}{\ensuremath{\exists}}
\newunicodechar{∀}{\ensuremath{\forall}}
\newunicodechar{∛}{\ensuremath{\sqrt[3]{\,}}}
\newunicodechar{∜}{\ensuremath{\sqrt[4]{\,}}}
\newunicodechar{⃗}{\ensuremath{{}^{\to}}}
\newunicodechar{ⁿ}{\ifmmode^{n}\else\textsuperscript{\ensuremath{n}}\fi}
\newunicodechar{ˣ}{\ifmmode^{x}\else\textsuperscript{\ensuremath{x}}\fi}
\newunicodechar{ᵇ}{\ifmmode^{b}\else\textsuperscript{\ensuremath{b}}\fi}
\newunicodechar{ʳ}{\ifmmode^{r}\else\textsuperscript{\ensuremath{r}}\fi}
\newunicodechar{ᵘ}{\ifmmode^{u}\else\textsuperscript{\ensuremath{u}}\fi}
\newunicodechar{ʸ}{\ifmmode^{y}\else\textsuperscript{\ensuremath{y}}\fi}
\newunicodechar{ᵃ}{\ifmmode^{a}\else\textsuperscript{\ensuremath{a}}\fi}
\newunicodechar{ᵉ}{\ifmmode^{e}\else\textsuperscript{\ensuremath{e}}\fi}
\newunicodechar{ᵏ}{\ifmmode^{k}\else\textsuperscript{\ensuremath{k}}\fi}
\newunicodechar{ᵖ}{\ifmmode^{p}\else\textsuperscript{\ensuremath{p}}\fi}
\newunicodechar{ᴺ}{\ifmmode^{N}\else\textsuperscript{\ensuremath{N}}\fi}
\newunicodechar{ᶜ}{\ifmmode^{c}\else\textsuperscript{\ensuremath{c}}\fi}
\newunicodechar{ʰ}{\ifmmode^{h}\else\textsuperscript{\ensuremath{h}}\fi}
\newunicodechar{ᵠ}{\ifmmode^{q}\else\textsuperscript{\ensuremath{q}}\fi}
\newunicodechar{ₙ}{\ifmmode_{n}\else\textsubscript{\ensuremath{n}}\fi}
\newunicodechar{ₐ}{\ifmmode_{a}\else\textsubscript{\ensuremath{a}}\fi}
\newunicodechar{ᵢ}{\ifmmode_{i}\else\textsubscript{\ensuremath{i}}\fi}
\newunicodechar{ᵣ}{\ifmmode_{r}\else\textsubscript{\ensuremath{r}}\fi}
\newunicodechar{ₖ}{\ifmmode_{k}\else\textsubscript{\ensuremath{k}}\fi}
\newunicodechar{ᵦ}{\ifmmode_{\beta}\else\textsubscript{\ensuremath{\beta}}\fi}
\newunicodechar{ₓ}{\ifmmode_{x}\else\textsubscript{\ensuremath{x}}\fi}
\newfontfamily\popdisplay{ArchivoBlack-Regular.ttf}[Path=../../../fonts/]
\newfontfamily\poplogo{SpaceGrotesk-Bold.ttf}[Path=../../../fonts/]
\newfontfamily\popsemi{PlusJakartaSans-SemiBold.ttf}[Path=../../../fonts/]
\newfontfamily\popxbold{PlusJakartaSans-ExtraBold.ttf}[Path=../../../fonts/]
\newfontfamily\popxboldls{PlusJakartaSans-ExtraBold.ttf}[Path=../../../fonts/,LetterSpace=3.0]

\setlist[enumerate]{leftmargin=1.7em,itemsep=5pt,topsep=4pt}
\setlist[itemize]{leftmargin=1.7em,itemsep=4pt,topsep=4pt,label={\color{poporange}\textbullet}}
\setlength{\parindent}{0pt}
\setlength{\parskip}{5pt}
\renewcommand{\arraystretch}{1.25}

% pgfplots : style Pop par défaut
\pgfplotsset{
  every axis/.append style={axis line style={popink,line width=1pt},tick style={popink},
    grid style={popline,line width=0.5pt},label style={font=\small},tick label style={font=\footnotesize}},
  popcurve/.style={poporange,line width=1.8pt,no markers,smooth},
}
% Même style disponible dans un \draw TikZ simple (hors pgfplots)
\tikzset{popcurve/.style={poporange,line width=1.8pt}}
\tikzset{popgrid/.style={popline,very thin}}

% ============================================================
% Pill / squiggle
% ============================================================
\newcommand{\poppill}[3]{%
  \tikz[baseline=-0.55ex]\node[draw=popink,line width=1.2pt,rounded corners=2.6pt,%
    fill=#2,inner xsep=6.5pt,inner ysep=3pt]{%
    \popxboldls\fontsize{7}{7}\selectfont\color{#3}\MakeUppercase{#1}};%
}
\newcommand{\popsquiggle}[1]{%
  \tikz[baseline=-0.4ex]{
    \draw[#1,line width=2.6pt,line cap=round]
      plot [smooth] coordinates {(0,0.09) (0.34,0.01) (0.68,0.21) (1.02,0.01) (1.36,0.10)};
  }%
}
% Mot-clé surligné (marqueur orange sous le texte)
\newcommand{\cle}[1]{{\popxbold\color{popdark}#1}}

% ============================================================
% En-tête / pied
% ============================================================
\pagestyle{fancy}
\fancyhf{}
\renewcommand{\headrulewidth}{0pt}
\renewcommand{\footrulewidth}{0pt}
\lhead{\raisebox{-0.15ex}{\poplogo\fontsize{13}{13}\selectfont\color{popink}OnBuch\color{poporange}+}}
\rhead{\poppill{\DOCMATIERE\ \textbullet\ \DOCNIVEAU\ \textbullet\ Leçon \DOCLECON}{white}{popink}}
\lfoot{\popxbold\fontsize{7.6}{7.6}\selectfont\color{popmuted}\MakeUppercase{Cours signé OnBuch+}\ \textbullet\ {\popsemi\DOCTITRE}}
\rfoot{\tikz[baseline=-0.6ex]\node[rounded corners=8.5pt,fill=popink,minimum height=17pt,minimum width=17pt,inner xsep=5pt,inner sep=0]{\popxbold\fontsize{8}{8}\selectfont\color{popcream}\thepage\,/\,\pageref*{LastPage}};}

% ============================================================
% Titres
% ============================================================
\newcommand{\chaptitre}[2]{%
  \begin{tcolorbox}[enhanced,colback=poporange,colframe=popink,boxrule=2.6pt,arc=11pt,
      drop shadow=popink,left=15pt,right=15pt,top=12pt,bottom=11pt,before skip=3pt,after skip=18pt]
    \poppill{#2}{popink}{popcream}\par\vspace{9pt}
    {\raggedright\hyphenpenalty=10000\exhyphenpenalty=10000\popdisplay\fontsize{22}{24}\selectfont\color{popcream}\MakeUppercase{#1}\par}
  \end{tcolorbox}
}
% Section numérotée : pastille orange avec le numéro + titre Archivo
\newcounter{popsec}
\newcommand{\coursec}[1]{%
  \stepcounter{popsec}\setcounter{popsub}{0}%
  \par\vspace{16pt}\noindent
  \begin{minipage}{\linewidth}
  \tikz[baseline=-0.7ex]\node[draw=popink,line width=1.6pt,fill=poporange,rounded corners=4pt,minimum size=22pt,inner sep=0]{\popdisplay\fontsize{12}{12}\selectfont\color{popcream}\thepopsec};%
  \hspace{8pt}{\popdisplay\fontsize{15}{17}\selectfont\color{popink}\MakeUppercase{#1}}\par
  \vspace{4pt}\noindent{\color{poporange}\rule{36pt}{3.6pt}}
  \end{minipage}\par\nopagebreak\vspace{8pt}
}
\newcounter{popsub}
\newcommand{\courssub}[1]{%
  \stepcounter{popsub}%
  \par\vspace{9pt}\noindent{\popxbold\fontsize{11.5}{13}\selectfont\color{popink}{\color{poporange}\thepopsec.\thepopsub}\hspace{6pt}#1}\par\nopagebreak\vspace{4pt}
}

% ============================================================
% Boîtes pédagogiques — même recette : fond blanc, bordure épaisse colorée,
% ombre dure encre, badge plein. Titre optionnel [#1].
% ============================================================
\tcbset{popbase/.style={breakable,enhanced,colback=white,boxrule=2.3pt,arc=9pt,
  shadow={3pt}{-3pt}{0pt}{popink},left=14pt,right=14pt,top=11pt,bottom=11pt,before skip=11pt,after skip=15pt,
  fontupper=\popsemi\fontsize{10.6}{14.5}\selectfont\color{popink},
  fontlower=\popsemi\fontsize{10.6}{14.5}\selectfont\color{popink}}}
% Coche dessinée (le glyphe ✓ n'existe pas dans les polices du document)
\newcommand{\popcheck}[1]{\tikz[baseline=-0.5ex]\draw[#1,line width=1.6pt,line cap=round,line join=round](-0.13,0.0)--(-0.04,-0.09)--(0.13,0.1);}
\providecommand{\coche}{\popcheck{popgreen}}
\providecommand{\resultat}[1]{\par\smallskip\noindent\fcolorbox{popgreen}{popgreenL}{\parbox{\dimexpr\linewidth-2\fboxsep-2\fboxrule\relax}{\textbf{Résultat.} #1}}\par\smallskip}
\newcommand{\popboxhead}[3]{\poppill{#1}{#2}{white}\ifx\relax#3\relax\else\hspace{6pt}{\popxbold\fontsize{10.6}{12}\selectfont\color{popink}#3}\fi\par\vspace{7pt}}

\newtcolorbox{aretenir}[1][]{popbase,colframe=popgreen,before upper={\popboxhead{À retenir}{popgreen}{#1}}}
\newtcolorbox{exemplebox}[1][]{popbase,colframe=poporange,before upper={\popboxhead{Exemple}{poporange}{#1}}}
\newtcolorbox{definition}[1][]{popbase,colframe=popblue,before upper={\popboxhead{Définition}{popblue}{#1}}}
\newtcolorbox{propriete}[1][]{popbase,colframe=poppurple,before upper={\popboxhead{Propriété}{poppurple}{#1}}}
\newtcolorbox{methode}[1][]{popbase,colframe=popgold,before upper={\popboxhead{Méthode}{popgold}{#1}}}
\newtcolorbox{attention}[1][]{popbase,colframe=poppink,before upper={\popboxhead{Attention}{poppink}{#1}}}
\newtcolorbox{experience}[1][]{popbase,colframe=popblue,colback=popblueL,before upper={\popboxhead{Expérience}{popblue}{#1}}}
% « Le savais-tu ? » : carte violette pleine (contexte, culture, Cameroun)
\newtcolorbox{savaistu}[1][]{popbase,colback=poppurple,colframe=popink,
  fontupper=\popsemi\fontsize{10.4}{14.2}\selectfont\color{white},
  before upper={\poppill{Le savais-tu\,?}{white}{poppurple}\ifx\relax#1\relax\else\hspace{6pt}{\popxbold\color{white}#1}\fi\par\vspace{7pt}}}
% Exercice résolu : énoncé en haut, \tcblower puis la solution sur fond crème
\newcounter{popexo}\newcounter{popexr}
\newtcolorbox{exoresolu}[1][]{popbase,colframe=poporange,colbacklower=popcream,
  segmentation style={popink,line width=1.2pt,dashed},
  before upper={\stepcounter{popexr}\popboxhead{Exercice résolu \thepopexr}{poporange}{#1}},
  before lower={\poppill{Solution}{popink}{popcream}\par\vspace{6pt}}}
% Exercice (non résolu en place) — la correction est dans la partie Corrigés
\newtcolorbox{exercice}[1][]{popbase,colframe=popink,
  before upper={\stepcounter{popexo}\popboxhead{Exercice \thepopexo}{popink}{#1}}}
% Corrigé
\newtcolorbox{corrige}[1][]{popbase,colframe=popgreen,colback=popgreenL,
  before upper={\popboxhead{Corrigé}{popgreen}{#1}}}
% Objectifs / prérequis / bilan
\newtcolorbox{objectifs}{popbase,colframe=popink,colback=poporangeL,
  before upper={\popboxhead{Objectifs de la leçon}{popink}{}}}
\newtcolorbox{prerequis}{popbase,colframe=popink,colback=popblueL,
  before upper={\popboxhead{Prérequis}{popblue}{}}}
\newtcolorbox{bilan}{popbase,colframe=popink,colback=white,boxrule=2.8pt,
  before upper={\popboxhead{Fiche bilan}{poporange}{L'essentiel à connaître pour l'examen}}}
% Figure encadrée avec légende
\newtcolorbox{popfigure}[1][]{enhanced,breakable=false,colback=white,colframe=popline,boxrule=1.4pt,arc=8pt,
  left=10pt,right=10pt,top=10pt,bottom=8pt,before skip=10pt,after skip=12pt,halign=center,
  lower separated=false,halign lower=center,
  fontlower=\popsemi\fontsize{9}{11.5}\selectfont\color{popmuted},
  #1}
\newcommand{\legende}[1]{\par\vspace{6pt}{\popsemi\fontsize{9}{11.5}\selectfont\color{popmuted}#1}}

% ============================================================
% Signature OnBuch+
% ============================================================
\newcommand{\poplogoinline}[1]{{\poplogo\fontsize{#1}{#1}\selectfont\color{popink}OnBuch\color{poporange}+}}
\newcommand{\signatureonbuch}{%
  \begin{tcolorbox}[enhanced,colback=popink,colframe=popink,boxrule=2.2pt,arc=10pt,
      drop shadow=poporange,left=14pt,right=14pt,top=10pt,bottom=10pt,before skip=0pt,after skip=0pt]
    \tikz[baseline=-0.7ex]\node[circle,fill=popgreen,draw=popcream,line width=1.3pt,minimum size=22pt,inner sep=0]{\popcheck{white}};%
    \hspace{9pt}\begin{minipage}[c]{0.62\linewidth}
      {\popxbold\fontsize{12}{13}\selectfont\color{popcream}Signé\ \textbullet\ {\poplogo OnBuch\color{poporange}+}}\par\vspace{2pt}
      {\raggedright\popsemi\fontsize{8}{10.5}\selectfont\color{popline}\MakeUppercase{Équipe pédagogique OnBuch+}\\\MakeUppercase{Conforme au programme officiel MINESEC}\par}
    \end{minipage}\hfill
    {\poplogo\fontsize{22}{22}\selectfont\color{popcream}OnBuch\color{poporange}+}
  \end{tcolorbox}%
}

% ============================================================
% Page de garde
% #1 titre · #2 matière · #3 niveau · #4 module · #5 n° leçon · #6 durée ·
% #7 description / objectifs (texte ou itemize)
% ============================================================
\newcommand{\pagegarde}[7]{%
  \thispagestyle{empty}
  \newgeometry{a4paper,margin=1.7cm,top=1.6cm,bottom=1.4cm}%
  \begin{tikzpicture}[remember picture,overlay]
    \fill[poporange!18] ([xshift=-3.2cm,yshift=-3.6cm]current page.north east) circle (4.6cm);
    \fill[poppurple!14] ([xshift=2.4cm,yshift=7.6cm]current page.south west) circle (3.8cm);
  \end{tikzpicture}%
  {\poplogo\fontsize{30}{30}\selectfont\color{popink}OnBuch\color{poporange}+}\hfill
  \raisebox{0.6ex}{\poppill{Cours signé \textbullet\ Premium}{popink}{popcream}}\par
  \vspace{1pt}\popsquiggle{poppurple}\par
  \vspace{8pt}
  \begin{tcolorbox}[enhanced,colback=poporange,colframe=popink,boxrule=2.8pt,arc=13pt,
      drop shadow=popink,left=18pt,right=18pt,top=12pt,bottom=13pt,after skip=10pt]
    \poppill{#2 \textbullet\ #3}{popink}{popcream}\hspace{5pt}\poppill{#4}{white}{popink}\par\vspace{8pt}
    {\popdisplay\fontsize{14}{14}\selectfont\color{popink}LEÇON #5}\par\vspace{4pt}
    {\raggedright\hyphenpenalty=10000\exhyphenpenalty=10000\popdisplay\fontsize{25}{27}\selectfont\color{popcream}\MakeUppercase{#1}\par}
    \vspace{8pt}
    {\popxbold\fontsize{9.5}{9.5}\selectfont\color{popink}\MakeUppercase{Durée conseillée : #6}}
  \end{tcolorbox}
  \begin{tcolorbox}[enhanced,colback=white,colframe=popink,boxrule=2.2pt,arc=10pt,
      drop shadow=popink,left=16pt,right=16pt,top=10pt,bottom=10pt,after skip=8pt]
    \poppill{Dans cette leçon}{poppurple}{white}\par\vspace{5pt}
    \popsemi\fontsize{9.3}{12.6}\selectfont\color{popink}#7
  \end{tcolorbox}
  \noindent\begin{minipage}[t]{0.487\linewidth}
    \begin{tcolorbox}[enhanced,colback=popgreen,colframe=popink,boxrule=2.2pt,arc=10pt,
        drop shadow=popink,left=11pt,right=11pt,top=10pt,bottom=10pt,height=3.1cm,valign=center]
      \tcbox[colback=white,colframe=popink,boxrule=1.3pt,arc=5pt,boxsep=0pt,nobeforeafter,
        left=3pt,right=3pt,top=3pt,bottom=3pt,valign=center]{\qrcode[height=1.9cm]{https://play.google.com/store/apps/details?id=com.luvvix.onbuch&hl=fr}}%
      \hspace{8pt}\begin{minipage}[c]{0.5\linewidth}
        {\popdisplay\fontsize{11}{12.5}\selectfont\color{white}\MakeUppercase{Télécharge OnBuch}}\par\vspace{4pt}
        {\raggedright\popsemi\fontsize{8.4}{11}\selectfont\color{white}Gratuit \textbullet\ Google Play\\Tuteur IA, annales, concours, résultats\par}
      \end{minipage}
    \end{tcolorbox}
  \end{minipage}\hfill
  \begin{minipage}[t]{0.487\linewidth}
    \begin{tcolorbox}[enhanced,colback=poppurple,colframe=popink,boxrule=2.2pt,arc=10pt,
        drop shadow=popink,left=11pt,right=11pt,top=10pt,bottom=10pt,height=3.1cm,valign=center]
      \tikz[baseline=-0.7ex]\node[circle,fill=white,draw=popink,line width=1.3pt,minimum size=1.6cm,inner sep=0]{\popdisplay\fontsize{24}{24}\selectfont\color{poppurple}+};%
      \hspace{8pt}\begin{minipage}[c]{0.6\linewidth}
        {\popdisplay\fontsize{11}{12.5}\selectfont\color{white}\MakeUppercase{Passe à OnBuch+}}\par\vspace{4pt}
        {\raggedright\popsemi\fontsize{8.4}{11}\selectfont\color{white}Tous les cours signés, Tuteur IA illimité, fiches PDF — onglet OnBuch+ de l'app\par}
      \end{minipage}
    \end{tcolorbox}
  \end{minipage}
  \vspace{8pt}
  \popcontenu
  \vfill
  \signatureonbuch
  \restoregeometry
  \newpage
}

% Rangée « Ce cours contient » (page de garde)
\newcommand{\poptile}[3]{% #1 couleur #2 gros texte #3 légende
  \begin{tcolorbox}[enhanced,colback=#1,colframe=popink,boxrule=2pt,arc=9pt,shadow={2.5pt}{-2.5pt}{0pt}{popink},
      left=6pt,right=6pt,top=9pt,bottom=9pt,halign=center,valign=center,height=1.75cm,width=\linewidth,nobeforeafter]
    {\popdisplay\fontsize{12.5}{13}\selectfont\color{white}\MakeUppercase{#2}}\par\vspace{3pt}
    {\centering\popsemi\fontsize{7.8}{9.5}\selectfont\color{white}#3\par}
  \end{tcolorbox}}
\newcommand{\popcontenu}{%
  \par\noindent{\popxboldls\fontsize{7.5}{7.5}\selectfont\color{popmuted}CE COURS SIGNÉ CONTIENT}\par\vspace{7pt}
  \noindent\begin{minipage}[t]{0.235\linewidth}\poptile{popblue}{Cours}{complet, illustré, pas à pas}\end{minipage}\hfill
  \begin{minipage}[t]{0.235\linewidth}\poptile{poporange}{Méthodes}{et exercices résolus}\end{minipage}\hfill
  \begin{minipage}[t]{0.235\linewidth}\poptile{poppink}{Exercices}{type examen + corrigés}\end{minipage}\hfill
  \begin{minipage}[t]{0.235\linewidth}\poptile{popgreen}{Fiche bilan}{l'essentiel à retenir}\end{minipage}\par}

% Sommaire
\newtcolorbox{sommaire}{enhanced,colback=white,colframe=popink,boxrule=2.2pt,arc=10pt,
  drop shadow=popink,left=18pt,right=18pt,top=13pt,bottom=13pt,before skip=6pt,after skip=20pt,
  before upper={\poppill{Sommaire}{poppurple}{white}\par\vspace{9pt}\popsemi\fontsize{10.6}{14.5}\selectfont\color{popink}}}

% Signature de fin de cours — poussée tout en bas de la dernière page
\newcommand{\findecours}{%
  \par\vfill\nopagebreak
  \begin{center}\popsquiggle{poporange}\end{center}\vspace{4pt}
  \signatureonbuch
}

%% FIGFIX-BEGIN v3 — correctifs globaux des figures (docs/onbuch-generation/tools/figfix)
\usepackage{adjustbox}\usepackage{etoolbox}
% toute figure TikZ/pgfplots plus large que la ligne est réduite à la largeur disponible
\BeforeBeginEnvironment{tikzpicture}{\begin{adjustbox}{max width=\linewidth}}
\AfterEndEnvironment{tikzpicture}{\end{adjustbox}}
% légendes pgfplots lisibles par-dessus les courbes
\pgfplotsset{every axis/.append style={legend style={fill=white,fill opacity=0.92,text opacity=1,draw=black!15,inner sep=3pt}}}
% étiquettes d'arêtes (syntaxe edge["..."]) avec fond blanc
\tikzset{every edge quotes/.append style={fill=white,inner sep=1.2pt}}
% bulles de cartes mentales : texte à la ligne dans la bulle, bulle un peu plus grande
\tikzset{every concept/.append style={text width=2.0cm,align=center,minimum size=2.3cm,inner sep=2pt}}
%% FIGFIX-END
\begin{document}

\pagegarde{Dipôles commandés, capteurs et chaînes électroniques}{Physique}{Tle D}{Module 3 : Électricité — évolutions temporelles des circuits électriques et électroniques}{P13}{3 h}{Cette leçon introduit les dipôles commandés (manuellement ou électriquement) et les capteurs qui transforment une grandeur physique (température, lumière, pression, ondes) en signal électrique. Elle aboutit à la notion de chaîne électronique — capteur, traitement, actionneur — illustrée par des applications concrètes : sirène commandée par la lumière, indicateur de niveau d'eau, poste radio.
\vspace{4pt}
\begin{multicols}{2}\small
\begin{itemize}
\item À la fin de cette leçon, je sais définir un dipôle commandé et un capteur, et expliquer le principe du captage d'une grandeur physique.
\item Je sais citer et décrire des dipôles commandés manuellement (rhéostat) et électriquement (relais, VDR, diode, transistor, électrodes).
\item Je sais interpréter les caractéristiques U = f(I, R) d'un rhéostat, R = f(θ) d'une thermistance et U = f(I, Φ) d'une photorésistance.
\item Je sais distinguer un capteur linéaire d'un capteur non linéaire à partir de sa courbe caractéristique.
\item Je sais identifier les trois éléments d'une chaîne électronique (capteur, traitement, actionneur) et la schématiser.
\item Je sais expliquer le fonctionnement d'une chaîne commandée par la température, la lumière ou un relais.
\end{itemize}\end{multicols}}

\begin{sommaire}
\begin{multicols}{2}
\begin{enumerate}
\item Dipôles commandés et principe du captage
\item Étude du rhéostat : caractéristique U = f(I, R)
\item Capteurs de température, de lumière, de pression et d'ondes
\item Le transistor en commutation et les électrodes
\item Chaînes électroniques : structure et exemples
\item Activité d'intégration
\item Méthodes et exercices résolus
\item Exercices
\item Corrigés des exercices
\item Fiche bilan
\end{enumerate}
\end{multicols}
\end{sommaire}

\begin{objectifs}
\begin{itemize}
\item À la fin de cette leçon, je sais définir un dipôle commandé et un capteur, et expliquer le principe du captage d'une grandeur physique.
\item Je sais citer et décrire des dipôles commandés manuellement (rhéostat) et électriquement (relais, VDR, diode, transistor, électrodes).
\item Je sais interpréter les caractéristiques U = f(I, R) d'un rhéostat, R = f(θ) d'une thermistance et U = f(I, Φ) d'une photorésistance.
\item Je sais distinguer un capteur linéaire d'un capteur non linéaire à partir de sa courbe caractéristique.
\item Je sais identifier les trois éléments d'une chaîne électronique (capteur, traitement, actionneur) et la schématiser.
\item Je sais expliquer le fonctionnement d'une chaîne commandée par la température, la lumière ou un relais.
\item Je sais expliquer le fonctionnement d'une antenne, d'un microphone et d'écouteurs téléphoniques.
\item Je sais analyser un montage à transistor utilisé en commutation (bloqué/saturé) dans une chaîne électronique.
\end{itemize}
\end{objectifs}

\begin{prerequis}
\begin{itemize}
\item Loi d'Ohm et caractéristique U = f(I) d'un dipôle (Première)
\item Puissance et énergie électriques, montages série et dérivation
\item Notion de résistance électrique et mesures au multimètre
\item Champ magnétique créé par un courant et force électromagnétique (notions de Première)
\item Notion d'onde électromagnétique et de son (Première)
\end{itemize}
\end{prerequis}

\newpage

\coursec{Dipôles commandés et principe du captage}

À Yaoundé, la sirène d'une école se déclenche automatiquement à la tombée de la nuit, sans que personne n'appuie sur un interrupteur. Dans un quartier de Douala, un réservoir d'eau déclenche une pompe dès que le niveau baisse. À la radio locale, la voix de l'animateur captée par un microphone ressort des écouteurs d'un auditeur à des centaines de kilomètres. Comment un simple composant électronique peut-il \emph{sentir} la lumière, la température ou le niveau d'eau et commander un appareil ? C'est toute la logique des \cle{dipôles commandés}, des \cle{capteurs} et des \cle{chaînes électroniques} que cette leçon va dévoiler.

\textbf{Problématique :} qu'est-ce qui distingue un dipôle ordinaire d'un dipôle commandé ? Comment une grandeur physique non électrique (lumière, température, pression) peut-elle être transformée en un signal électrique exploitable ?

\courssub{Du dipôle passif au dipôle commandé}

\textbf{Observation.} Prends un résistor ordinaire de résistance $R = \SI{100}{\ohm}$ : qu'il fasse jour ou nuit, chaud ou froid, sa résistance reste pratiquement constante. Sa caractéristique $U = RI$ est une droite passant par l'origine : c'est un \cle{dipôle passif} dont le paramètre $R$ ne dépend d'aucune grandeur extérieure. Prends maintenant une photorésistance : exposée au soleil, sa résistance vaut quelques centaines d'ohms ; placée dans l'obscurité, elle grimpe à plusieurs centaines de kilo-ohms. Son paramètre électrique $R$ est \emph{piloté} par une grandeur extérieure, l'éclairement. C'est un dipôle commandé.

\begin{definition}[Dipôle commandé]
Un \cle{dipôle commandé} est un dipôle dont une grandeur électrique caractéristique (résistance $R$, intensité $I$ ou tension $U$) est \textbf{contrôlée par une autre grandeur}, appelée \emph{grandeur de commande}.
\begin{itemize}
\item Si la grandeur de commande est une action de l'opérateur (rotation d'un bouton, translation d'un curseur), la commande est \cle{manuelle}.
\item Si la grandeur de commande est une grandeur électrique (courant ou tension), la commande est \cle{électrique}.
\item Si la grandeur de commande est une grandeur physique non électrique (température, lumière, pression, niveau, ondes), le dipôle commandé joue le rôle de \cle{capteur}.
\end{itemize}
\end{definition}

\courssub{Commande manuelle : le rhéostat}

\textbf{Description.} Un \cle{rhéostat} est une résistance variable constituée d'un fil résistif (ou d'une piste de carbone) enroulé sur un support isolant, sur lequel glisse un \cle{curseur} $C$. Le dipôle possède trois bornes : les deux extrémités $A$ et $B$ de la piste, et le curseur $C$. En déplaçant le curseur, on fait varier la longueur de fil traversée par le courant, donc la résistance, conformément à la loi de Pouillet :
\[ R = \rho \, \dfrac{\ell}{S} \]
où $\rho$ est la résistivité du matériau ($\unit{\ohm \cdot m}$), $\ell$ la longueur de fil utilisée (m) et $S$ sa section ($\unit{m^2}$). L'analyse dimensionnelle est immédiate : $[\rho]\times\dfrac{[\ell]}{[S]} = \dfrac{\Omega\cdot\mathrm{m}\cdot\mathrm{m}}{\mathrm{m}^2} = \Omega$. La résistance $R_{AC}$ entre l'extrémité $A$ et le curseur $C$ est donc proportionnelle à la position du curseur.

\textbf{Deux branchements possibles.}
\begin{itemize}
\item \textbf{Montage série (rhéostat)} : on utilise seulement \emph{deux bornes} ($A$ et $C$). Placé en série avec un récepteur, le rhéostat fait varier l'intensité du courant : $I = \dfrac{E}{R_{AC} + R_{\text{charge}}}$. Quand $R_{AC}$ augmente, $I$ diminue.
\item \textbf{Montage potentiométrique} : on utilise les \emph{trois bornes}. La piste $AB$ est branchée aux bornes du générateur ; on prélève entre $C$ et $B$ une fraction de la tension : $U_{CB} = \dfrac{R_{CB}}{R_{AB}}\,E$. C'est un \emph{diviseur de tension réglable} : la tension de sortie varie de $0$ à $E$.
\end{itemize}

\begin{popfigure}
\begin{center}
\begin{tikzpicture}[scale=1]
% Piste résistive
\draw[very thick,popblue] (0,0) -- (5,0);
\foreach \x in {0.25,0.75,...,4.75} \draw[popblue] (\x,0) -- (\x+0.25,0.35);
\draw[very thick,popblue] (0,0) -- (5,0);
% Bornes A et B
\draw[fill] (0,0) circle (0.06) node[below left=2pt] {$A$};
\draw[fill] (5,0) circle (0.06) node[below right=2pt] {$B$};
\draw (0,0) -- (0,-0.8); \draw (5,0) -- (5,-0.8);
% Curseur
\draw[fill,poporange] (2.6,0) circle (0.05);
\draw[very thick,poporange,-{Stealth}] (2.6,1.1) -- (2.6,0.05);
\draw[fill] (2.6,1.1) circle (0.06) node[above] {$C$ (curseur)};
\draw (2.6,1.1) -- (2.6,1.6);
\draw[{Stealth}-{Stealth},popdark] (1.9,1.35) -- (3.3,1.35) node[midway,above,font=\small] {déplacement};
% Symbole normalisé
\begin{scope}[xshift=7.2cm,yshift=0.2cm]
\draw[thick,popdark] (0,0) rectangle (2.2,0.7);
\draw (-0.8,0.35) -- (0,0.35); \draw (2.2,0.35) -- (3.0,0.35);
\draw[thick,poporange,-{Stealth}] (0.4,1.3) -- (1.5,0.72);
\node[font=\small] at (1.1,-0.4) {symbole du rhéostat};
\end{scope}
\end{tikzpicture}
\end{center}
\legende{Rhéostat à curseur : la résistance utilisée $R_{AC}$ dépend de la position du curseur $C$ sur la piste résistive. À droite, symbole normalisé.}
\end{popfigure}

\begin{attention}[Rhéostat ou potentiomètre ?]
Ne confonds pas les deux montages ! Le \textbf{rhéostat} n'utilise que \emph{deux bornes} et se branche \emph{en série} : il règle le \emph{courant}. Le \textbf{potentiomètre} utilise les \emph{trois bornes} en diviseur de tension : il règle la \emph{tension}. Au Bac, compter les bornes réellement connectées sur le schéma est le réflexe qui sauve des points.
\end{attention}

\begin{exoresolu}[Caractéristique d'un rhéostat]
Un rhéostat de résistance totale $R_{AB} = \SI{1,0}{k\ohm}$ est alimenté par un générateur de tension $E = \SI{12}{V}$ (montage potentiométrique). Le curseur est placé de sorte que $R_{CB} = \SI{400}{\ohm}$.
\begin{enumerate}
\item Calculer la tension $U_{CB}$ aux bornes de la partie $CB$.
\item Quelle puissance est dissipée dans la piste $AB$ ?
\end{enumerate}
\tcblower
\textbf{1.} Diviseur de tension : $U_{CB} = \dfrac{R_{CB}}{R_{AB}} E = \dfrac{400}{1000} \times 12 = \SI{4,8}{V}$.
\textbf{2.} La piste $AB$ est branchée directement sur le générateur : $P = \dfrac{E^2}{R_{AB}} = \dfrac{12^2}{1000} = \SI{0,144}{W} = \SI{144}{mW}$.
\end{exoresolu}

\courssub{Commande électrique : relais, VDR, diode, transistor, électrodes}

\textbf{Le relais électromagnétique.} Un relais comporte deux circuits \emph{électriquement isolés} l'un de l'autre (séparation galvanique) :
\begin{itemize}
\item un \emph{circuit de commande} : une bobine enroulée autour d'un noyau de fer doux (électroaimant), parcourue par un faible courant $i$ ;
\item un \emph{circuit de puissance} : un contact mobile (palette) attiré par le noyau lorsque la bobine est alimentée, qui ferme (contact \emph{repos} / \emph{travail}) ou ouvre le circuit d'un appareil puissant (moteur, sirène, lampe).
\end{itemize}

\begin{popfigure}
\begin{center}
\begin{tikzpicture}[scale=0.95]
% Bobine (électroaimant)
\draw[thick,popblue] (0,0) rectangle (0.6,2.2);
\node[popblue,font=\small,rotate=90] at (0.3,1.1) {bobine};
\foreach \y in {0.2,0.6,1.0,1.4,1.8} \draw[popblue] (0.6,\y) arc (-90:90:0.12);
% Circuit de commande
\draw (0,0) -- (-1.6,0) to[battery1,l={$E_1$}] (-1.6,2.2) -- (0,2.2);
\node[font=\small,popblue] at (-2.5,2.2) {circuit de commande};
% Noyau + palette
\draw[thick,popdark] (1.2,2.0) -- (2.6,2.0);
\draw[very thick,poporange] (1.3,1.75) -- (3.0,1.45) node[pos=0.9,below,font=\small] {palette};
\draw[fill] (1.3,1.75) circle (0.05);
% Contact fixe (travail)
\draw[fill,popdark] (3.0,1.15) circle (0.06);
% Circuit de puissance
\draw (3.0,1.15) -- (4.6,1.15) to[lamp,l_={$M$}] (4.6,-0.6) -- (1.3,-0.6) -- (1.3,1.7);
\node[font=\small,poppurple,align=center] at (3.2,-1.0) {circuit de puissance};
\draw[dashed,popmuted] (0.9,2.5) -- (0.9,-0.9);
\node[font=\small,popmuted] at (0.9,2.8) {isolation galvanique};
\end{tikzpicture}
\end{center}
\legende{Relais électromagnétique : le faible courant de la bobine (à gauche) attire la palette et ferme le circuit de puissance (à droite). Les deux circuits sont galvaniquement isolés.}
\end{popfigure}

\begin{propriete}[Fonctionnement du relais (admis)]
Lorsqu'un courant, même faible, parcourt la bobine du relais, l'électroaimant attire la palette et \textbf{commande la fermeture (ou l'ouverture) d'un circuit de puissance}. Les circuits de commande et de puissance sont \textbf{galvaniquement séparés} : aucun courant ne passe de l'un à l'autre. Le relais permet donc de commander un appareil sous \SI{220}{\volt} (réseau ENEO) à partir d'un circuit de commande sous quelques volts seulement, en toute sécurité.
\end{propriete}

\textbf{La VDR (varistance).} Une \cle{VDR} (\emph{Voltage Dependent Resistor}) est un dipôle dont la résistance dépend fortement de la tension à ses bornes : très grande (quelques $\si{M\ohm}$) tant que $U$ reste sous une tension seuil $U_n$, elle s'effondre brutalement (quelques $\si{\ohm}$) au-delà. La VDR est donc \emph{commandée par la tension}. Elle protège les circuits contre les surtensions transitoires (foudre, commutation).

\textbf{La diode.} Une \cle{diode} ne conduit le courant que si la tension à ses bornes dans le sens direct dépasse la \emph{tension de seuil} $U_S \approx \SI{0,6}{\volt}$ (silicium). Elle est commandée par la tension appliquée : bloquée en dessous du seuil ou en sens inverse, passante au-dessus. En commutation, elle agit comme un interrupteur commandé par la polarité de la tension.

\textbf{Le transistor bipolaire (NPN) en commutation.} Un \cle{transistor} bipolaire NPN possède trois électrodes : l'\emph{émetteur} $E$, le \emph{collecteur} $C$ et la \emph{base} $B$. En régime de \textbf{commutation}, il fonctionne en \emph{tout-ou-rien} entre deux états stables :
\begin{itemize}
\item \textbf{État bloqué (ouvert)} : $I_B = 0 \Rightarrow I_C = 0$. La tension collecteur-émetteur $U_{CE} = V_{CC}$ (tension d'alimentation). Le transistor équivaut à un interrupteur \emph{ouvert}.
\item \textbf{État saturé (fermé)} : $I_B \ge I_{B,\text{sat}}$. Le transistor conduit le courant maximal imposé par le circuit extérieur : $I_C = \dfrac{V_{CC} - U_{CE,\text{sat}}}{R_C}$ avec $U_{CE,\text{sat}} \approx \SI{0,1}{\volt}$. Le transistor équivaut à un interrupteur \emph{fermé} (résistance résiduelle très faible).
\end{itemize}
Le courant de base nécessaire pour saturer est $I_{B,\text{sat}} = \dfrac{I_{C,\text{sat}}}{\beta_{\min}}$ (on prend souvent un facteur de sursaturation $k \approx 3$ à $10$, soit $I_B = k \cdot I_{B,\text{sat}}$). Un courant de base de quelques microampères suffit à commander un courant de collecteur de plusieurs milliampères (voire ampères avec un transistor de puissance) : le transistor est un \emph{interrupteur commandé en courant}.

\begin{popfigure}
\begin{center}
\begin{tikzpicture}[scale=1]
% Schéma de commande transistor NPN (charge résistive + relais)
\draw (0,0) to[battery1,l_={$V_{CC}$}] (0,3);
\draw (0,3) to[R=$R_C$,-*] (3,3);
\draw (3,3) -- (4,3);
\draw (4,3) to[npn,anchor=B] (4,0) node[npn] (Q) {};
\draw (Q.E) -- (4,0) -- (0,0);
% Commande base
\draw (Q.B) -- ++(-1.5,0) to[R=$R_B$] ++(-2,0) to[battery1,l_={$V_{cmd}$}] ++(0,-1.5) -| (Q.E);
\node[align=center,font=\small] at (-2,-0.75) {Commande\\logique};
\end{tikzpicture}
\end{center}
\legende{Montage commutateur transistor NPN : commande en base, charge en collecteur (résistance, relais, LED...).}
\end{popfigure}

\begin{methode}[Dimensionner la résistance de base $R_B$ pour saturer un transistor NPN]
\begin{enumerate}
\item Déterminer le courant de collecteur souhaité en saturation : $I_{C,\text{sat}} = \dfrac{V_{CC} - U_{CE,\text{sat}}}{R_C}$ (avec $U_{CE,\text{sat}} \approx \SI{0,1}{V}$).
\item Choisir un gain $\beta_{\min}$ (datasheet, typiquement $50$--$100$ en puissance). Calculer $I_{B,\text{min}} = \dfrac{I_{C,\text{sat}}}{\beta_{\min}}$.
\item Appliquer un facteur de sursaturation $k$ (souvent $k=5$ à $10$) : $I_B = k \cdot I_{B,\text{min}}$.
\item Appliquer la loi des mailles à la boucle de base : $V_{cmd} = R_B I_B + U_{BE}$ (avec $U_{BE} \approx \SI{0,7}{V}$ en conduction). En déduire $R_B = \dfrac{V_{cmd} - 0,7}{I_B}$.
\item Choisir la valeur normalisée la plus proche (inférieure pour garantir la saturation).
\end{enumerate}
\end{methode}

\textbf{Les électrodes.} Deux électrodes métalliques plongées dans un liquide conducteur (eau contenant des sels dissous) laissent passer un courant dès que le liquide établit le contact entre elles. La conduction est donc \emph{commandée par la présence du liquide} : c'est le principe de l'indicateur de niveau d'eau du réservoir de Douala. Le liquide joue le rôle de résistance variable (quasi nulle quand les électrodes sont immergées, infinie sinon).

\begin{exoresolu}[Commande d'un relais par transistor]
On veut commander un relais de bobine $R_C = \SI{220}{\ohm}$ sous $V_{CC} = \SI{12}{V}$ par un transistor NPN ($\beta_{\min}=50$, $U_{CE,\text{sat}}=\SI{0,1}{V}$, $U_{BE}=\SI{0,7}{V}$). La commande logique fournit $V_{cmd} = \SI{5}{V}$ (niveau haut). Calculer $R_B$ pour assurer la saturation (facteur $k=5$). Prévoir la diode de roue libre.
\tcblower
\textbf{1. Courant de collecteur en saturation :}
$I_{C,\text{sat}} = \dfrac{12 - 0,1}{220} \approx \SI{54,1}{mA}$.
\textbf{2. Courant de base minimum :}
$I_{B,\text{min}} = \dfrac{54,1}{50} \approx \SI{1,08}{mA}$.
\textbf{3. Courant de base imposé (sursaturation $k=5$) :}
$I_B = 5 \times 1,08 = \SI{5,4}{mA}$.
\textbf{4. Résistance de base :}
$R_B = \dfrac{5 - 0,7}{5,4 \times 10^{-3}} \approx \SI{796}{\ohm}$. On choisit $R_B = \SI{680}{\ohm}$ (valeur normalisée E12 inférieure).
\textbf{5. Diode de roue libre :} Une diode (ex. 1N4007) en parallèle sur la bobine (cathode au $+12V$) protège le transistor de la surtension d'auto-induction ($u_L = -L\,di/dt$) à l'ouverture.
\end{exoresolu}

\begin{savaistu}[La VDR, gardienne de tes appareils]
Au Cameroun, les orages et les fluctuations du réseau ENEO provoquent des \emph{surtensions} brutales qui détruisent téléviseurs et téléphones en charge. La \cle{VDR} (varistance), placée en parallèle à l'entrée des multiprises parafoudre, présente une résistance énorme en fonctionnement normal ; mais dès que la tension dépasse son seuil, sa résistance s'effondre : elle devient brusquement conductrice et dévie le courant de la surtension, protégeant ainsi les appareils branchés en aval. C'est un dipôle commandé par la tension, héros discret de nos installations électriques !
\end{savaistu}

\courssub{Capteurs de température, de lumière, de pression et d'ondes}

\begin{definition}[Capteur]
Un \cle{capteur} est un dispositif qui transforme une \textbf{grandeur physique} (température, lumière, pression, ondes mécaniques ou électromagnétiques) en une \textbf{grandeur électrique mesurable} (résistance, tension ou courant), appelée \emph{signal électrique}.
\end{definition}

Le principe général du \cle{captage} se décompose en trois étapes : la grandeur physique modifie une propriété électrique d'un composant ; cette modification se traduit en variation de résistance, de tension ou de courant ; le signal obtenu est ensuite exploité (mesuré, amplifié, traité).

\begin{popfigure}
\begin{center}
\begin{tikzpicture}[node distance=0.5cm and 1.1cm]
\node[draw=popblue,fill=popblueL,rounded corners,thick,align=center,minimum height=1cm] (gp) {Grandeur\\physique\\(lumière, $\theta$, pression)};
\node[draw=poporange,fill=poporangeL,rounded corners,thick,align=center,minimum height=1cm,right=of gp] (cap) {\textbf{Capteur}\\(transducteur)};
\node[draw=popgreen,fill=popgreenL,rounded corners,thick,align=center,minimum height=1cm,right=of cap] (sig) {Signal\\électrique\\($R$, $U$ ou $I$)};
\draw[-{Stealth},very thick,popdark] (gp) -- (cap);
\draw[-{Stealth},very thick,popdark] (cap) -- (sig);
\end{tikzpicture}
\end{center}
\legende{Schéma fonctionnel du captage : la grandeur physique est convertie en signal électrique exploitable.}
\end{popfigure}

On distingue deux grandes familles de capteurs :
\begin{itemize}
\item Les \cle{capteurs passifs} : leur \emph{résistance} varie avec la grandeur physique ; ils nécessitent une alimentation extérieure (montage en diviseur de tension ou pont de Wheatstone) pour fournir un signal de tension ou courant. Exemples : la thermistance CTN (température), la photorésistance (lumière), la jauge de contrainte (pression).
\item Les \cle{capteurs actifs} : ils se comportent en \emph{générateurs} et produisent directement une tension ou un courant à partir de la grandeur physique, sans alimentation auxiliaire. Exemples : le microphone dynamique (le son fait vibrer une bobine dans un champ magnétique $\rightarrow$ induction), l'antenne (l'onde électromagnétique y induit un courant alternatif), la photodiode (effet photovoltaïque), le capteur piézoélectrique.
\end{itemize}

\courssub{Capteurs de température : Thermistance CTN et sonde KTY}

\begin{itemize}
\item \textbf{Thermistance CTN (\emph{Coefficient de Température Négatif}).} Composant en oxydes métalliques (Mn, Ni, Co) dont la résistance \emph{diminue} fortement quand la température $\theta$ augmente. Caractéristique $R = f(\theta)$ : courbe décroissante, \textbf{non linéaire} (exponentielle approchée par loi de Steinhart-Hart : $\frac{1}{T} = A + B\ln R + C(\ln R)^3$). Très sensible, peu coûteuse, utilisée pour mesures ponctuelles (thermostats, sondes moteur, électroménager).
\item \textbf{Sonde KTY (silicium).} Résistance qui \emph{augmente} quasi-linéairement avec la température (CTP). Caractéristique $R = f(\theta)$ : droite quasi-parfaite sur une large plage ($-50^\circ\text{C}$ à $+150^\circ\text{C}$). Précise, stable, utilisée en instrumentation industrielle et automobile.
\end{itemize}

\begin{attention}[Linéarité du capteur]
Un capteur est \textbf{linéaire} si son signal de sortie $S$ est proportionnel à la grandeur mesurée $X$ : $S = k X + b$ (droite). La thermistance CTN est \textbf{non linéaire} (courbe exponentielle) ; la sonde KTY est \textbf{linéaire}. Au Bac, on te demandera de \emph{distinguer} les deux à partir de la courbe $R=f(\theta)$ ou $U=f(\theta)$ : droite $\Rightarrow$ linéaire ; courbe $\Rightarrow$ non linéaire.
\end{attention}

\courssub{Capteur de lumière : Photorésistance (LDR)}

Une \cle{photorésistance} (LDR, \emph{Light Dependent Resistor}) est un composant à base de sulfure de cadmium (CdS) dont la résistance \emph{diminue} quand l'éclairement $\Phi$ (en lux) augmente.
Caractéristique $R = f(\Phi)$ : courbe décroissante, \textbf{non linéaire} (souvent loi de puissance $R \propto \Phi^{-\gamma}$ sur échelle log-log). Résistance noire (obscurité) : $\sim \SI{1}{M\ohm}$ ; pleine lumière : $\sim \SI{100}{\ohm}$--$\SI{1}{k\ohm}$.
Utilisation : détection jour/nuit (réverbères, sirènes d'école), mesure d'éclairement, barrières optiques.

\courssub{Capteur de pression : Jauge de contrainte}

Une \cle{jauge de contrainte} (jauge de déformation) est un fil métallique ou une piste de silicium collée sur un support déformable. Sous l'effet d'une pression (force), le support se déforme, allongeant ou raccourcissant le fil $\Rightarrow$ variation de résistance $\Delta R$ (loi de Pouillet + effet piézorésistif).
Le signal $\Delta R / R$ est très faible ($\sim 10^{-3}$) : il nécessite un \emph{pont de Wheatstone} et un amplificateur d'instrumentation. Application : capteurs de pression pneus, pesons électroniques (marchés, hôpitaux), monitoring barrages (Lom Pangar, Nachtigal).

\courssub{Capteurs d'ondes : Antenne, Microphone, Écouteur}

\begin{itemize}
\item \textbf{Antenne (réception).} Un conducteur parcouru par une onde électromagnétique (champ électrique variable) voit apparaître aux bornes une \textbf{tension alternative} de même fréquence (induction électromagnétique, loi de Faraday). C'est un capteur \emph{actif} (générateur de tension alternative). Exemple : antenne radio FM (88--108 MHz) sur le toit d'un immeuble à Yaoundé.
\item \textbf{Microphone (son $\rightarrow$ électricité).}
    \begin{itemize}
    \item \emph{Dynamique (bobine mobile) :} membrane + bobine dans champ magnétique permanent. Vibration $\rightarrow$ induction $\rightarrow$ tension alternative proportionnelle à la vitesse de la membrane. Capteur actif.
    \item \emph{Électret / Piézoélectrique :} variation de capacité ou effet piézo $\rightarrow$ tension. Nécessite souvent un transistor FET intégré (alimentation fantôme).
    \end{itemize}
\item \textbf{Écouteur / Haut-parleur (électricité $\rightarrow$ son).} Inverse du microphone dynamique : courant alternatif dans bobine $\rightarrow$ force de Laplace $\rightarrow$ vibration membrane $\rightarrow$ onde sonore. C'est un \emph{actionneur} (voir chaîne électronique).
\end{itemize}

\begin{exemplebox}[Caractéristiques $U = f(I, \Phi)$ d'une photorésistance en diviseur]
Une photorésistance $R_P$ forme un diviseur de tension avec une résistance fixe $R_1 = \SI{10}{k\ohm}$ sous $E = \SI{5}{V}$. $U_S$ est la tension aux bornes de $R_1$ (sortie).
\begin{itemize}
\item Jour ($\Phi$ fort) : $R_P = \SI{1}{k\ohm} \Rightarrow U_S = 5 \times \frac{10}{11} \approx \SI{4,55}{V}$.
\item Nuit ($\Phi$ faible) : $R_P = \SI{1}{M\ohm} \Rightarrow U_S = 5 \times \frac{10}{1010} \approx \SI{0,05}{V}$.
\end{itemize}
La tension de sortie $U_S$ \emph{augmente} avec la lumière. On peut comparer $U_S$ à un seuil (comparateur) pour commander un relais.
\end{exemplebox}

\begin{aretenir}[L'essentiel sur les capteurs]
\begin{itemize}
\item \textbf{Capteur passif} : $R = f(\text{grandeur physique})$ (CTN, LDR, jauge). Nécessite polarisation.
\item \textbf{Capteur actif} : Génère $U$ ou $I$ (antenne, microphone dynamique, piézo).
\item \textbf{Linéaire} : courbe droite ($KTY$, jauge dans petit domaine). \textbf{Non linéaire} : courbe courbe ($CTN$, $LDR$).
\item \textbf{Température} : CTN ($R\downarrow$ quand $\theta\uparrow$, non linéaire), KTY ($R\uparrow$ quand $\theta\uparrow$, linéaire).
\item \textbf{Lumière} : LDR ($R\downarrow$ quand $\Phi\uparrow$, non linéaire).
\item \textbf{Ondes} : Antenne (EM $\rightarrow$ $u(t)$), Microphone (Son $\rightarrow$ $u(t)$), Écouteur ($i(t)$ $\rightarrow$ Son).
\end{itemize}
\end{aretenir}

\courssub{Chaînes électroniques : structure et exemples}

\begin{definition}[Chaîne électronique]
Une \cle{chaîne électronique} est un ensemble de dipôles commandés reliés entre eux pour réaliser une fonction automatique : \textbf{Acquisition} (capteur) $\rightarrow$ \textbf{Traitement} (conditionnement, amplification, comparaison, logique) $\rightarrow$ \textbf{Action} (actionneur).
\end{definition}

\begin{popfigure}
\begin{center}
\begin{tikzpicture}[node distance=0.4cm and 2.3cm, every node/.style={font=\small}]
\node[draw=popblue,fill=popblueL,rounded corners,thick,align=center,minimum height=1.2cm] (cap) {CAPTEUR\\Grandeur physique\\$\rightarrow$ Signal él.};
\node[draw=poporange,fill=poporangeL,rounded corners,thick,align=center,minimum height=1.2cm,right=of cap] (tra) {TRAITEMENT\\Amplification,\\Comparaison, Logique,\\Interface};
\node[draw=popgreen,fill=popgreenL,rounded corners,thick,align=center,minimum height=1.2cm,right=of tra] (act) {ACTIONNEUR\\Signal él. $\rightarrow$\\Action physique};
\draw[-{Stealth},very thick,popdark] (cap) -- node[above, align=center, font=\scriptsize] {signal\\d'entrée} (tra);
\draw[-{Stealth},very thick,popdark] (tra) -- node[above, align=center, font=\scriptsize] {signal de\\commande} (act);
\node[draw=poppurple,fill=poppurple!10,rounded corners,thick,align=center,below=0.8cm of tra, font=\small] (alim) {ALIMENTATION\\($V_{CC}$, $GND$)};
\draw[dashed,poppurple] (alim) -- (cap); \draw[dashed,poppurple] (alim) -- (tra); \draw[dashed,poppurple] (alim) -- (act);
\end{tikzpicture}
\end{center}
\legende{Structure générale d'une chaîne électronique : Capteur $\rightarrow$ Traitement $\rightarrow$ Actionneur. L'alimentation est omise sur les schémas fonctionnels mais indispensable.}
\end{popfigure}

\courssub{Exemple 1 : Commande d'une sirène par photorésistance et transistor (Détecteur crépusculaire)}

\textbf{Situation :} Allumer une sirène (ou réverbère) automatiquement la nuit.
\textbf{Chaîne :}
\begin{enumerate}
\item \textbf{Capteur :} Photorésistance $R_P$ + Résistance fixe $R_1$ en diviseur sous $V_{CC}$. Sortie : tension $U_S$ aux bornes de $R_1$. Nuit $\Rightarrow R_P \uparrow \Rightarrow U_S \downarrow$ (si $R_P$ en haut) ou $U_S \uparrow$ (si $R_P$ en bas). Choisissons $R_P$ en série avec $R_1$ au $+V_{CC}$, sortie sur $R_1$ vers masse : Nuit $\Rightarrow R_P \uparrow \Rightarrow U_S \downarrow$.
\item \textbf{Traitement :} \textbf{Comparateur} (ampli-op en comparateur ou transistor en commutation). Seuil $U_{seuil}$ fixe (ex. $\SI{2,5}{V}$). Si $U_S < U_{seuil}$ (nuit) $\Rightarrow$ Sortie traitement = Haut ($\approx V_{CC}$).
\item \textbf{Interface de puissance :} Transistor NPN + Relais. Base commandée par sortie traitement. Relais en collecteur avec diode roue libre.
\item \textbf{Actionneur :} Sirène (ou lampe) sur contact relais (circuit puissance \SI{220}{V} ou \SI{12}{V}).
\end{enumerate}

\begin{popfigure}
\begin{center}
\begin{tikzpicture}[scale=0.9, every node/.style={font=\small}]
\draw (0,5) node[left] {$+V_{CC}$} -- (7,5);
\draw (0,5) to[R, l_={$R_1$}] (0,2.5) to[photoresistor, l_={$R_P$}] (0,0);
\node[npn] (Q) at (3.5,2.5) {};
\draw (0,2.5) to[R, l={$R_B$}] (Q.B);
\draw (Q.C) to[L, l^=bobine du relais] (Q.C |- 0,5);
\draw (Q.C) -- ++(1.4,0) coordinate (d1) to[D, l_={$D$}] (d1 |- 0,5);
\draw (Q.E) -- (Q.E |- 0,0);
\draw (7,5) to[nos, l=contact du relais] (7,2.5) to[lamp, l=sirène] (7,0) -- (0,0);
\draw (3.5,0) node[ground] {};
\end{tikzpicture}
\end{center}
\legende{Détecteur crépusculaire : le jour, $R_P$ est faible, la tension base-émetteur est trop faible et le transistor est bloqué. La nuit, $R_P$ devient très grande, la tension de la base augmente, le transistor conduit, la bobine du relais est alimentée et son contact ferme le circuit de la sirène. La diode $D$ (roue libre) protège le transistor de la surtension à l'ouverture.}
\end{popfigure}

\courssub{Exemple 2 : Indicateur de niveau d'eau par électrodes et transistor}

\textbf{Situation :} Allumer une LED (ou buzzer) quand l'eau atteint un niveau haut dans un réservoir (Douala).
\textbf{Chaîne :}
\begin{enumerate}
\item \textbf{Capteur :} Deux électrodes $E_1$ (commun, tout en bas) et $E_2$ (niveau haut). L'eau (conductrice) relie $E_1$ et $E_2$ $\Rightarrow$ résistance $R_{eau} \approx \SI{10}{k\ohm}$--$\SI{100}{k\ohm}$. Air $\Rightarrow R = \infty$.
\item \textbf{Traitement + Interface :} Un unique transistor NPN. Base reliée à $E_2$ via $R_B$, Émetteur à la masse ($E_1$), Collecteur à LED + $R_{LED}$ sous $V_{CC}$.
\item \textbf{Fonctionnement :}
    \begin{itemize}
    \item Eau absente : circuit base ouvert $\Rightarrow I_B=0 \Rightarrow$ Transistor bloqué $\Rightarrow$ LED éteinte.
    \item Eau présente : courant de base $I_B = \frac{V_{CC} - U_{BE}}{R_B + R_{eau}}$ circule $\Rightarrow$ Transistor saturé $\Rightarrow$ LED allumée.
    \end{itemize}
\end{enumerate}

\begin{popfigure}
\begin{center}
\begin{tikzpicture}[scale=0.9, every node/.style={font=\small}]
\draw (0,5) node[left] {$+V_{CC} = \SI{9}{V}$} -- (5,5);
\node[npn] (Q) at (4.5,2) {};
\draw (Q.C) to[leD, l_={$D_1$}] ++(0,1.4) to[R, l_={$R_{\text{LED}}$}] (Q.C |- 0,5);
\draw (Q.E) -- (Q.E |- 0,0) node[ground] {};
\draw (0,5) -- (0,3.6) node[left] {électrode $E_2$};
\draw[popblue, dashed] (0,3.6) to[R, l_={$R_{\text{eau}}$}] (0,2) node[left] {électrode $E_1$};
\draw (0,2) to[R, l={$R_B$}] (Q.B);
\node[popblue, align=center] at (-2.2,2.8) {eau\\conductrice};
\end{tikzpicture}
\end{center}
\legende{Indicateur de niveau d'eau : l'eau ferme le circuit de base du transistor. $R_B$ limite $I_B$ ; $R_{eau}$ varie avec la conductivité de l'eau.}
\end{popfigure}

\courssub{Exemple 3 : Antenne et écouteur d'un poste radio (Réception AM/FM)}

\textbf{Chaîne de réception simplifiée :}
\begin{enumerate}
\item \textbf{Capteur (Antenne) :} L'onde radioélectrique (champ E variable) induit une tension alternative $u_{ant}(t)$ de très faible amplitude ($\mu\text{V}$--$\text{mV}$) aux bornes de l'antenne (bobine ferrite pour GO/PO, fil pour FM). Fréquence porteuse $f_p$ (ex. \SI{900}{kHz} AM, \SI{100}{MHz} FM).
\item \textbf{Traitement HF (Haute Fréquence) :}
    \begin{itemize}
    \item \emph{Sélection / Accord :} Circuit $LC$ accordé sur $f_p$ (filtre passe-bande) $\rightarrow$ sélectionne la station désirée, rejette les autres.
    \item \emph{Amplification HF :} Amplificateur à faible bruit (transistor/FET) amplifie $u_{ant}$.
    \end{itemize}
\item \textbf{Détection (Démodulation) :}
    \begin{itemize}
    \item \emph{AM (Modulation d'Amplitude) :} Détecteur à diode (redresseur + filtre RC) $\rightarrow$ récupère l'enveloppe = signal audio (BF).
    \item \emph{FM (Modulation de Fréquence) :} Discriminateur / Détecteur de rapport (PLL) $\rightarrow$ récupère variations de fréquence = signal audio (BF).
    \end{itemize}
\item \textbf{Traitement BF (Basse Fréquence) :} Amplificateur de puissance (étage de sortie push-pull ou classe D) adapte l'impédance et amplifie le signal audio pour l'écouteur.
\item \textbf{Actionneur (Écouteur / Haut-parleur) :} Convertit le courant audio alternatif en ondes sonores (force de Laplace sur bobine mobile).
\end{enumerate}

\begin{savaistu}[Radio communautaire au Cameroun]
Les radios rurales (ex. Radio Environnement, stations locales) utilisent des chaînes de réception simples. L'antenne filaire sur le toit capte l'onde ; un circuit $LC$ variable (condensateur variable = potentiomètre capacitif) accorde la station ; une diode germanium ($U_S \approx \SI{0,2}{V}$) détecte l'AM ; un casque piézoélectrique haute impédance (sans alimentation) restitue le son. C'est l'électronique « à l'état pur » : passive, robuste, économe en énergie.
\end{savaistu}

\begin{methode}[Analyser une chaîne électronique (Bac)]
\begin{enumerate}
\item \textbf{Identifier les trois blocs} : Capteur (quelle grandeur physique ? passif/actif ?), Traitement (quel composant actif ? transistor, ampli-op, logique ?), Actionneur (quelle action physique ?).
\item \textbf{Suivre le signal} : Écrire la nature du signal à chaque étape (grandeur physique $\rightarrow$ $R$ ou $U$ ou $I$ $\rightarrow$ $U$ logique $\rightarrow$ $I$ puissance $\rightarrow$ mouvement/son/lumière).
\item \textbf{Expliquer le rôle de chaque composant clé} : Pourquoi ce diviseur ? Pourquoi ce transistor (commutateur/ampli) ? Pourquoi ce relais (isolation/puissance) ? Pourquoi cette diode (roue libre/redressement) ?
\item \textbf{Préciser les conditions de fonctionnement} : Seuil de commutation, saturation, polarisation.
\end{enumerate}
\end{methode}

\begin{exoresolu}[Analyse d'une chaîne : Détecteur de gelée pour agriculture]
Un capteur de température (thermistance CTN $R_{CTN}$) est placé en diviseur de tension avec $R_1 = \SI{10}{k\ohm}$ sous $V_{CC}=\SI{5}{V}$. La tension $U_S$ aux bornes de $R_1$ entre sur l'entrée inverseuse d'un comparateur (ampli-op) dont l'entrée non-inverseuse reçoit une tension de référence $U_{ref} = \SI{2,5}{V}$. La sortie du comparateur commande la base d'un transistor NPN qui actionne un relais alimentant une résistance chauffante dans une serre.
Données : $R_{CTN}(0^\circ\text{C}) = \SI{32}{k\ohm}$, $R_{CTN}(10^\circ\text{C}) = \SI{10}{k\ohm}$.
\begin{enumerate}
\item Déterminer $U_S$ à $0^\circ\text{C}$ et $10^\circ\text{C}$.
\item Le chauffage doit s'allumer si $\theta < 2^\circ\text{C}$ (risque de gel). Le comparateur sort $V_{CC}$ si $U_- < U_+$. Le montage est-il correct ?
\item Calculer $R_B$ pour saturer le transistor ($I_{C,\text{sat}}=\SI{20}{mA}$, $\beta_{\min}=100$, $k=5$, $V_{cmd}=V_{CC}=\SI{5}{V}$).
\end{enumerate}
\tcblower
\textbf{1. Diviseur de tension :} $U_S = V_{CC} \frac{R_1}{R_1 + R_{CTN}}$.
$\theta = 0^\circ\text{C}$ : $U_S = 5 \times \frac{10}{10+32} = \frac{50}{42} \approx \SI{1,19}{V}$.
$\theta = 10^\circ\text{C}$ : $U_S = 5 \times \frac{10}{10+10} = \SI{2,5}{V}$.
\textbf{2. Logique :} À $0^\circ\text{C}$, $U_S = \SI{1,19}{V} < U_{ref} (\SI{2,5}{V}) \Rightarrow U_- < U_+ \Rightarrow$ Sortie comparateur = $V_{CC}$ (Haut) $\Rightarrow$ Transistor ON $\Rightarrow$ Relais ON $\Rightarrow$ Chauffage ON. \textbf{Correct.}
À $10^\circ\text{C}$, $U_S = \SI{2,5}{V} = U_{ref}$ (seuil). Pour $\theta > 10^\circ\text{C}$, $U_S > U_{ref} \Rightarrow$ Sortie = $0$ (Bas) $\Rightarrow$ Chauffage OFF.
\textbf{3. Base resistor :}
$I_{B,\text{min}} = 20/100 = \SI{0,2}{mA}$.
$I_B = 5 \times 0,2 = \SI{1,0}{mA}$.
$R_B = \frac{5 - 0,7}{1,0 \times 10^{-3}} = \SI{4,3}{k\ohm}$. Choix : $\SI{3,9}{k\ohm}$ ou $\SI{4,7}{k\ohm}$ (préférer $\SI{3,9}{k\ohm}$ pour marge).
\end{exoresolu}

\begin{exercice}[Chaîne de commande d'une pompe par niveau d'eau]
On souhaite commander une pompe \SI{220}{V} / \SI{500}{W} pour remplir un réservoir. Deux électrodes $E_{bas}$ et $E_{haut}$ détectent le niveau. La chaîne utilise : diviseur résistif (électrodes + $R_1$), comparateur (ampli-op), relais.
\begin{enumerate}
\item Compléter le schéma bloc de la chaîne (Capteur, Traitement, Actionneur, Alimentation).
\item Expliquer pourquoi un relais est indispensable entre le comparateur (sortie $\pm \SI{12}{V}$) et la pompe (\SI{220}{V}).
\item Si l'eau est peu conductrice ($R_{eau} = \SI{200}{k\ohm}$) et $R_1 = \SI{100}{k\ohm}$, $V_{CC}=\SI{12}{V}$, calculer la tension sur l'entrée du comparateur quand l'eau touche l'électrode. Est-ce suffisant pour un seuil à $\SI{6}{V}$ ?
\end{enumerate}
\end{exercice}

\begin{corrige}[Exercice : Chaîne de commande pompe]
\begin{enumerate}
\item \textbf{Blocs :}
    \begin{itemize}
    \item \textbf{Capteur :} Électrodes + $R_1$ (diviseur) $\rightarrow$ Grandeur : Niveau d'eau $\rightarrow$ Signal : Tension $U_S$.
    \item \textbf{Traitement :} Comparateur (Ampli-op) $\rightarrow$ Compare $U_S$ à $U_{ref}$ ($\SI{6}{V}$) $\rightarrow$ Sortie logique ($\pm \SI{12}{V}$).
    \item \textbf{Interface/Puissance :} Transistor driver + Relais (bobine \SI{12}{V}, contact \SI{220}{V}).
    \item \textbf{Actionneur :} Pompe \SI{220}{V}.
    \item \textbf{Alimentation :} Secteur \SI{220}{V} (pour pompe) + Transfo/Redresseur/Régulateur $\pm \SI{12}{V}$ (pour électronique).
    \end{itemize}
\item \textbf{Séparation galvanique :} Le comparateur fonctionne en TBT (\SI{12}{V}). La pompe est sur le secteur \SI{220}{V}. Le relais assure l'isolation électrique totale (sécurité personnes, protection électronique contre surtensions pompe). Un transistor seul ne supporte pas \SI{220}{V} et ne isole pas.
\item \textbf{Calcul :} $U_S = V_{CC} \frac{R_1}{R_1 + R_{eau}} = 12 \times \frac{100}{100+200} = 12 \times \frac{1}{3} = \SI{4}{V}$.
    $U_S = \SI{4}{V} < U_{ref} = \SI{6}{V}$. Si la logique est « Pompe ON si Niveau BAS (électrode haute à l'air) $\Rightarrow U_S$ haut », il faut inverser le diviseur ou les entrées du comparateur. Ici, eau présente $\Rightarrow U_S=\SI{4}{V}$ (Bas). Si on veut Pompe ON quand réservoir VIDE (électrode haute à l'air $\Rightarrow R_{eau}=\infty \Rightarrow U_S=\SI{12}{V}$), alors $U_S > U_{ref}$ déclenche la pompe. Le calcul montre qu'avec de l'eau, $U_S$ baisse. Le schéma fonctionne pour détecter l'absence d'eau (niveau bas).
\end{enumerate}
\end{corrige}

\begin{aretenir}[Synthèse de la leçon P13]
\begin{itemize}
\item \textbf{Dipôle commandé} : $R$, $I$ ou $U$ piloté par grandeur de commande (manuelle, électrique, physique).
\item \textbf{Rhéostat (2 bornes)} : règle $I$ (série). \textbf{Potentiomètre (3 bornes)} : règle $U$ (diviseur).
\item \textbf{Relais} : Isolation galvanique, commande puissance par signal faible.
\item \textbf{VDR} : Protection surtension ($R \downarrow$ brutal si $U > U_{seuil}$).
\item \textbf{Transistor NPN commutation} : Bloqué ($I_B=0$) / Saturé ($I_B > I_{C}/\beta_{min}$). $R_B$ calculée pour sursaturation ($k=5$--$10$). Diode roue libre sur charge inductive.
\item \textbf{Capteur passif} : $R=f(X)$ (CTN, LDR, Jauge). \textbf{Capteur actif} : Générateur $U=f(X)$ (Antenne, Micro, Piézo).
\item \textbf{Linéaire} (KTY, Jauge petit domaine) vs \textbf{Non linéaire} (CTN, LDR).
\item \textbf{Chaîne électronique} : Capteur $\rightarrow$ Traitement (Conditionnement, Comparaison, Logique) $\rightarrow$ Actionneur. Exemples : Sirène crépusculaire (LDR), Niveau eau (Électrodes), Radio (Antenne $\rightarrow$ Accord $\rightarrow$ Détection $\rightarrow$ Amplification BF $\rightarrow$ HP).
\end{itemize}
\end{aretenir}

\coursec{Étude du rhéostat : caractéristique U = f(I, R)}

Dans la section précédente, nous avons défini le \cle{dipôle commandé} : un dipôle dont la caractéristique dépend d'une grandeur de commande. Le \cle{rhéostat} est le plus ancien et le plus simple d'entre eux : c'est un \textbf{conducteur ohmique dont la résistance est réglable manuellement}, par déplacement d'un curseur le long d'un fil résistant bobiné. Commandé à la main, il permet de piloter un circuit — luminosité d'une lampe, vitesse d'un ventilateur — sans jamais le débrancher. Étudions-le en profondeur : c'est le premier exemple de caractéristique $U = f(I, R)$ que le programme te demande de savoir interpréter.

\courssub{Le rhéostat en série : un limiteur de courant}

\begin{definition}[Rhéostat]
Un \cle{rhéostat} est un conducteur ohmique à résistance variable, constitué d'un fil résistant (alliage nickel-chrome) enroulé sur un support isolant cylindrique, parcouru par un \textbf{curseur} mobile $C$. Utilisé entre une extrémité $A$ et le curseur $C$, il présente une résistance $R$ réglable entre $0$ et une valeur maximale $R_{\max}$ (typiquement de quelques ohms à quelques centaines d'ohms), tout en supportant un courant maximal $I_{\max}$ fixé par le constructeur.
\end{definition}

Pour commander l'intensité dans un dipôle récepteur (lampe, moteur), on monte le rhéostat \textbf{en série} avec ce dipôle. Le montage d'étude est le suivant : un générateur de f.é.m. $E$, un ampèremètre, la lampe et le rhéostat forment une maille unique ; un voltmètre mesure la tension $U$ aux bornes du rhéostat.

\begin{popfigure}
\begin{center}
\begin{circuitikz}[european, scale=1.0, transform shape]
\draw (0,0) to[V=$E$] (0,3)
      to[ammeter] (2.5,3)
      to[lamp, l_={$L$}] (5.5,3)
      to[vR, l_={$R$}] (8,3)
      -- (8,0) -- (0,0);
\draw (8,3) -- (9.5,3) to[voltmeter] (9.5,0) -- (8,0);
\draw[->, thick, poporange] (6.2,2.4) -- (7.4,2.4) node[midway, above]{$i$};
\end{circuitikz}
\end{center}
\legende{Montage du rhéostat $R$ en série avec une lampe $L$ alimentée par un générateur $E$. L'ampèremètre mesure $i$, le voltmètre mesure la tension $U$ aux bornes du rhéostat.}
\end{popfigure}

Appliquons la loi des mailles (loi d'additivité des tensions) à ce circuit série, la lampe étant assimilée à un conducteur ohmique de résistance $r$ :
\[ E = U_{\text{lampe}} + U = r\,i + R\,i = (r + R)\,i \qquad\Longrightarrow\qquad i = \frac{E}{r + R}. \]

\begin{propriete}[Rôle du rhéostat en série]
Dans un circuit série alimenté par un générateur de f.é.m. $E$ et contenant un récepteur de résistance $r$, l'intensité vaut
\[ i = \frac{E}{r + R}. \]
Augmenter $R$ \textbf{diminue} l'intensité $i$ dans tout le circuit : le rhéostat joue le rôle de \textbf{limiteur (ou régleur) de courant}. C'est une conséquence directe de la loi d'Ohm et de la loi des mailles.
\end{propriete}

Vérification dimensionnelle : $[E] = \text{V}$, $[r + R] = \Omega$, donc $[i] = \text{V}/\Omega = \text{A}$ : la relation est homogène.

\begin{attention}[Protéger avant de mesurer]
Au démarrage d'une expérience, on règle toujours le rhéostat à sa valeur \textbf{maximale} $R_{\max}$, puis on diminue progressivement. Si l'on partait de $R = 0$, l'intensité initiale $i = E/r$ pourrait dépasser le calibre de l'ampèremètre ou griller la lampe. C'est un réflexe de TP exigé au Baccalauréat.
\end{attention}

\courssub{Caractéristique U = f(I) à résistance R fixée}

Fixons le curseur : le rhéostat se comporte alors comme un conducteur ohmique ordinaire de résistance $R$. En faisant varier l'intensité $I$ qui le traverse (par exemple grâce à un générateur réglable) et en relevant la tension $U$ à ses bornes, on obtient sa caractéristique.

\begin{propriete}[Caractéristique du rhéostat réglé à R]
Pour une position fixée du curseur, la caractéristique $U = f(I)$ du rhéostat est la \textbf{droite passant par l'origine}
\[ U = R\,I, \]
de pente égale à la résistance $R$ choisie. C'est la traduction graphique immédiate de la \textbf{loi d'Ohm} : $U$ et $I$ sont proportionnelles, le coefficient de proportionnalité étant $R$.
\end{propriete}

\emph{Démonstration (immédiate mais exigible)} : le rhéostat est un conducteur ohmique ; par définition, il vérifie la loi d'Ohm $U = R I$. La fonction $I \mapsto R I$ est linéaire ; sa représentation graphique est donc une droite passant par l'origine, de coefficient directeur $\dfrac{\Delta U}{\Delta I} = R$. $\blacksquare$

\courssub{Réseau de caractéristiques U = f(I, R) : un faisceau de droites}

Faisons maintenant varier la grandeur de \textbf{commande}, c'est-à-dire la position du curseur, donc $R$. Pour chaque valeur $R_1$, $R_2$, $R_3$\ldots on obtient une nouvelle droite. L'ensemble forme un \textbf{faisceau de droites concourantes en l'origine}, paramétré par $R$ : c'est la signature graphique d'un dipôle \emph{commandé}.

\begin{popfigure}
\begin{center}
\begin{tikzpicture}
\begin{axis}[
    width=0.85\linewidth, height=7cm,
    xlabel={$I$ (A)}, ylabel={$U$ (V)},
    xmin=0, xmax=2.2, ymin=0, ymax=13,
    xtick={0,0.5,1,1.5,2}, ytick={0,2,4,6,8,10,12},
    grid=both, grid style={popline!40},
    legend style={at={(0.03,0.97)}, anchor=north west, font=\small},
    axis line style={popdark},
]
\addplot[popcurve, popblue, domain=0:2, samples=2, thick] {2*x};
\addlegendentry{$R_1 = \SI{2}{\ohm}$}
\addplot[popcurve, poporange, domain=0:2, samples=2, thick] {4*x};
\addlegendentry{$R_2 = \SI{4}{\ohm}$}
\addplot[popcurve, poppurple, domain=0:2, samples=2, thick] {6*x};
\addlegendentry{$R_3 = \SI{6}{\ohm}$}
\addplot[mark=*, poppink, only marks] coordinates {(1,6)};
\node[popdark, anchor=west, font=\small] at (axis cs:1.05,6) {$M(1\ \text{A}\,;\,6\ \text{V})$};
\draw[dashed, popmuted] (axis cs:1,0) -- (axis cs:1,6) -- (axis cs:0,6);
\end{axis}
\end{tikzpicture}
\end{center}
\legende{Réseau de caractéristiques $U = R\,I$ d'un rhéostat pour trois positions du curseur ($R_1 < R_2 < R_3$). Plus $R$ est grande, plus la droite est pentue. Le point de fonctionnement $M$ appartient à la droite $R_3 = U/I = \SI{6}{\ohm}$.}
\end{popfigure}

\begin{aretenir}[Lecture du faisceau de droites]
\begin{itemize}
\item Toutes les droites passent par l'origine : si $I = 0$, alors $U = 0$, quelle que soit $R$.
\item La pente de chaque droite est égale à la résistance correspondante : $R_1 < R_2 < R_3 \iff$ droites de plus en plus inclinées vers l'axe des tensions.
\item \textbf{À intensité $I$ imposée}, $U = R I$ augmente quand $R$ augmente.
\item \textbf{À tension $U$ imposée}, $I = U/R$ diminue quand $R$ augmente : c'est précisément le rôle de limitation du courant.
\end{itemize}
\end{aretenir}

\begin{methode}[Déterminer R à partir d'un point de fonctionnement]
On donne un point de fonctionnement $M(I_0\,;\,U_0)$ du rhéostat sur le réseau $U = f(I, R)$.
\begin{enumerate}
\item Relever les coordonnées du point : $I_0$ en ampères, $U_0$ en volts.
\item Calculer le quotient $R = \dfrac{U_0}{I_0}$ (loi d'Ohm) : c'est la résistance du rhéostat pour cette position du curseur.
\item Vérifier graphiquement que $M$ appartient bien à la droite de pente $R$ passant par l'origine.
\end{enumerate}
\emph{Exemple} : pour le point $M(\SI{1}{A}\,;\,\SI{6}{V})$ de la figure, $R = \dfrac{6}{1} = \SI{6}{\ohm}$ : le point appartient à la caractéristique $R_3$.
\end{methode}

\begin{exoresolu}[Régler le rhéostat pour protéger une lampe]
\textbf{Énoncé.} Un générateur de f.é.m. $E = \SI{12}{V}$ (résistance interne négligeable) alimente en série une lampe de résistance $r = \SI{4}{\ohm}$ et un rhéostat de résistance réglable $R$.
\begin{enumerate}
\item Quelle est l'intensité du courant lorsque $R = 0$ ? La lampe, qui supporte au plus $I_{\max} = \SI{1}{A}$, est-elle en danger ?
\item Déterminer la valeur de $R$ qui limite l'intensité à $I = \SI{1}{A}$.
\item Calculer alors la tension $U$ aux bornes du rhéostat et la puissance qu'il dissipe.
\end{enumerate}
\tcblower
\textbf{Solution.}
\begin{enumerate}
\item Loi des mailles : $i = \dfrac{E}{r + R}$. Pour $R = 0$ : $i = \dfrac{12}{4} = \SI{3}{A}$. C'est \textbf{trois fois} l'intensité maximale admissible : la lampe grille immédiatement.
\item On impose $I = \SI{1}{A}$ :
\[ I = \frac{E}{r + R} \iff r + R = \frac{E}{I} \iff R = \frac{E}{I} - r = \frac{12}{1} - 4 = \SI{8}{\ohm}. \]
\item Tension aux bornes du rhéostat : $U = R I = 8 \times 1 = \SI{8}{V}$ (la lampe reçoit les $\SI{4}{V}$ restants : $8 + 4 = 12 = E$, la loi des mailles est vérifiée). Puissance dissipée par effet Joule :
\[ P = U\,I = R\,I^2 = 8 \times 1^2 = \SI{8}{W}. \]
Le rhéostat doit donc supporter au moins $\SI{8}{W}$ : c'est pourquoi les rhéostats de laboratoire sont de gros cylindres bobinés, conçus pour évacuer la chaleur.
\end{enumerate}
\end{exoresolu}

\courssub{Applications : gradateur de lumière et variateur de vitesse}

Le principe « résistance série réglable » se retrouve partout dans la vie quotidienne :

\begin{itemize}
\item \textbf{Gradateur de lumière} : dans une lampe de chevet à variateur, la rotation du bouton déplace le curseur d'un rhéostat (ou, dans les versions modernes, règle un composant électronique équivalent). Augmenter $R$ diminue $I$, donc la puissance $P_L = r I^2$ de la lampe : l'éclat baisse.
\item \textbf{Variateur de vitesse d'un ventilateur} : les ventilateurs de plafond très répandus à Douala et Yaoundé possèdent souvent un boîtier à plusieurs positions. Chaque position insère une résistance différente en série avec le moteur : plus la résistance insérée est grande, plus le courant (et donc la vitesse de rotation) est faible.
\item \textbf{Rhéostat de démarrage des moteurs} : dans les ateliers, on insère un rhéostat au démarrage d'un moteur pour limiter l'appel de courant, puis on le réduit progressivement à zéro.
\end{itemize}

\begin{savaistu}[Pourquoi les gradateurs modernes ne chauffent plus]
Le gradateur à rhéostat gaspille de l'énergie : dans l'exercice résolu, $\SI{8}{W}$ sur les $\SI{12}{W}$ fournis par le générateur sont perdus en chaleur dans le rhéostat ! Les gradateurs modernes utilisent un \textbf{triac}, un interrupteur électronique ultra-rapide qui découpe la tension alternative : presque aucune énergie n'est dissipée. C'est le même souci d'efficacité énergétique qui guide ENEO et SONATREL dans la modernisation du réseau camerounais.
\end{savaistu}

\courssub{Complément : le montage potentiométrique}

Le montage série présente une limitation fondamentale : l'intensité vaut $i = \dfrac{E}{r + R}$, qui ne s'annule jamais, même pour $R$ très grande. Pour obtenir une tension réglable \textbf{de 0 à $E$}, on utilise le rhéostat en \textbf{potentiomètre} : ses deux extrémités $A$ et $B$ sont branchées aux bornes du générateur, et l'on prélève la tension $U_S$ entre le curseur $C$ et l'extrémité $B$.

\begin{popfigure}
\begin{center}
\begin{circuitikz}[european, scale=1.0, transform shape]
\draw (0,0) to[V=$E$] (0,3) -- (2,3);
\draw (2,3) to[R, l_={$R_{AC}$}] (4,3);
\draw (4,3) to[R, l_={$R_{CB}$}] (6,3) -- (6,0) -- (0,0);
\draw (4,3) -- (4,4.2) -- (7.5,4.2);
\draw (7.5,4.2) to[voltmeter, l_={$U_S$}] (7.5,0) -- (6,0);
\fill[poppink] (4,3) circle (2pt) node[below left, popdark]{$C$};
\node[popdark] at (2,3.5){$A$};
\node[popdark] at (6.35,3.2){$B$};
\end{circuitikz}
\end{center}
\legende{Montage potentiométrique : la résistance totale $AB$ du rhéostat est aux bornes du générateur ; la tension de sortie $U_S$ est prélevée entre le curseur $C$ et l'extrémité $B$.}
\end{popfigure}

Lorsque le circuit de sortie ne débite presque aucun courant (voltmètre, entrée d'un amplificateur), les deux portions $R_{AC}$ et $R_{CB}$ sont parcourues par le même courant et forment un \textbf{pont diviseur de tension} :
\[ U_S = \frac{R_{CB}}{R_{AC} + R_{CB}}\,E = \frac{R_{CB}}{R_{\max}}\,E. \]

\begin{propriete}[Tension de sortie du potentiomètre]
En montage potentiométrique à vide, la tension de sortie est proportionnelle à la fraction de résistance prélevée :
\[ U_S = \frac{R_{CB}}{R_{\max}}\,E, \qquad 0 \leq U_S \leq E. \]
Curseur en $B$ : $U_S = 0$ ; curseur en $A$ : $U_S = E$. La tension de sortie balaye \textbf{toute la plage de 0 à $E$}.
\end{propriete}

\begin{attention}[Rhéostat série ou potentiomètre : ne pas confondre]
\begin{itemize}
\item En \textbf{montage série} (2 bornes utilisées), le rhéostat règle le \textbf{courant} entre $i_{\max} = E/r$ et une valeur minimale non nulle : il \textbf{ne peut jamais annuler le courant}.
\item En \textbf{montage potentiométrique} (3 bornes utilisées), il règle la \textbf{tension} de sortie et peut l'annuler : $U_S = 0$ est atteint lorsque le curseur est en $B$.
\end{itemize}
Au Baccalauréat, lis attentivement l'énoncé : le nombre de bornes connectées trahit le montage utilisé, et prétendre qu'un rhéostat série peut éteindre complètement la lampe est une erreur classique sanctionnée.
\end{attention}

\begin{exercice}[Point de fonctionnement sur un réseau de caractéristiques]
Un rhéostat est traversé par un courant d'intensité $I = \SI{0,50}{A}$ et la tension à ses bornes vaut $U = \SI{7,5}{V}$.
\begin{enumerate}
\item Déterminer la résistance $R$ correspondant à cette position du curseur.
\item Le curseur est déplacé de sorte que la résistance double. À intensité inchangée, que vaut la nouvelle tension ?
\item Le réseau $U = f(I, R)$ de ce rhéostat comporte les droites $R = \SI{5}{\ohm}$, $\SI{10}{\ohm}$, $\SI{15}{\ohm}$. Sur quelle droite se trouve le point de fonctionnement initial ?
\end{enumerate}
\end{exercice}

\begin{corrige}[Exercice : point de fonctionnement]
\begin{enumerate}
\item Loi d'Ohm : $R = \dfrac{U}{I} = \dfrac{7,5}{0,50} = \SI{15}{\ohm}$.
\item Nouvelle résistance $R' = 2R = \SI{30}{\ohm}$. À $I = \SI{0,50}{A}$ : $U' = R' I = 30 \times 0,50 = \SI{15}{V}$. La tension double, conformément à la proportionnalité $U = RI$ à $I$ fixée.
\item Le point $(\SI{0,50}{A}\,;\,\SI{7,5}{V})$ vérifie $U/I = \SI{15}{\ohm}$ : il appartient à la droite $R = \SI{15}{\ohm}$ du réseau.
\end{enumerate}
\end{corrige}

\begin{aretenir}[L'essentiel sur le rhéostat]
\begin{itemize}
\item Le rhéostat est un \textbf{conducteur ohmique à résistance réglable manuellement} ; monté en série, il limite l'intensité : $i = \dfrac{E}{r + R}$.
\item À $R$ fixée, sa caractéristique est la droite $U = RI$ passant par l'origine (loi d'Ohm) ; le réseau $U = f(I, R)$ est un \textbf{faisceau de droites} dont la pente croît avec $R$.
\item À $I$ imposée, $U$ augmente avec $R$ ; à $U$ imposée, $I$ diminue quand $R$ augmente.
\item En montage \textbf{potentiométrique}, la tension de sortie $U_S = \dfrac{R_{CB}}{R_{\max}}E$ varie de $0$ à $E$ : seul ce montage permet d'annuler la grandeur réglée.
\end{itemize}
\end{aretenir}

Le rhéostat est commandé \emph{à la main}. Mais comment commander un dipôle par la température d'un local, par la lumière ambiante ou par une pression, sans intervention humaine ? C'est précisément l'objet de la section suivante, consacrée aux \textbf{capteurs} : thermistances, photorésistances et autres traducteurs du monde physique en grandeurs électriques.

\coursec{Capteurs de température, de lumière, de pression et d'ondes}

Dans la section précédente, nous avons étudié le rhéostat : un dipôle dont la résistance est commandée \emph{manuellement}, par le déplacement d'un curseur. La caractéristique $U = f(I, R)$ nous a montré que la résistance joue le rôle de \cle{grandeur de commande}. Mais dans la vie quotidienne, on souhaite souvent que ce soit la \emph{nature} elle-même qui commande le circuit : la température d'une chambre, la lumière du jour, le niveau d'eau d'un forage. C'est exactement le rôle des \cle{capteurs} : ce sont des dipôles dont une grandeur électrique (résistance, tension, courant) varie sous l'effet d'une grandeur physique extérieure. Nous allons maintenant passer en revue les grandes familles de capteurs au programme, en insistant sur l'exploitation de leurs courbes caractéristiques, compétence exigible au Baccalauréat.

\courssub{Capteurs de température : thermistance CTN et sonde KTY}

\textbf{Observation.} Le thermomètre électronique que l'infirmière place sous ton aisselle affiche $37{,}2~^{\circ}$C en quelques secondes. Aucun mercure, aucun cadran : à l'intérieur, un petit composant voit sa résistance varier avec la température, et un circuit électronique traduit cette variation en un nombre affiché.

\begin{definition}[Thermistance CTN]
Une \cle{thermistance CTN} (à \emph{coefficient de température négatif}) est un dipôle résistif dont la résistance \textbf{diminue lorsque la température augmente}. Elle est fabriquée à partir d'oxydes métalliques frittés (matériaux semi-conducteurs).
\end{definition}

\begin{definition}[Sonde KTY]
Une \cle{sonde KTY} est une résistance au silicium dont la résistance \textbf{augmente de façon quasi linéaire} lorsque la température augmente. On dit qu'elle se comporte comme une thermistance à coefficient de température positif (CTP).
\end{definition}

\begin{propriete}[Évolutions des résistances (admise)]
Pour une thermistance CTN, la résistance décroît \emph{fortement} et \emph{non linéairement} avec la température : une loi approchée est
\[ R(\theta) = R_0 \, e^{\,B\left(\frac{1}{T} - \frac{1}{T_0}\right)}, \]
où $T$ est la température absolue en kelvin ($T = \theta + 273$ avec $\theta$ en degrés Celsius), $T_0$ une température de référence et $B$ une constante caractéristique du matériau, homogène à une température (l'exponentielle est bien sans dimension). Pour une sonde KTY, la résistance croît de façon quasi linéaire : $R(\theta) \approx R_{25}\left[1 + \alpha(\theta - 25)\right]$ avec $\alpha \approx 7{,}9 \times 10^{-3}~^{\circ}$C$^{-1}$. Cette propriété est admise ; seule son exploitation graphique est exigible.
\end{propriete}

\begin{experience}[Tracé de la caractéristique $R = f(\theta)$ d'une CTN]
\textbf{Matériel :} une thermistance CTN, un ohmmètre, un bain-marie, un thermomètre, un agitateur.

\textbf{Protocole :} plonger la thermistance et le thermomètre dans l'eau. Relever $R$ pour différentes températures $\theta$ de $20~^{\circ}$C à $80~^{\circ}$C, en homogénéisant le bain avant chaque mesure et en attendant la stabilisation de la température.

\textbf{Observations :} $R$ diminue rapidement quand $\theta$ augmente ; la diminution est d'autant plus lente que la température est élevée.

\textbf{Interprétation :} la courbe obtenue est décroissante et \emph{non rectiligne} : la CTN est un capteur \textbf{non linéaire}. La même expérience avec une sonde KTY donnerait une droite croissante : la KTY est un capteur quasi \textbf{linéaire}.
\end{experience}

\begin{popfigure}
\begin{center}
\begin{tikzpicture}
\begin{axis}[
  width=0.92\linewidth, height=6.5cm,
  xlabel={$\theta$ ($^{\circ}$C)}, ylabel={$R$ (k$\Omega$)},
  xmin=0, xmax=100, ymin=0, ymax=12,
  xtick={0,20,40,60,80,100}, ytick={0,2,4,6,8,10,12},
  grid=both, grid style={popline!40},
  legend style={at={(0.97,0.97)}, anchor=north east, font=\small}
]
\addplot[popcurve,domain=5:100,samples=200] {12*exp(-0.045*(x-10))};
\addlegendentry{CTN}
\addplot[popcurve,poporange,domain=0:100,samples=50] {1 + 0.008*x};
\addlegendentry{KTY}
\end{axis}
\end{tikzpicture}
\end{center}
\legende{Caractéristiques $R = f(\theta)$ : la CTN (courbe décroissante, non linéaire) et la KTY (droite croissante, quasi linéaire).}
\end{popfigure}

\begin{exemplebox}[La CTN est-elle linéaire ?]
Une CTN vaut $R_1 = 10~\text{k}\Omega$ à $\theta_1 = 20~^{\circ}$C et $R_2 = 2~\text{k}\Omega$ à $\theta_2 = 60~^{\circ}$C. Si elle était linéaire, le taux de variation serait constant :
\[ \frac{\Delta R}{\Delta \theta} = \frac{2 - 10}{60 - 20} = \frac{-8}{40} = -0{,}20~\text{k}\Omega/^{\circ}\text{C}. \]
On prédirait alors à $40~^{\circ}$C : $R = 10 - 0{,}20 \times 20 = 6~\text{k}\Omega$. Or la mesure réelle donne environ $4{,}5~\text{k}\Omega$ : la variation n'est \textbf{pas proportionnelle} à la température. De plus, le rapport $R/\theta$ n'est pas constant ($10/20 = 0{,}5$ mais $2/60 \approx 0{,}033$). La CTN est donc un capteur \textbf{non linéaire}.
\end{exemplebox}

\courssub{Capteurs de lumière : photorésistance et photodiode}

\textbf{Observation.} À Yaoundé, certains lampadaires s'allument tout seuls à la tombée de la nuit. Personne n'actionne d'interrupteur : un capteur de lumière détecte la décroissance de l'éclairement et commande l'allumage.

\begin{definition}[Photorésistance (LDR)]
Une \cle{photorésistance} (ou LDR, \emph{Light Dependent Resistor}) est un dipôle dont la résistance \textbf{diminue lorsque l'éclairement $\Phi$ augmente}. Dans l'obscurité, sa résistance est très grande (plusieurs M$\Omega$) ; en pleine lumière, elle tombe à quelques centaines d'ohms.
\end{definition}

\begin{definition}[Photodiode]
Une \cle{photodiode} est une diode utilisée en polarisation inverse : le courant inverse qui la traverse est \textbf{proportionnel à l'éclairement} qu'elle reçoit. Elle convertit directement un flux lumineux en courant électrique.
\end{definition}

\textbf{Caractéristique $U = f(I, \Phi)$ d'une photorésistance.} À éclairement $\Phi$ fixé, la photorésistance se comporte comme un conducteur ohmique : $U = R(\Phi)\,I$, donc sa caractéristique est une \emph{droite passant par l'origine} dont la pente vaut $R(\Phi)$. Quand l'éclairement augmente, $R$ diminue, donc la pente de la droite diminue : on obtient un \cle{réseau de droites} paramétré par $\Phi$.

\begin{popfigure}
\begin{center}
\begin{tikzpicture}
\begin{axis}[
  width=0.92\linewidth, height=6.5cm,
  xlabel={$I$ (mA)}, ylabel={$U$ (V)},
  xmin=0, xmax=10, ymin=0, ymax=10,
  xtick={0,2,4,6,8,10}, ytick={0,2,4,6,8,10},
  grid=both, grid style={popline!40},
  legend style={at={(0.03,0.97)}, anchor=north west, font=\small}
]
\addplot[popcurve,domain=0:10,samples=20] {1.5*x};
\addlegendentry{$\Phi_1$ (faible)}
\addplot[popcurve,poporange,domain=0:10,samples=20] {0.8*x};
\addlegendentry{$\Phi_2$ (moyen)}
\addplot[popcurve,popgreen,domain=0:10,samples=20] {0.3*x};
\addlegendentry{$\Phi_3$ (fort)}
\end{axis}
\end{tikzpicture}
\end{center}
\legende{Réseau de caractéristiques $U = f(I, \Phi)$ d'une photorésistance pour $\Phi_1 < \Phi_2 < \Phi_3$ : la pente $R(\Phi)$ diminue quand l'éclairement augmente.}
\end{popfigure}

\begin{attention}[Ne pas inverser les évolutions]
Deux sens de variation à mémoriser sans se tromper :
\begin{itemize}
\item Thermistance CTN : $\theta$ augmente $\Rightarrow$ $R$ \textbf{diminue} ;
\item Photorésistance : $\Phi$ augmente $\Rightarrow$ $R$ \textbf{diminue}.
\end{itemize}
Une erreur classique au Bac consiste à affirmer que la résistance d'une LDR augmente avec la lumière. Astuce : plus le capteur est « excité » (chaud, éclairé), plus il « laisse passer » le courant, donc plus sa résistance est faible. Attention aussi à ne pas confondre photorésistance (sa \emph{résistance} varie) et photodiode (c'est le \emph{courant inverse} qui varie).
\end{attention}

\begin{popfigure}
\begin{center}
\begin{tikzpicture}[scale=1.0]
% Thermistance
\draw[thick] (0,0) rectangle (1.6,0.8);
\draw[thick] (-0.8,0.4) -- (0,0.4); \draw[thick] (1.6,0.4) -- (2.4,0.4);
\draw[thick,->] (2.0,-0.3) -- (-0.4,1.1);
\node[below] at (0.8,-0.5) {Thermistance};
% Photorésistance
\begin{scope}[xshift=4.5cm]
\draw[thick] (0,0) rectangle (1.6,0.8);
\draw[thick] (-0.8,0.4) -- (0,0.4); \draw[thick] (1.6,0.4) -- (2.4,0.4);
\draw[thick,->,poporange] (-0.5,1.5) -- (0.2,0.9);
\draw[thick,->,poporange] (0.3,1.5) -- (1.0,0.9);
\node[below] at (0.8,-0.5) {Photorésistance};
\end{scope}
% Photodiode
\begin{scope}[xshift=9cm]
\draw[thick] (-0.8,0.4) -- (0.2,0.4); \draw[thick] (1.0,0.4) -- (2.4,0.4);
\draw[thick] (0.2,0.4) -- (0.6,0.8) -- (0.6,0.0) -- cycle;
\draw[thick] (0.6,0.0) -- (0.6,0.8);
\draw[thick] (1.0,0.0) -- (1.0,0.8);
\draw[thick,->,poporange] (0.0,1.5) -- (0.45,0.95);
\draw[thick,->,poporange] (0.8,1.5) -- (1.25,0.95);
\node[below] at (0.8,-0.5) {Photodiode};
\end{scope}
\end{tikzpicture}
\end{center}
\legende{Symboles normalisés : thermistance (rectangle barré d'une flèche oblique), photorésistance (rectangle éclairé par deux flèches), photodiode (diode recevant des flèches lumineuses).}
\end{popfigure}

\courssub{Capteurs de pression}

\begin{definition}[Capteur de pression]
Un \cle{capteur de pression} convertit une pression (ou une force) en une grandeur électrique. Deux technologies au programme :
\begin{itemize}
\item la \cle{jauge de contrainte} : un fil conducteur collé sur une membrane se déforme sous l'effet de la pression ; or la résistance d'un fil vaut $R = \rho \dfrac{\ell}{S}$ : l'allongement ($\ell$ augmente, $S$ diminue) fait donc \textbf{augmenter $R$}. La déformation mécanique est traduite en variation de résistance ;
\item la \cle{capsule capacitive} : deux armatures dont l'une est une membrane flexible forment un condensateur ; la pression rapproche les armatures, et comme $C = \varepsilon \dfrac{S}{e}$, la diminution de la distance $e$ fait \textbf{augmenter la capacité $C$}. La pression est traduite en variation de capacité.
\end{itemize}
\end{definition}

Ces capteurs équipent les manomètres électroniques, les stations météorologiques (pression atmosphérique) et les balances électroniques (la masse posée déforme une jauge de contrainte).

\courssub{Capteurs d'ondes : antenne et microphone}

\begin{definition}[Antenne]
Une \cle{antenne} est un conducteur qui capte les ondes électromagnétiques (radio, télévision, téléphonie mobile). Le champ électrique variable de l'onde met en mouvement les électrons du conducteur : il apparaît aux bornes de l'antenne une \textbf{tension induite} alternative, image fidèle de l'onde reçue. C'est le phénomène d'induction électromagnétique.
\end{definition}

\begin{definition}[Microphone]
Un \cle{microphone} convertit une onde sonore en tension électrique. Dans le microphone électrodynamique, l'onde sonore fait vibrer une membrane solidaire d'une bobine placée dans le champ d'un aimant : le mouvement de la bobine y induit une tension alternative de même fréquence que le son. L'onde mécanique est ainsi transformée en signal électrique.
\end{definition}

\begin{savaistu}[Des thermistances partout autour de nous]
Les thermomètres médicaux électroniques utilisés dans les hôpitaux de Douala et de Yaoundé contiennent une thermistance CTN : sa résistance chute quand la température corporelle monte, et un microcontrôleur affiche la température en moins de dix secondes. Les climatiseurs, les réfrigérateurs et les ventilateurs automatiques utilisent aussi des CTN pour réguler la température. Quant aux antennes, elles sont au cœur du réseau : chaque relais MTN ou Orange visible sur les collines du Cameroun capte et émet des ondes électromagnétiques grâce à des panneaux d'antennes.
\end{savaistu}

\courssub{Capteur linéaire et capteur non linéaire}

\begin{definition}[Capteur linéaire / non linéaire]
Un capteur est dit \cle{linéaire} si la grandeur électrique de sortie $s$ est une fonction affine (en pratique proportionnelle) de la grandeur physique d'entrée $m$ : $s = a\,m + b$. Graphiquement, sa courbe caractéristique $s = f(m)$ est une \textbf{droite}. Sinon, le capteur est \cle{non linéaire}. La KTY et la photodiode sont (quasi) linéaires ; la CTN et la photorésistance sont non linéaires.
\end{definition}

\textbf{Conséquence pratique.} Avec un capteur linéaire, le signal électrique est proportionnel à la grandeur mesurée : l'étalonnage est simple (deux points suffisent) et la précision est uniforme sur toute la plage de mesure. Avec un capteur non linéaire, la sensibilité varie : une CTN est très sensible à basse température (forte pente) mais peu sensible à haute température, et l'étalonnage exige de nombreux points de mesure.

\begin{methode}[Décider de la linéarité d'un capteur]
\begin{enumerate}
\item Tracer (ou observer) la courbe caractéristique : grandeur électrique en ordonnée, grandeur physique en abscisse.
\item Si la courbe est une droite, le capteur est linéaire ; sinon, il est non linéaire.
\item Pour confirmer quantitativement, calculer le rapport $\dfrac{\Delta s}{\Delta m}$ (ou le rapport $s/m$) pour plusieurs couples de mesures : s'il est \textbf{constant}, il y a proportionnalité, donc linéarité.
\item Conclure en rédigeant : « la courbe est une droite / n'est pas une droite, donc le capteur est linéaire / non linéaire ».
\end{enumerate}
\end{methode}

\begin{exoresolu}[Exploitation de la caractéristique d'une photorésistance]
On mesure la résistance d'une photorésistance sous trois éclairements : $R_1 = 9~\text{k}\Omega$ pour $\Phi_1 = 100~\text{lx}$, $R_2 = 3~\text{k}\Omega$ pour $\Phi_2 = 500~\text{lx}$, $R_3 = 1~\text{k}\Omega$ pour $\Phi_3 = 1500~\text{lx}$.

1. La photorésistance est-elle un capteur linéaire ?

2. On l'insère dans un circuit sous tension $U = 6~\text{V}$ constante. Calculer l'intensité du courant pour chaque éclairement et conclure.
\tcblower
\textbf{1. Linéarité.} Calculons les taux de variation entre points successifs :
\[ \frac{\Delta R}{\Delta \Phi}\bigg|_{1\to 2} = \frac{3 - 9}{500 - 100} = -0{,}015~\text{k}\Omega/\text{lx} ; \qquad \frac{\Delta R}{\Delta \Phi}\bigg|_{2\to 3} = \frac{1 - 3}{1500 - 500} = -0{,}002~\text{k}\Omega/\text{lx}. \]
Les taux ne sont pas égaux : la variation n'est pas proportionnelle, la courbe $R = f(\Phi)$ n'est pas une droite. La photorésistance est un capteur \textbf{non linéaire}.

\textbf{2. Intensités.} D'après la loi d'Ohm, $I = \dfrac{U}{R}$ :
\[ I_1 = \frac{6}{9000} \approx 0{,}67~\text{mA} ; \qquad I_2 = \frac{6}{3000} = 2{,}0~\text{mA} ; \qquad I_3 = \frac{6}{1000} = 6{,}0~\text{mA}. \]
Quand l'éclairement augmente, $R$ diminue et le courant augmente : le circuit électronique associé pourra détecter cette augmentation de courant pour, par exemple, commander l'extinction d'un lampadaire au lever du jour.
\end{exoresolu}

\begin{aretenir}[L'essentiel de la section]
\begin{itemize}
\item CTN : $\theta$ augmente $\Rightarrow$ $R$ diminue, de façon forte et non linéaire ; KTY : $R$ croît quasi linéairement avec $\theta$.
\item Photorésistance : $\Phi$ augmente $\Rightarrow$ $R$ diminue ; sa caractéristique $U = f(I)$ à $\Phi$ fixé est une droite de pente $R(\Phi)$. Photodiode : courant inverse proportionnel à l'éclairement.
\item Capteurs de pression : la déformation modifie $R$ (jauge de contrainte) ou $C$ (capsule capacitive).
\item Antenne : onde électromagnétique $\rightarrow$ tension induite ; microphone : onde sonore $\rightarrow$ tension électrique.
\item Un capteur est linéaire si sa courbe caractéristique est une droite (signal proportionnel à la grandeur mesurée) ; sinon il est non linéaire.
\end{itemize}
\end{aretenir}

\coursec{Le transistor en commutation et les électrodes}

Dans la section précédente, nous avons rencontré toute une famille de capteurs : thermistances, photorésistances, antennes, microphones. Tous ont un point commun : ils délivrent une information électrique \emph{faible}. Une photorésistance éclairée fait varier une tension de quelques volts, mais elle est incapable d'alimenter directement une sirène, un moteur ou une lampe puissante. Il manque un \textbf{maillon intermédiaire} : un composant capable de recevoir l'ordre donné par le capteur et d'ouvrir ou de fermer le circuit de puissance de l'actionneur. Ce maillon, c'est le \cle{transistor}, véritable « interrupteur commandé électriquement ». C'est lui qui fait le lien entre le monde du capteur et celui de l'actionneur dans toute chaîne électronique.

\courssub{Présentation du transistor bipolaire NPN}

\textbf{Intuition.} Imagine un robinet dont l'ouverture serait commandée non pas par ta main, mais par un faible filet d'eau : plus le filet de commande coule, plus le robinet principal laisse passer d'eau. Le transistor fonctionne exactement ainsi, avec des courants électriques à la place de l'eau.

\begin{definition}[Transistor bipolaire NPN]
Le \cle{transistor bipolaire} est un composant électronique à trois bornes, appelées \cle{électrodes} :
\begin{itemize}
\item la \cle{base} $B$ : électrode de \textbf{commande} ;
\item le \cle{collecteur} $C$ : électrode d'\textbf{entrée} du courant principal ;
\item l'\cle{émetteur} $E$ : électrode de \textbf{sortie} du courant principal.
\end{itemize}
Dans un transistor \cle{NPN}, un petit courant injecté dans la base ($I_B$) autorise et contrôle un courant beaucoup plus grand circulant du collecteur vers l'émetteur ($I_C$). C'est un \cle{dipôle commandé électriquement} (on considère le dipôle collecteur--émetteur, commandé par la base).
\end{definition}

La loi des nœuds appliquée au transistor donne immédiatement la relation entre les trois courants :
\[ I_E = I_B + I_C. \]
Comme $I_B$ est très petit devant $I_C$ (typiquement cent fois plus petit), on a en pratique $I_E \approx I_C$.

\begin{popfigure}
\begin{center}
\begin{circuitikz}[european, scale=1.2]
% Transistor NPN
\node[npn] (T) at (0,0) {};
% Connexions et labels
\draw (T.B) -- ++(-1.5,0) node[left] {$B$ (base)};
\draw (T.C) -- ++(0,1.5) node[above] {$C$ (collecteur)};
\draw (T.E) -- ++(0,-1.5) node[below] {$E$ (émetteur)};
% Courants
\draw[-{Stealth}, popblue, very thick] (-1.0,0.5) -- (T.B) node[midway, above] {$I_B$};
\draw[-{Stealth}, poporange, very thick] (0,2.5) -- (T.C) node[midway, right] {$I_C$};
\draw[-{Stealth}, popgreen, very thick] (T.E) -- (0,-2.5) node[midway, right] {$I_E$};
% Tensions
\draw[-{Stealth}, poppurple, thick] (T.B) to[bend right=45] node[midway, left] {$V_{BE}$} (T.E);
\draw[-{Stealth}, popgold, thick] (T.C) to[bend left=45] node[midway, right] {$V_{CE}$} (T.E);
% Gain label
\node[popmuted, align=center] at (2.5,0) {$I_C = \beta I_B$\\(régime linéaire)};
\end{circuitikz}
\end{center}
\legende{Symbole normalisé du transistor NPN (style européen) : courants $I_B$, $I_C$, $I_E$ (avec $I_E = I_B + I_C$) et tensions $V_{BE}$ (base--émetteur) et $V_{CE}$ (collecteur--émetteur). La flèche sur l'émetteur indique le sens conventionnel du courant $I_E$ (sortant pour NPN).}
\end{popfigure}

\begin{attention}[La jonction base--émetteur se comporte comme une diode]
Entre la base et l'émetteur, le transistor présente une \textbf{jonction PN} : elle ne devient passante que si $V_{BE}$ atteint environ \SI{0,7}{V} (tension de seuil pour le silicium). En dessous, $I_B = 0$ et le transistor reste bloqué. Au-dessus, $V_{BE}$ reste pratiquement bloquée à \SI{0,7}{V} : c'est le courant $I_B$ qui augmente, pas la tension.
\end{attention}

\courssub{Le transistor comme interrupteur commandé : régimes bloqué et saturé}

\begin{propriete}[Gain en courant et régimes de fonctionnement (admise)]
En \cle{régime linéaire} (amplification analogique), le courant de collecteur est proportionnel au courant de base :
\[ I_C = \beta \cdot I_B, \]
où $\beta$ (sans dimension) est le \cle{gain en courant} du transistor, de l'ordre de $100$ à $300$ pour les transistors courants. Cette relation est \textbf{admise} au programme. En \cle{commutation} (logique, tout-ou-rien), on n'utilise que les deux états extrêmes :
\begin{itemize}
\item \textbf{Transistor bloqué} : $I_B = 0 \Rightarrow I_C = 0$ et $V_{CE} = E$ (tension d'alimentation). Le transistor se comporte comme un \textbf{interrupteur ouvert} entre $C$ et $E$ (résistance quasi infinie).
\item \textbf{Transistor saturé} : $I_B$ suffisamment grand $\Rightarrow V_{CE} \approx V_{CE,\text{sat}} \approx \SI{0,2}{V}$ et le courant $I_C$ n'est limité que par le circuit extérieur (résistance de charge, résistance interne du générateur). Le transistor se comporte comme un \textbf{interrupteur fermé} (résistance quasi nulle).
\end{itemize}
\end{propriete}

\begin{aretenir}[Condition de saturation]
Pour garantir la saturation, le courant de base doit vérifier :
\[ I_B \geq \frac{I_{C,\text{sat}}}{\beta}, \]
où $I_{C,\text{sat}}$ est le courant que le circuit du collecteur laisserait passer si le transistor était parfaitement fermé ($V_{CE} = 0$). En pratique, on choisit $I_B$ deux à cinq fois supérieur à cette valeur minimale pour une saturation franche et robuste (facteur de sursaturation).
\end{aretenir}

\begin{exemplebox}[Calcul du courant de base minimal et résistance de base]
On veut saturer un transistor de gain $\beta = 100$ dont le collecteur alimente un relais consommant $I_{C,\text{sat}} = \SI{50}{mA}$ sous \SI{12}{V}.
\begin{enumerate}
\item Courant de base minimal théorique : $I_{B,\min} = \dfrac{I_{C,\text{sat}}}{\beta} = \dfrac{\num{50e-3}}{100} = \SI{0,5}{mA}$.
\item On vise un courant de base réel $I_B = \SI{1}{mA}$ (facteur 2) pour une saturation franche.
\item La tension de commande (sortie capteur ou générateur) vaut $U = \SI{5}{V}$. La résistance de base $R_B$ doit limiter le courant :
\[ R_B = \frac{U - V_{BE}}{I_B} = \frac{5 - 0{,}7}{\num{1e-3}} = \SI{4,3}{k\Omega}. \]
On choisit la valeur normalisée la plus proche \textbf{inférieure} pour garantir $I_B \ge \SI{1}{mA}$ : $R_B = \SI{3,9}{k\Omega}$ (ou \SI{4,7}{k\Omega} si on accepte $I_B \approx \SI{0,9}{mA}$, encore bien au-dessus de \SI{0,5}{mA}).
\end{enumerate}
Le transistor réalise ici une \textbf{amplification de puissance} : une puissance de commande $P_B \approx V_{BE} I_B \approx \SI{0,7}{mW}$ pilote une puissance de charge $P_C \approx \SI{12}{V} \times \SI{50}{mA} = \SI{600}{mW}$.
\end{exemplebox}

\begin{attention}[Toujours protéger la base par une résistance]
Ne relie \textbf{jamais} la base directement au pôle $+$ d'un générateur de tension. Comme $V_{BE}$ reste bloquée à \SI{0,7}{V}, toute la tension excédentaire ($U - 0,7$) s'appliquerait sur la résistance interne quasi nulle de la jonction : le courant de base serait limité seulement par les résistances parasites, atteignant des valeurs destructrices (souvent $> \SI{100}{mA}$). Le transistor est \textbf{détruit instantanément}. La résistance de base $R_B$ est la « soupape de sécurité » qui fixe $I_B = (U - V_{BE})/R_B$.
\end{attention}

\courssub{Montage type : capteur en pont diviseur, actionneur sur le collecteur}

Le montage universel d'une chaîne de commande à transistor comporte deux étages distincts :
\begin{itemize}
\item \textbf{Étage de commande (entrée)} : le capteur (photorésistance, thermistance CTN, électrodes) est associé à une résistance fixe $R_1$ pour former un \cle{pont diviseur de tension} alimentant la base (éventuellement via une résistance de protection $R_B$). Quand la grandeur physique varie (lumière, température, niveau d'eau), la résistance du capteur varie, donc la tension de base $V_B$ varie, ce qui fait apparaître ou disparaître $I_B$.
\item \textbf{Étage de puissance (sortie)} : l'actionneur (lampe, LED, relais, sirène, moteur) est placé en série avec le collecteur, alimenté par la source de puissance $E$ (souvent plus élevée que la tension de commande). Il n'est traversé par un courant significatif que si le transistor est saturé.
\end{itemize}

\begin{popfigure}
\begin{center}
\begin{circuitikz}[european, scale=0.95, transform shape]
% Alimentation puissance (collecteur)
\draw (0,0) node[ground]{};
\draw (0,0) to[battery1, l_={$E_{\text{puissance}}$}] (0,4);
\draw (0,4) -- (2,4);
% Pont diviseur sur la base (commande)
\draw (2,4) to[R=$R_1$] (2,2.5);
\draw (2,2.5) to[photoresistor, l_={$R_{\text{capteur}}$}] (2,1) -- (2,0) -- (0,0);
% Résistance de base
\draw (2,2.5) -- (3.5,2.5) to[R=$R_B$] (5,2.5);
% Transistor
\node[npn, anchor=B] (T) at (5,2.5) {};
% Collecteur -> Actionneur (LED + Rlim)
\draw (T.C) -- (5,3.5) to[R=$R_{\text{lim}}$] (5,4) to[led, l_=LED] (3,4) -- (2,4);
% Emetteur -> Masse
\draw (T.E) -- (5,0) -- (2,0);
% Annotations
\node[popmuted, align=center] at (1,3.5) {\small Étage\\ \small puissance};
\node[popmuted, align=center] at (1,1.5) {\small Pont diviseur\\ \small (capteur + $R_1$)};
\node[popmuted, align=center] at (6.5,2.5) {\small Transistor\\ \small NPN};
\node[popblue] at (4.2,2.8) {$I_B$};
\draw[-{Stealth}, popblue, thick] (3.8,2.6) -- (4.6,2.6);
\node[poporange] at (5.5,3.2) {$I_C$};
\draw[-{Stealth}, poporange, thick] (5,3.4) -- (5,2.9);
\end{circuitikz}
\end{center}
\legende{Schéma principe d'une chaîne de commande : le capteur (ici photorésistance) dans un pont diviseur pilote la base ; l'actionneur (LED + résistance de limitation) est en série sur le collecteur. $R_B$ protège la base.}
\end{popfigure}

\begin{methode}[Analyser un montage à transistor en commutation (3 étapes)]
\begin{enumerate}
\item \textbf{État du capteur} : D'après les conditions physiques (lumière forte/faible, température haute/basse, niveau d'eau), déterminer si la résistance du capteur $R_{\text{cap}}$ est \textbf{grande} ou \textbf{petite} par rapport à $R_1$.
\item \textbf{Tension et courant de base} : Calculer la tension au point milieu du pont diviseur (tension à vide $V_{B0}$) :
\[ V_{B0} = E_{\text{commande}} \cdot \frac{R_{\text{bas}}}{R_1 + R_{\text{bas}}} \]
(attention à l'ordre : $R_{\text{bas}}$ est la résistance reliée à la masse).
\begin{itemize}
\item Si $V_{B0} < \SI{0,7}{V}$ : la jonction BE est bloquée, $I_B = 0$.
\item Si $V_{B0} > \SI{0,7}{V}$ : la jonction BE conduit. Le courant de base est approximativement $I_B \approx \dfrac{V_{B0} - 0,7}{R_B}$ (si $R_B$ existe) ou limité par la résistance interne du pont. Vérifier la condition de saturation $I_B \geq I_{C,\text{sat}}/\beta$.
\end{itemize}
\item \textbf{État du transistor et de l'actionneur} : Conclure.
\begin{itemize}
\item $I_B = 0$ ou insuffisant $\rightarrow$ Transistor \textbf{bloqué} $\rightarrow$ Interrupteur C-E \textbf{ouvert} $\rightarrow$ Actionneur \textbf{au repos} (éteint).
\item $I_B \geq I_{C,\text{sat}}/\beta$ $\rightarrow$ Transistor \textbf{saturé} ($V_{CE} \approx \SI{0,2}{V}$) $\rightarrow$ Interrupteur C-E \textbf{fermé} $\rightarrow$ Actionneur \textbf{alimenté} (allumé).
\end{itemize}
\end{enumerate}
\end{methode}

\courssub{Les électrodes : capteur de niveau d'un liquide conducteur}

\begin{definition}[Électrodes (capteur de niveau)]
Des \cle{électrodes} sont deux conducteurs métalliques (souvent inox ou graphite) plongés dans un liquide, sans contact direct entre eux. Elles forment un capteur \textbf{tout-ou-rien} (non linéaire) :
\begin{itemize}
\item \textbf{Niveau bas} : le liquide ne recouvre pas les électrodes. L'air (isolant) est entre elles. Circuit de base \textbf{ouvert} : $I_B = 0$.
\item \textbf{Niveau haut} : le liquide recouvre les deux électrodes. L'eau courante (eau de forage, eau de pluie, eau de rivière) contient des sels minéraux dissous $\rightarrow$ ions mobiles $\rightarrow$ \textbf{solution conductrice}. Le liquide joue le rôle d'une résistance $R_{\text{eau}}$ (quelques $k\Omega$ à quelques $100\,k\Omega$) qui \textbf{ferme le circuit de base} : $I_B$ apparaît.
\end{itemize}
\end{definition}

\begin{attention}[Eau pure vs eau courante]
L'eau \emph{distillée} (pure, déionisée) est un excellent isolant ($R \to \infty$) : le montage \textbf{ne fonctionne pas} avec de l'eau pure. Il fonctionne grâce à la conductivité ionique de l'eau « réelle » (eau du robinet, eau de puits, eau de mer). C'est un capteur de \textbf{présence d'eau conductrice}.
\end{attention}

C'est le principe de l'\textbf{indicateur de niveau d'eau} d'un réservoir (château d'eau, cuve de forage, bac de récupération) : deux électrodes fixées à la hauteur voulue commandent, via un transistor, une LED témoin, une alarme sonore (buzzer) ou le démarrage d'une pompe (via un relais de puissance).

\begin{exoresolu}[Veilleuse automatique à photorésistance]
Un montage utilise un pont diviseur formé de $R_1 = \SI{100}{k\Omega}$ (en haut, reliée au $+$) et d'une photorésistance $R_\phi$ (en bas, reliée à la masse), alimenté sous $E = \SI{9}{V}$. La base du transistor ($\beta = 200$) est reliée au point milieu via une résistance de protection $R_B = \SI{22}{k\Omega}$. Le collecteur alimente une LED rouge avec sa résistance de limitation, le courant de saturation souhaité est $I_{C,\text{sat}} = \SI{20}{mA}$. Caractéristiques de la photorésistance : jour $R_\phi = \SI{1}{k\Omega}$, nuit $R_\phi = \SI{100}{k\Omega}$. Déterminer l'état de la LED le jour et la nuit.
\tcblower
\textbf{Étape 1 — État du capteur.}
Le jour (forte lumière), $R_\phi = \SI{1}{k\Omega}$ (faible). La nuit (obscurité), $R_\phi = \SI{100}{k\Omega}$ (grande).

\textbf{Étape 2 — Tension de base $V_{B0}$ (pont diviseur à vide).}
La formule du pont diviseur (capteur en bas) donne :
\[ V_{B0} = E \cdot \frac{R_\phi}{R_1 + R_\phi}. \]
\begin{itemize}
\item \textbf{Le jour} : $V_{B0} = 9 \times \dfrac{1}{100 + 1} \approx \SI{0,089}{V}$.
Comme $V_{B0} < \SI{0,7}{V}$, la jonction base--émetteur est \textbf{bloquée}. $I_B = 0$.
\item \textbf{La nuit} : $V_{B0} = 9 \times \dfrac{100}{100 + 100} = \SI{4,5}{V}$.
Comme $V_{B0} > \SI{0,7}{V}$, la jonction conduit. Le courant de base est limité par $R_B$ :
\[ I_B \approx \frac{V_{B0} - V_{BE}}{R_B} = \frac{4{,}5 - 0{,}7}{\num{22e3}} \approx \SI{173}{\micro A}. \]
\end{itemize}

\textbf{Étape 3 — Condition de saturation.}
Courant de base minimal requis :
\[ I_{B,\min} = \frac{I_{C,\text{sat}}}{\beta} = \frac{\num{20e-3}}{200} = \SI{0,1}{mA} = \SI{100}{\micro A}. \]
La nuit, $I_B \approx \SI{173}{\micro A} \geq \SI{100}{\micro A}$ : condition remplie, \textbf{saturation franche}.

\textbf{Conclusion.}
\begin{itemize}
\item \textbf{Le jour} : Transistor \textbf{bloqué} ($I_B=0$), $V_{CE} = \SI{9}{V}$, LED \textbf{éteinte}.
\item \textbf{La nuit} : Transistor \textbf{saturé} ($V_{CE} \approx \SI{0,2}{V}$), LED \textbf{allumée}.
\end{itemize}
On obtient une \textbf{veilleuse automatique} qui s'allume à la tombée de la nuit. Le choix de $R_1 = \SI{100}{k\Omega}$ (au lieu de \SI{10}{k\Omega}) assure un $V_{B0}$ bien inférieur au seuil le jour, sans avoir besoin de « résistance talon » supplémentaire.
\end{exoresolu}

\begin{exercice}[Indicateur de niveau d'eau par électrodes]
On réalise un indicateur de niveau pour un réservoir d'eau de forage. Deux électrodes sont placées au niveau maximal autorisé. La résistance de l'eau entre les électrodes immergées est estimée à $R_{\text{eau}} = \SI{50}{k\Omega}$. Le circuit de base est alimenté sous $E = \SI{12}{V}$ à travers une résistance $R_1 = \SI{100}{k\Omega}$ en série avec la résistance de l'eau $R_{\text{eau}}$ (le point milieu entre $R_1$ et $R_{\text{eau}}$ est relié à la base via une résistance de protection $R_B = \SI{10}{k\Omega}$). Le transistor a un gain $\beta = 150$. Le buzzer (actionneur) placé au collecteur exige un courant de saturation $I_{C,\text{sat}} = \SI{30}{mA}$ pour sonner correctement.
\textbf{Le buzzer sonne-t-il quand l'eau atteint les électrodes ?} Si non, proposer une modification simple sur $R_1$ pour qu'il sonne.
\end{exercice}

\begin{corrige}[Exercice : Indicateur de niveau d'eau]
\textbf{1. Circuit de base quand l'eau touche les électrodes.}
Le circuit série est : $E \rightarrow R_1 \rightarrow R_{\text{eau}} \rightarrow R_B \rightarrow \text{Base} \rightarrow \text{Émetteur (masse)}$.
La tension aux bornes de la série $R_1 + R_{\text{eau}} + R_B$ est $E - V_{BE} \approx 12 - 0,7 = \SI{11,3}{V}$.
Résistance totale du circuit de base : $R_{\text{tot}} = R_1 + R_{\text{eau}} + R_B = 100 + 50 + 10 = \SI{160}{k\Omega}$.
Courant de base réel :
\[ I_B = \frac{E - V_{BE}}{R_{\text{tot}}} = \frac{11{,}3}{\num{160e3}} \approx \SI{70,6}{\micro A}. \]

\textbf{2. Condition de saturation.}
\[ I_{B,\min} = \frac{I_{C,\text{sat}}}{\beta} = \frac{\num{30e-3}}{150} = \SI{0,2}{mA} = \SI{200}{\micro A}. \]

\textbf{3. Comparaison.}
$I_B \approx \SI{71}{\micro A} < I_{B,\min} = \SI{200}{\micro A}$.
Le transistor n'est \textbf{pas saturé} (il est en régime linéaire faiblement conducteur). Le buzzer ne sonne pas (ou émet un son très faible).

\textbf{4. Correction du montage.}
Il faut augmenter $I_B$ en diminuant la résistance série $R_1$ (car $R_{\text{eau}}$ est imposée par l'eau et $R_B$ protège la base).
Essayons $R_1 = \SI{10}{k\Omega}$ : $R_{\text{tot}} = 10 + 50 + 10 = \SI{70}{k\Omega}$.
$I_B = 11{,}3 / 70\,000 \approx \SI{161}{\micro A}$ (encore $< 200\,\mu A$).
Essayons $R_1 = \SI{4,7}{k\Omega}$ (valeur normalisée) : $R_{\text{tot}} = 4,7 + 50 + 10 = \SI{64,7}{k\Omega}$.
$I_B = 11{,}3 / 64\,700 \approx \SI{175}{\micro A}$ (toujours limite).
Essayons $R_1 = \SI{2,2}{k\Omega}$ : $R_{\text{tot}} = 2,2 + 50 + 10 = \SI{62,2}{k\Omega}$.
$I_B = 11{,}3 / 62\,200 \approx \SI{182}{\micro A}$.
Il faut descendre $R_1$ à environ \SI{1}{k\Omega} ou moins pour atteindre \SI{200}{\micro A} avec ce $R_{\text{eau}}$ et ce $R_B$.
\textbf{Meilleure solution pratique :} Diminuer $R_B$ (si le driver le permet) ou utiliser un transistor à gain $\beta$ plus élevé (Darlington), ou amplifier le signal avec un étage intermédiaire. Ici, avec ce transistor, il faut $R_1 \approx \SI{0}{k\Omega}$ (court-circuit) pour avoir $I_B = 11,3 / 60\,000 \approx \SI{188}{\micro A}$, encore insuffisant. \textbf{Conclusion :} Ce montage simple avec un seul transistor standard est mal adapté à une eau aussi peu conductrice ($50\,k\Omega$) pour piloter \SI{30}{mA}. Il faut soit un transistor Darlington ($\beta \approx 10\,000$), soit un montage à deux transistors, soit accepter un buzzer à plus faible courant.
\end{corrige}

\begin{savaistu}[Le transistor : l'invention qui a changé le monde]
Le transistor a été inventé en 1947 aux laboratoires Bell Labs (USA) par John Bardeen, Walter Brattain et William Shockley (prix Nobel de physique 1956), pour remplacer les tubes à vide, gros, fragiles, chauffants et gourmands en énergie. Miniaturisé puis gravé par milliards sur des puces de silicium (technologie CMOS), il est devenu la brique élémentaire de toute l'électronique numérique et analogique moderne. Un smartphone 2024 contient plusieurs \textbf{dizaines de milliards} de transistors (gravure 3 nm), chacun plus petit qu'un virus (\SI{30}{nm} de grille). Les réseaux MTN et Orange, les décodeurs TNT, les radios des villages électrifiés par les barrages de Nachtigal, Lom Pangar ou Edéa, les systèmes de paiement Mobile Money : tous fonctionnent grâce à des transistors en commutation, exactement comme celui que tu viens d'étudier, mais intégrés par millions dans des circuits intégrés.
\end{savaistu}

\begin{aretenir}[L'essentiel de la section : Transistor en commutation et électrodes]
\begin{itemize}
\item Le transistor NPN possède trois électrodes : base $B$ (commande), collecteur $C$ (entrée puissance), émetteur $E$ (sortie puissance) ; loi des nœuds : $I_E = I_B + I_C \approx I_C$.
\item Jonction base--émetteur = diode : seuil $V_{BE} \approx \SI{0,7}{V}$ (Si). $V_{BE}$ ne dépasse pratiquement pas \SI{0,7}{V}.
\item Deux états en commutation (admis) :
  \begin{itemize}
  \item \textbf{Bloqué} : $I_B = 0 \Rightarrow I_C = 0$, $V_{CE} = E$ (interrupteur ouvert).
  \item \textbf{Saturé} : $I_B \geq I_{C,\text{sat}}/\beta \Rightarrow V_{CE} \approx \SI{0,2}{V}$, $I_C = I_{C,\text{sat}}$ imposé par le circuit extérieur (interrupteur fermé).
  \end{itemize}
\item \textbf{Résistance de base $R_B$ obligatoire} pour limiter $I_B$ et protéger le transistor ($I_B = (U_{\text{cde}} - 0,7)/R_B$).
\item \textbf{Montage type} : Capteur + $R_1$ en pont diviseur sur la base (étage commande) ; Actionneur en série sur le collecteur (étage puissance).
\item \textbf{Électrodes} : capteur tout-ou-rrien de niveau pour liquide \textbf{conducteur} (eau courante). Fermeture du circuit de base par conduction ionique quand le niveau monte.
\item \textbf{Chaîne électronique} = Capteur (info physique $\to$ élec.) $\to$ Traitement (ici transistor en commutation) $\to$ Actionneur (info élec. $\to$ action physique).
\end{itemize}
\end{aretenir}

\coursec{Chaînes électroniques : structure et exemples}

Dans la section précédente, nous avons appris à utiliser le transistor comme un interrupteur commandé : bloqué ou saturé selon le courant injecté dans sa base, grâce à un capteur (photorésistance, thermistance, électrodes). Nous disposons donc désormais de toutes les briques élémentaires. Il reste à les assembler. C'est exactement ce que fait tout appareil électronique autour de toi — la lampe automatique d'une cour à Douala, l'alarme d'une banque, le poste radio qui diffuse CRTV — : il \emph{acquiert} une information, la \emph{traite}, puis \emph{agit}. Cette organisation universelle porte un nom : la \cle{chaîne électronique}.

\courssub{Définition et structure d'une chaîne électronique}

\begin{definition}[Chaîne électronique]
Une \cle{chaîne électronique} est un ensemble organisé de composants qui \textbf{acquiert} une information issue du milieu extérieur (grandeur physique : lumière, température, pression, niveau, onde), la \textbf{traite} (amplification, comparaison à un seuil, commutation) et \textbf{agit} sur le milieu (allumer, sonner, chauffer, déplacer).
\end{definition}

\begin{definition}[Actionneur]
Un \cle{actionneur} est le dispositif de sortie qui convertit le signal électrique traité en une \textbf{action} perceptible ou utile : lumière (lampe, LED), son (sirène, écouteur, haut-parleur), mouvement (moteur, électroaimant), chaleur (résistance chauffante).
\end{definition}

Toute chaîne électronique se décompose en \textbf{trois blocs fonctionnels} :

\begin{itemize}
\item le \cle{capteur} (bloc d'\textbf{acquisition}) : il convertit la grandeur physique d'entrée en signal électrique (tension, courant ou résistance variable) ;
\item le \cle{circuit de traitement} : il exploite ce signal — amplification d'un signal faible, comparaison à une tension de référence, commutation par transistor ou relais ;
\item l'\cle{actionneur} (bloc d'\textbf{action}) : il produit l'effet final.
\end{itemize}

\begin{popfigure}
\begin{center}
\begin{tikzpicture}[font=\small,
bloc/.style={draw=popblue, thick, fill=popblueL, rounded corners=3pt, minimum width=3.1cm, minimum height=1.5cm, align=center},
fleche/.style={-{Stealth[length=3mm]}, thick, popdark}]
\node[bloc, align=center] (cap) at (0,0) {\textbf{Capteur}\\ acquisition};
\node[bloc, align=center] (trt) at (4.6,0) {\textbf{Traitement}\\ amplification,\\ comparaison, commutation};
\node[bloc, align=center] (act) at (9.6,0) {\textbf{Actionneur}\\ action};
\draw[fleche] (cap) -- (trt) node[midway, above, popdark]{signal électrique};
\draw[fleche] (trt) -- (act) node[midway, above, popdark]{signal de commande};
\draw[fleche, poporange] (-3.4,0) -- (cap) node[midway, above, poporange, align=center]{grandeur physique\\ (lumière, $\theta$, niveau, onde)};
\draw[fleche, popgreen] (act) -- (13.0,0) node[midway, above, popgreen, align=center]{effet utile\\ (lumière, son, mouvement)};
\end{tikzpicture}
\end{center}
\legende{Schématisation normalisée d'une chaîne électronique par blocs fonctionnels.}
\end{popfigure}

\begin{methode}[Analyser une chaîne électronique]
\begin{enumerate}
\item \textbf{Identifier les trois blocs} : quel est le capteur ? quel composant traite (transistor, relais, amplificateur) ? quel est l'actionneur ?
\item \textbf{Suivre le signal de l'entrée vers la sortie} : comment la grandeur physique modifie-t-elle la résistance ou la tension délivrée par le capteur ?
\item \textbf{Préciser l'état de l'élément de commutation} : transistor bloqué ($I_B \approx 0$, $I_C \approx 0$) ou saturé ($I_C$ maximal, limité par le circuit extérieur) ; relais au repos ou collé.
\item \textbf{Conclure sur l'état de l'actionneur} dans chaque situation (jour/nuit, niveau bas/haut, etc.).
\end{enumerate}
\end{methode}

\courssub{Exemple 1 : commande d'une sirène par photorésistance et transistor}

Une entreprise de Yaoundé veut qu'une sirène retentisse \textbf{la nuit} en cas d'intrusion détectée par coupure d'un faisceau lumineux. Le montage repose sur une photorésistance (LDR) placée dans un \cle{pont diviseur de tension} qui polarise la base d'un transistor ; le transistor commande un relais qui alimente la sirène.

\begin{popfigure}
\begin{center}
\begin{circuitikz}[european, scale=0.92, transform shape]
\draw (0,4) node[left]{$+E$} to[short,o-] (0,3.6);
\draw (0,3.6) to[R=$R_1$] (0,1.8);
\draw (0,1.8) to[photoresistor, l_={$R_{\text{LDR}}$}] (0,0) node[ground]{};
\draw (0,1.8) -- (1.6,1.8) node[above, popdark]{$V_B$};
\draw (1.6,1.8) to[R=$R_B$] (3.4,1.8);
\draw (3.4,1.8) -- (3.4,1.2);
\node[npn, rotate=-90] (T) at (3.4,0.4) {};
\draw (T.E) node[ground]{};
\draw (T.C) -- (3.4,2.6);
\draw (3.4,2.6) to[L, l_={bobine relais}] (3.4,4.4) to[short,-o] (3.4,4.8) node[above]{$+E$};
\draw[poporange, thick] (5.2,3.4) to[nos, l_={contact relais}] (7.2,3.4);
\draw (5.2,3.4) to[short,-o] (5.2,4.4) node[above]{$+E'$};
\draw (7.2,3.4) -- (7.2,2.2) to[lamp, l_={sirène}] (7.2,0.6) node[ground]{};
\end{circuitikz}
\end{center}
\legende{Chaîne sirène : la LDR et $R_1$ forment un pont diviseur ; la tension $V_B$ commande le transistor, qui commande le relais, qui alimente la sirène.}
\end{popfigure}

\textbf{Analyse du pont diviseur.} Le même courant traverse $R_1$ et $R_{\text{LDR}}$ (le courant de base, très faible, est négligé devant le courant du pont) ; la tension de base est donc donnée par la formule du pont diviseur :
\[
V_B = E \times \frac{R_{\text{LDR}}}{R_1 + R_{\text{LDR}}}.
\]
Vérification dimensionnelle : un rapport de résistances est sans dimension, donc $V_B$ est bien homogène à une tension. On remarque que $V_B$ est une fonction \textbf{croissante} de $R_{\text{LDR}}$ : plus il fait sombre, plus $V_B$ est grande.

\begin{exemplebox}[Chaîne sirène : analyse jour/nuit chiffrée]
Prenons $E = \SI{12}{V}$, $R_1 = \SI{10}{k\Omega}$, et un seuil de saturation du transistor à $V_B \approx \SI{2}{V}$.

\textbf{De jour} : $R_{\text{LDR}} = \SI{1}{k\Omega}$.
\[
V_B = 12 \times \frac{1}{10 + 1} \approx \SI{1,1}{V}.
\]
Cette tension est insuffisante : le courant de base est quasi nul, le transistor est \textbf{bloqué}, $I_C \approx 0$, le relais reste au repos, la sirène est \textbf{éteinte}.

\textbf{De nuit} : $R_{\text{LDR}} = \SI{100}{k\Omega}$.
\[
V_B = 12 \times \frac{100}{10 + 100} \approx \SI{10,9}{V}.
\]
La base est fortement polarisée à travers $R_B$ : le transistor est \textbf{saturé}, $I_C$ devient important, la bobine du relais est alimentée, le contact se ferme et la sirène est \textbf{activée}.
\end{exemplebox}

La chaîne se lit donc : lumière (entrée) $\to$ LDR (capteur) $\to$ pont diviseur + transistor (traitement : comparaison au seuil, commutation) $\to$ relais + sirène (action).

\begin{attention}[Rôle de la résistance $R_B$]
La résistance $R_B$ placée entre le pont diviseur et la base \textbf{limite le courant de base} : sans elle, quand $V_B$ devient grande, le courant $I_B$ dépasserait la valeur maximale admissible et détruirait la jonction base--émetteur du transistor. Ne l'oublie jamais dans un schéma, et ne l'oublie pas dans ta rédaction au Bac.
\end{attention}

\courssub{Exemple 2 : indicateur de niveau d'eau par électrodes et transistor}

L'eau ordinaire (non distillée) conduit faiblement le courant grâce aux ions qu'elle contient : deux électrodes métalliques plongées dans l'eau forment un conducteur de résistance élevée mais finie. On exploite cette propriété pour détecter un niveau.

\begin{popfigure}
\begin{center}
\begin{tikzpicture}[font=\small, scale=0.95]
% Cuve
\draw[thick, popdark] (0,0) -- (0,3.4) (3.4,0) -- (3.4,3.4) (0,0) -- (3.4,0);
\fill[popblue!25] (0.06,0.06) rectangle (3.34,2.2);
\draw[popblue] (0.06,2.2) -- (3.34,2.2) node[midway, above, popblue]{niveau d'eau};
% Électrodes
\draw[thick, poporange] (1.0,3.8) -- (1.0,1.0);
\draw[thick, poporange] (2.4,3.8) -- (2.4,2.6);
\node[poporange, right] at (2.45,3.6) {électrode de détection};
\node[poporange, left] at (0.95,3.6) {électrode commune};
% Circuit : électrode commune -> R_B -> base
\draw (1.0,3.8) -- (1.0,4.6) -- (5.6,4.6) to[R=$R_B$] (7.4,4.6);
\draw (7.4,4.6) -- (7.4,4.0);
\node[npn, rotate=-90] (T) at (7.4,3.2) {};
\draw (T.E) node[ground]{};
\draw (T.C) -- (7.4,5.6);
\draw (7.4,5.6) to[D, l_=LED] (7.4,7.0) to[R=$R$] (7.4,8.2) to[short,-o] (7.4,8.6) node[above]{$+E$};
% Électrode de détection -> +E
\draw (2.4,3.8) -- (2.4,8.6) -- (7.4,8.6);
\end{tikzpicture}
\end{center}
\legende{Indicateur de niveau d'eau : quand l'eau touche l'électrode de détection, le courant de base circule à travers l'eau et la LED s'allume.}
\end{popfigure}

\textbf{Analyse.} L'électrode de détection est reliée au pôle positif $+E$ de l'alimentation ; l'électrode commune, plongée en permanence dans l'eau, est reliée à la base du transistor à travers $R_B$.

\begin{itemize}
\item \textbf{Niveau bas} (eau sous l'électrode de détection) : le circuit de base est ouvert, $I_B = 0$, le transistor est \textbf{bloqué} : la LED est éteinte.
\item \textbf{Niveau haut} (l'eau relie les deux électrodes) : un faible courant de base $I_B$ circule de $+E$ vers la base à travers l'eau ; le transistor amplifie ce courant ($I_C = \beta\, I_B$) et devient passant : la LED \textbf{s'allume}, signalant que le niveau est atteint.
\end{itemize}

C'est le transistor qui rend le dispositif sensible : le courant minuscule traversant l'eau (quelques microampères à quelques dizaines de microampères) suffit, après amplification ($\beta$ de l'ordre de $100$), à alimenter la LED (quelques milliampères).

\begin{savaistu}[Pompes de forage protégées au Cameroun]
De nombreux forages équipés de pompes électriques au Cameroun (régions du Nord, de l'Adamaoua, quartiers périphériques de Yaoundé et Douala) utilisent des \textbf{électrodes de niveau} : quand le niveau d'eau descend sous l'électrode basse, le transistor se bloque et coupe la pompe. Cela protège le moteur contre la \emph{marche à sec}, qui détruirait rapidement la pompe — une application directe, et vitale, de cette leçon !
\end{savaistu}

\courssub{Exemple 3 : chaîne commandée par relais}

Le \cle{relais électromagnétique} joue un rôle précieux : il \textbf{sépare le circuit de commande (basse puissance) du circuit de puissance}. La structure est toujours la même :
\[
\text{capteur} \longrightarrow \text{transistor} \longrightarrow \text{bobine du relais} \longrightarrow \text{contact} \longrightarrow \text{actionneur de puissance}.
\]

Quand le capteur rend le transistor saturé, la bobine est parcourue par un courant, crée un champ magnétique qui attire une palette ferromagnétique : le contact se ferme et alimente l'actionneur (moteur de portail, lampe \SI{220}{V}, compresseur). Ainsi, un courant de quelques milliampères issu d'une thermistance commande un appareil de plusieurs centaines de watts sur le réseau ENEO, sans que le circuit fragile du capteur soit exposé au secteur.

\courssub{Exemple 4 : la chaîne radio — antenne, microphone, écouteur}

La radio est la chaîne électronique la plus complète que tu utilises chaque jour. Elle se décompose ainsi :
\[
\text{antenne (capteur d'ondes)} \longrightarrow \text{circuit d'accord + amplification (traitement)} \longrightarrow \text{écouteur (actionneur)}.
\]

\textbf{Fonctionnement de l'antenne.} Une onde électromagnétique émise par un émetteur (CRTV, stations FM, relais MTN ou Orange) est constituée de champs électrique $\vec{E}$ et magnétique $\vec{B}$ oscillants. En traversant le conducteur de l'antenne, le champ $\vec{E}$ met en mouvement les électrons libres : il apparaît aux bornes de l'antenne une \textbf{tension induite variable}, image fidèle (mais très faible) de l'onde reçue. L'antenne est donc bien un \emph{capteur} : elle convertit une onde en signal électrique.

\begin{attention}[L'antenne ne choisit pas la station !]
L'antenne capte \textbf{toutes les ondes} qui la traversent, de toutes les fréquences. C'est le \textbf{circuit d'accord} (association bobine--condensateur $LC$, réglé par le bouton de syntonisation) qui \emph{sélectionne} la fréquence de la station désirée par résonance. Ne rédige jamais « l'antenne sélectionne la station » : la sélection est une fonction du bloc de \emph{traitement}, pas du capteur.
\end{attention}

\begin{propriete}[Microphone et écouteur : conversions réciproques (admise)]
Dans un \cle{microphone électrodynamique}, le son fait vibrer une membrane solidaire d'une bobine mobile placée dans le champ magnétique d'un aimant : il apparaît aux bornes de la bobine une \textbf{tension induite image du son} (conversion son $\to$ électricité, par induction électromagnétique).

L'\cle{écouteur} réalise la conversion \textbf{inverse} : la tension variable issue de l'amplificateur parcourt la bobine mobile ; la force électromagnétique (force de Laplace) qui s'exerce sur elle fait vibrer la membrane, qui reproduit le son (conversion électricité $\to$ son).
\end{propriete}

\begin{popfigure}
\begin{center}
\begin{tikzpicture}[font=\small, scale=1.0]
% Aimant
\fill[popgold!60] (0,-0.9) rectangle (0.7,0.9);
\node[popdark] at (0.35,0) {\textbf{N}};
\fill[popgold!30] (2.6,-0.9) rectangle (3.3,0.9);
\node[popdark] at (2.95,0) {\textbf{S}};
\draw[-{Stealth}, poppurple, thick] (0.75,0.4) -- (2.55,0.4);
\draw[-{Stealth}, poppurple, thick] (0.75,-0.4) -- (2.55,-0.4);
\node[poppurple, above] at (1.65,0.45) {$\vec{B}$};
% Bobine mobile
\draw[poporange, very thick] (1.45,-0.7) rectangle (1.85,0.7);
\node[poporange, below] at (1.65,-0.75) {bobine mobile};
% Membrane
\draw[popblue, very thick, decorate, decoration={coil, aspect=0.4, segment length=6pt}] (1.85,0) -- (4.6,0);
\node[popblue, right] at (4.6,0) {membrane};
% Fils
\draw (1.45,-0.7) -- (1.45,-1.6) -- (0.6,-1.6) to[short,-o] (0.6,-2.0);
\draw (1.85,-0.7) -- (1.85,-1.6) -- (2.7,-1.6) to[short,-o] (2.7,-2.0);
\node[below] at (1.65,-2.05) {tension variable $u(t)$};
% Flèche vibration
\draw[{Stealth}-{Stealth}, popgreen, thick] (4.9,-0.5) -- (4.9,0.5) node[midway, right, popgreen]{vibration $\to$ son};
\end{tikzpicture}
\end{center}
\legende{Schéma simplifié d'un écouteur électrodynamique : la bobine mobile, plongée dans le champ $\vec{B}$ de l'aimant, est solidaire de la membrane.}
\end{popfigure}

\textbf{Justification physique.} La bobine parcourue par le courant $i(t)$, placée dans le champ $\vec{B}$ de l'aimant, subit une force de Laplace $\vec{F} = i\,\vec{\ell} \wedge \vec{B}$, proportionnelle à $i$. La membrane, solidaire de la bobine, est donc mise en mouvement par une force qui reproduit les variations du signal électrique : elle comprime l'air au rythme du signal, et l'oreille perçoit le son. Dans le microphone, c'est le phénomène d'induction qui joue : le mouvement de la bobine dans le champ crée une f.é.m. induite $e = -\dfrac{d\Phi}{dt}$ proportionnelle à la vitesse de la membrane, donc image du son reçu.

\begin{exoresolu}[Identifier les blocs d'une chaîne]
Énoncé : un hall d'immeuble à Bonanjo s'éclaire automatiquement à la tombée de la nuit grâce à une photorésistance, un transistor, un relais et une lampe alimentée en \SI{220}{V}. Identifier les trois blocs fonctionnels et décrire l'état de chaque élément le jour et la nuit.
\tcblower
Solution détaillée :
\begin{itemize}
\item \textbf{Capteur} : la photorésistance (lumière $\to$ résistance variable).
\item \textbf{Traitement} : le pont diviseur et le transistor (comparaison au seuil, commutation), puis la bobine du relais.
\item \textbf{Actionneur} : la lampe (électricité $\to$ lumière).
\end{itemize}
\textbf{Le jour} : $R_{\text{LDR}}$ faible ($\approx \SI{1}{k\Omega}$), la tension de base est faible, le transistor est bloqué, la bobine du relais n'est pas alimentée, le contact reste ouvert : la lampe est éteinte.
\textbf{La nuit} : $R_{\text{LDR}}$ grande ($\approx \SI{100}{k\Omega}$), la tension de base dépasse le seuil, le transistor sature, le relais colle, le contact ferme le circuit \SI{220}{V} : la lampe s'allume. Le relais assure ici l'isolement entre la commande basse tension et la puissance du secteur.
\end{exoresolu}

\begin{aretenir}[L'essentiel]
\begin{itemize}
\item Une chaîne électronique = \textbf{capteur} (acquisition) $\to$ \textbf{traitement} (amplification, comparaison, commutation) $\to$ \textbf{actionneur} (action).
\item L'actionneur convertit le signal électrique en lumière, son, mouvement ou chaleur.
\item Le transistor amplifie et commute ; le relais isole la commande de la puissance.
\item L'antenne convertit l'onde en tension ; le microphone convertit le son en tension (induction) ; l'écouteur convertit la tension en son (force de Laplace).
\item C'est le circuit d'accord, et non l'antenne, qui sélectionne la station de radio.
\end{itemize}
\end{aretenir}

\newpage

\coursec{Activité d'intégration}

\courssub{Objectifs de l'activité}

À l'issue de cette activité, tu seras capable de :
\begin{itemize}
\item exploiter les caractéristiques expérimentales de capteurs réels : courbe $R = f(\theta)$ d'une thermistance CTN et tableau $R = f(\Phi)$ d'une photorésistance ;
\item distinguer un capteur \cle{linéaire} d'un capteur \cle{non linéaire} et justifier ta réponse ;
\item dimensionner un pont diviseur de tension et une résistance de base pour faire fonctionner un transistor en \cle{commutation} (bloqué / saturé) à un seuil choisi ;
\item schématiser une \cle{chaîne électronique} complète (capteur $\rightarrow$ traitement $\rightarrow$ actionneur) en blocs fonctionnels et en circuit électrique ;
\item analyser de manière critique un système technique complet et proposer des améliorations.
\end{itemize}

\courssub{Situation}

Monsieur Kamdem est maraîcher à Bafoussam, dans la région de l'Ouest. Il cultive des tomates sous une serre de \SI{80}{m^2}. La chaleur de la mi-journée, les nuits fraîches et les coupures d'eau fréquentes menacent sa production. Il te demande, à toi élève de Terminale C, de concevoir un système électronique simple et peu coûteux, alimenté par une batterie de $E = \SI{12}{V}$ rechargée par un petit panneau solaire, qui réalise trois fonctions :

\begin{enumerate}
\item \textbf{Ventilation} : un ventilateur doit se déclencher automatiquement dès que la température intérieure dépasse $\theta_s = \SI{35}{\celsius}$ ;
\item \textbf{Éclairage} : une lampe doit s'allumer automatiquement à la tombée de la nuit ;
\item \textbf{Alarme de niveau} : un buzzer doit retentir lorsque le niveau d'eau du bac de réserve descend sous deux électrodes de contrôle.
\end{enumerate}

\textbf{Données techniques.}

\textbf{Capteur de température : thermistance CTN.} Le constructeur fournit les mesures suivantes :

\begin{center}
\begin{tabular}{|l|c|c|c|c|c|c|c|}
\hline
\rowcolor{popblueL} $\theta$ (\si{\celsius}) & 10 & 20 & 25 & 30 & 35 & 40 & 50 \\
\hline $R_{\text{CTN}}$ (\si{k\ohm}) & 20,0 & 12,5 & 10,0 & 8,0 & 6,5 & 5,3 & 3,6 \\
\hline
\end{tabular}
\end{center}

\textbf{Capteur de lumière : photorésistance (LDR).}

\begin{center}
\begin{tabular}{|l|c|c|c|c|c|}
\hline
\rowcolor{popblueL} Éclairement $\Phi$ (lux) & 10 (nuit) & 100 & 500 & 1000 & 5000 (plein jour) \\
\hline $R_{\text{LDR}}$ (\si{k\ohm}) & 500 & 50 & 12 & 6 & 1,5 \\
\hline
\end{tabular}
\end{center}

\textbf{Transistors utilisés (NPN, montage émetteur commun)} : gain $\beta = 200$ ; tension base-émetteur en conduction $V_{BE} = \SI{0,7}{V}$ ; tension de saturation $V_{CE,\text{sat}} \approx \SI{0}{V}$. Pour garantir une saturation franche, on impose un courant de base $I_B = 2\,I_{B,\min}$ (coefficient de sursaturation $k = 2$).

\textbf{Actionneurs} (branchés entre le $+\SI{12}{V}$ et le collecteur) :
\begin{itemize}
\item ventilateur : $I_{C,\text{sat}} = \SI{240}{mA}$ ;
\item lampe : $I_{C,\text{sat}} = \SI{120}{mA}$ ;
\item buzzer : $I_{C,\text{sat}} = \SI{30}{mA}$.
\end{itemize}

\textbf{Seuils choisis.} Ventilation : basculement à $\theta_s = \SI{35}{\celsius}$. Éclairage : basculement à $\Phi_s = 100$ lux. Alarme : l'eau du bac, légèrement conductrice, présente une résistance d'environ $R_{\text{eau}} = \SI{10}{k\ohm}$ entre les deux électrodes immergées ; le circuit est ouvert ($R \to \infty$) quand le niveau descend sous les électrodes.

\textbf{Principe retenu pour chaque voie} : le capteur est associé à un résistor fixe en \cle{pont diviseur de tension} ; la tension du point milieu attaque la base du transistor à travers une résistance de base $R_B$.

\courssub{Consignes}

\textbf{Partie A — Exploitation des caractéristiques des capteurs}

\begin{enumerate}
\item Tracer sur papier millimétré la courbe $R_{\text{CTN}} = f(\theta)$. La thermistance est-elle un capteur linéaire ? Justifier. Quel est le signe de son coefficient de température et que signifie « CTN » ?
\item À partir du tableau de la LDR, calculer les rapports $\dfrac{R_{\text{LDR}}}{\Phi}$ pour quelques points. La photorésistance est-elle un capteur linéaire en $\Phi$ ? Proposer une grandeur transformée (par exemple $\log R$ en fonction de $\log \Phi$) qui rendrait la caractéristique presque linéaire.
\item Relever les valeurs des capteurs aux seuils : $R_{\text{CTN}}(\SI{35}{\celsius})$ et $R_{\text{LDR}}(100\ \text{lux})$.
\end{enumerate}

\textbf{Partie B — Dimensionnement des trois voies}

\begin{enumerate}
\setcounter{enumi}{3}
\item \textbf{Voie ventilation.} On monte la CTN en haut du pont diviseur (entre $+\SI{12}{V}$ et la base) et un résistor $R_1$ en bas. Montrer que la tension de base au seuil vaut $V_{B,s} = E\,\dfrac{R_1}{R_1 + R_{\text{CTN},s}}$. En négligeant le courant de base devant le courant du pont, choisir $V_{B,s} = \SI{1,4}{V}$ (marge au-dessus de $V_{BE}$) et calculer $R_1$.
\item Calculer le courant de base minimal $I_{B,\min} = \dfrac{I_{C,\text{sat}}}{\beta}$ du ventilateur, puis la résistance de base $R_B = \dfrac{V_{B,s} - V_{BE}}{k\,I_{B,\min}}$ avec $k = 2$. Vérifier que l'hypothèse de la question 4 est légitime.
\item \textbf{Voie éclairage.} On souhaite que la lampe s'allume quand il fait nuit. Faut-il placer la LDR en haut ou en bas du pont ? Justifier par le sens de variation de $R_{\text{LDR}}$. Avec le même seuil $V_{B,s} = \SI{1,4}{V}$, calculer le résistor associé puis $R_B$ pour la lampe.
\item \textbf{Voie alarme.} On adopte le montage suivant : les deux électrodes (résistance $R_{\text{eau}} = \SI{10}{k\ohm}$ quand elles sont immergées, circuit ouvert sinon) sont placées entre $+\SI{12}{V}$ et la base ; un résistor $R_3 = \SI{47}{k\ohm}$ relie la base à la masse ; le buzzer est branché entre $+\SI{12}{V}$ et le collecteur. Calculer $V_B$ bac plein et en déduire l'état du transistor puis la tension aux bornes du buzzer. Que devient $V_B$ quand le niveau est bas ? Conclure sur le fonctionnement de l'alarme et vérifier que le transistor est bien saturé bac plein.
\end{enumerate}

\textbf{Partie C — Schématisation et analyse}

\begin{enumerate}
\setcounter{enumi}{7}
\item Représenter les trois chaînes électroniques en \textbf{blocs fonctionnels} (capteur $\rightarrow$ traitement $\rightarrow$ actionneur), en précisant la grandeur physique d'entrée et la grandeur de sortie de chaque bloc.
\item Dessiner le schéma électrique complet de la voie ventilation.
\item Monsieur Kamdem craint que le ventilateur ne s'allume et ne s'éteigne sans cesse autour de \SI{35}{\celsius}. Proposer une explication du phénomène et une piste d'amélioration (on pourra citer le rôle d'un relais ou d'un montage à hystérésis).
\end{enumerate}

\courssub{Ressources à mobiliser}

\begin{itemize}
\item \textbf{Pont diviseur de tension} (courant de base négligé) : $V_B = E\,\dfrac{R_{\text{bas}}}{R_{\text{haut}} + R_{\text{bas}}}$.
\item \textbf{Transistor en commutation} : bloqué si $V_{BE} < \SI{0,7}{V}$ ; saturé si $I_B \geq I_{B,\min} = \dfrac{I_{C,\text{sat}}}{\beta}$ ; résistance de base $R_B = \dfrac{V_B - V_{BE}}{k\,I_{B,\min}}$.
\item \textbf{Capteur linéaire} : la caractéristique (grandeur électrique en fonction de la grandeur physique) est une droite ; sinon le capteur est non linéaire.
\item \textbf{Chaîne électronique} : capteur (conversion d'une grandeur physique en signal électrique), unité de traitement (comparaison à un seuil, amplification, commutation), actionneur (conversion du signal électrique en action).
\item \textbf{Électrodes} : l'eau non distillée conduit faiblement le courant ; deux électrodes immergées ferment le circuit, à découvert elles l'ouvrent.
\end{itemize}

\courssub{Démarche et corrigé commenté}

\begin{exoresolu}[Corrigé de l'activité d'intégration]
Énoncé : concevoir les trois voies (ventilation à \SI{35}{\celsius}, éclairage nocturne, alarme de niveau bas) de la serre de Monsieur Kamdem avec $E = \SI{12}{V}$, $\beta = 200$, $k = 2$, $V_{B,s} = \SI{1,4}{V}$.
\tcblower
\textbf{Partie A — Caractéristiques des capteurs.}

\textbf{1.} Le tracé de $R_{\text{CTN}} = f(\theta)$ donne une courbe décroissante, concave vers le haut : ce n'est \textbf{pas une droite}, la thermistance est un capteur \cle{non linéaire}. La résistance diminue quand la température augmente : le coefficient de température est \textbf{négatif} (CTN = coefficient de température négatif ; en anglais NTC).

\textbf{2.} Pour la LDR : $\dfrac{500}{10} = 50$ ; $\dfrac{50}{100} = 0,5$ ; $\dfrac{1,5}{5000} = \num{3e-4}$ (en \si{k\ohm}/lux). Le rapport n'est pas constant : la LDR est \textbf{non linéaire} en $\Phi$. En traçant $\log R$ en fonction de $\log \Phi$, les points s'alignent presque : la caractéristique suit une loi de puissance $R \propto \Phi^{-\alpha}$, linéarisable en échelle logarithmique.

\textbf{3.} Lecture aux seuils : $R_{\text{CTN},s} = \SI{6,5}{k\ohm}$ à \SI{35}{\celsius} ; $R_{\text{LDR},s} = \SI{50}{k\ohm}$ à 100 lux.

\textbf{Partie B — Dimensionnement.}

\textbf{4.} Le pont étant parcouru par le même courant (courant de base négligé), la loi d'additivité des tensions donne :
\[ V_{B,s} = E\,\dfrac{R_1}{R_1 + R_{\text{CTN},s}} \quad\Longrightarrow\quad R_1 = R_{\text{CTN},s}\,\dfrac{V_{B,s}}{E - V_{B,s}} = 6,5\times\dfrac{1,4}{12 - 1,4} = \SI{0,86}{k\ohm}. \]
On choisit la valeur normalisée $R_1 = \SI{820}{\ohm}$ (ajustable avec un potentiomètre pour régler finement le seuil). Quand $\theta$ dépasse \SI{35}{\celsius}, $R_{\text{CTN}}$ diminue, donc $V_B$ augmente au-dessus du seuil : le transistor se sature et le ventilateur tourne. C'est bien le comportement voulu.

\textbf{5.} Ventilateur : $I_{B,\min} = \dfrac{\SI{240}{mA}}{200} = \SI{1,2}{mA}$. Avec $k = 2$ : $I_B = \SI{2,4}{mA}$, d'où
\[ R_B = \dfrac{V_{B,s} - V_{BE}}{k\,I_{B,\min}} = \dfrac{1,4 - 0,7}{\num{2,4e-3}} \approx \SI{292}{\ohm} \;\Rightarrow\; R_B = \SI{270}{\ohm}\ \text{(normalisée)}. \]
Vérification de l'hypothèse : le courant du pont vaut $\dfrac{E}{R_{\text{CTN},s} + R_1} = \dfrac{12}{6500 + 820} \approx \SI{1,6}{mA}$, du même ordre que $I_B$ : l'approximation est grossière, d'où l'intérêt du potentiomètre d'ajustage et de la sursaturation $k = 2$ qui garantit le basculement malgré l'imprécision. (Pour un dimensionnement fin, on refait le calcul avec le courant de base retranché.)

\textbf{6.} La nuit, $R_{\text{LDR}}$ \emph{augmente}. Pour que $V_B$ augmente la nuit, il faut placer la \textbf{LDR en bas} du pont (entre base et masse) et le résistor fixe $R_2$ en haut : alors $V_B = E\,\dfrac{R_{\text{LDR}}}{R_2 + R_{\text{LDR}}}$ croît avec $R_{\text{LDR}}$. Au seuil :
\[ R_2 = R_{\text{LDR},s}\,\dfrac{E - V_{B,s}}{V_{B,s}} = 50\times\dfrac{10,6}{1,4} \approx \SI{379}{k\ohm} \;\Rightarrow\; R_2 = \SI{390}{k\ohm}. \]
Lampe : $I_{B,\min} = \dfrac{120}{200} = \SI{0,6}{mA}$ ; $R_B = \dfrac{0,7}{\num{1,2e-3}} \approx \SI{580}{\ohm} \Rightarrow \SI{560}{\ohm}$.

\textbf{7.} Voie alarme. Montage : électrodes entre $+\SI{12}{V}$ et la base, $R_3 = \SI{47}{k\ohm}$ entre base et masse, buzzer entre $+\SI{12}{V}$ et le collecteur.

\emph{Bac plein} : les électrodes immergées équivalent à $R_{\text{eau}} = \SI{10}{k\ohm}$ ; le pont $\{R_{\text{eau}}, R_3\}$ donne
\[ V_B = E\,\dfrac{R_3}{R_{\text{eau}} + R_3} = 12\times\dfrac{47}{10 + 47} \approx \SI{9,9}{V} \gg V_{BE}. \]
Le transistor conduit. Vérifions la saturation : le courant délivré à la base vaut, par la loi des nœuds (la base étant portée à $V_{BE} = \SI{0,7}{V}$),
\[ I_B = \dfrac{E - V_{BE}}{R_{\text{eau}}} - \dfrac{V_{BE}}{R_3} = \dfrac{12 - 0,7}{\num{10e3}} - \dfrac{0,7}{\num{47e3}} \approx \SI{1,13}{mA} - \SI{0,015}{mA} \approx \SI{1,1}{mA}, \]
très supérieur à $I_{B,\min} = \dfrac{\SI{30}{mA}}{200} = \SI{0,15}{mA}$ : le transistor est \textbf{saturé}, $V_{CE} \approx \SI{0}{V}$. Le buzzer, branché entre $+\SI{12}{V}$ et le collecteur, voit une tension quasi nulle à ses bornes : il est \textbf{silencieux}. C'est bien ce qu'on veut quand la réserve est pleine.

\emph{Niveau bas} : les électrodes sont à découvert, le circuit de base est ouvert ; $R_3$ ramène la base à la masse : $V_B = \SI{0}{V} < V_{BE}$, le transistor est \textbf{bloqué}, aucun courant de collecteur ne circule. Les \SI{12}{V} de la batterie s'appliquent alors intégralement aux bornes du buzzer : il \textbf{sonne}. L'alarme fonctionne comme demandé, sans résistance de base supplémentaire : la limitation du courant de base est assurée par $R_{\text{eau}}$ et $R_3$.

\begin{attention}[Piège classique de la voie alarme]
On pourrait croire qu'il faut « saturer le transistor pour faire sonner le buzzer ». Ici c'est l'inverse : le buzzer est alimenté quand le transistor est \textbf{bloqué} (le collecteur est alors au potentiel $+E$ à travers le buzzer). Avant tout calcul, toujours se demander : dans quel état du transistor l'actionneur est-il alimenté ?
\end{attention}

\textbf{Partie C — Schématisation.}

\textbf{8.} Blocs fonctionnels (identiques en structure pour les trois voies) :
\begin{center}
\begin{tikzpicture}[node distance=0.35cm, every node/.style={draw, rounded corners, minimum height=0.9cm, align=center, font=\small}]
\node[fill=popblueL, align=center] (c) {Capteur\\ CTN / LDR / électrodes};
\node[fill=poporangeL, right=of c, align=center] (t) {Traitement\\ pont diviseur + transistor};
\node[fill=popgreenL, right=of t, align=center] (a) {Actionneur\\ ventilateur / lampe / buzzer};
\draw[-{Stealth}, thick] (c) -- (t) node[midway, above, draw=none, font=\scriptsize]{signal électrique};
\draw[-{Stealth}, thick] (t) -- (a) node[midway, above, draw=none, font=\scriptsize]{courant de commande};
\node[draw=none, left=0.15cm of c, font=\scriptsize] {$\theta$, $\Phi$, niveau};
\end{tikzpicture}
\end{center}

\textbf{9.} Schéma de la voie ventilation :
\begin{center}
\begin{circuitikz}[european, scale=0.9, transform shape]
\draw (0,0) node[below]{masse} -- (0,0.5);
\draw (0,4) to[short, -o] (-1.5,4) node[left]{$+\SI{12}{V}$};
\draw (0,4) to[R, l={$R_{\text{CTN}}$}] (0,2.2) to[R, l={$R_1$}] (0,0);
\draw (0,2.2) to[R, l_={$R_B$}] (2.2,2.2) -- (2.2,1.55);
\draw (2.2,1.5) node[npn, anchor=B](Q){};
\node[right] at (Q) {T};
\draw (Q.E) -- (2.2,0) -- (0,0);
\draw (Q.C) -- (2.2,3.2) to[generic, l=ventilateur] (2.2,4) -- (0,4);
\end{circuitikz}
\end{center}

\textbf{10.} Sans hystérésis, quand $\theta$ frôle \SI{35}{\celsius}, le moindre bruit fait osciller $V_B$ autour du seuil : le ventilateur « pompe ». Solutions : intercaler un \textbf{relais électromagnétique} (ses contacts créent une commutation franche et isolent la commande de la puissance) ou un montage comparateur à \textbf{hystérésis} (deux seuils distincts, par exemple marche à \SI{36}{\celsius}, arrêt à \SI{33}{\celsius}).
\end{exoresolu}

\courssub{Bilan de l'activité}

Cette activité t'a permis de mettre en évidence plusieurs idées essentielles de la leçon :

\begin{itemize}
\item Les capteurs réels (CTN, LDR) sont le plus souvent \textbf{non linéaires} : l'exploitation de leurs caractéristiques impose la lecture graphique ou tabulaire, et parfois un changement de variable pour linéariser.
\item Le \textbf{pont diviseur de tension} est l'interface universelle entre un capteur résistif et l'étage de traitement : le choix de la position du capteur (en haut ou en bas) détermine le sens de la commande.
\item Le \textbf{transistor en commutation} se dimensionne par deux calculs simples : $I_{B,\min} = I_{C,\text{sat}}/\beta$ puis $R_B$, avec un coefficient de sursaturation pour garantir un basculement franc.
\item Selon le branchement de l'actionneur, c'est l'état \textbf{saturé} ou l'état \textbf{bloqué} du transistor qui l'alimente : l'analyse des deux états (bac plein / bac vide) est indispensable avant de conclure.
\item Toute chaîne électronique se ramène à la même architecture \textbf{capteur $\rightarrow$ traitement $\rightarrow$ actionneur}, qu'il s'agisse d'une serre de Bafoussam, d'un indicateur de niveau d'eau ou d'un poste radio.
\item Un système technique réel exige une \textbf{analyse critique} : stabilité du seuil, hystérésis, tolérances des composants, alimentation autonome (panneau solaire et batterie) — autant de compétences transposables à tout projet d'électronique appliquée.
\end{itemize}

\newpage

\coursec{Méthodes et exercices résolus}

\courssub{Méthode 1 : Interpréter la caractéristique $U = f(I, R)$ d'un rhéostat}

\begin{methode}[Exploiter un rhéostat en montage potentiométrique ou en série]
\begin{enumerate}
\item \textbf{Identifier le montage.} En série (2 bornes utilisées), le rhéostat limite le courant ; en potentiomètre (3 bornes), il fournit une tension réglable entre $0$ et $E$.
\item \textbf{Écrire la loi d'Ohm et la loi des mailles.} En série avec un récepteur $R_0$ sous un générateur $E$ : $I = \dfrac{E}{R_0 + R_h}$ et $U_{R_h} = R_h I$.
\item \textbf{Analyser les positions extrêmes du curseur.} Curseur en bout de course : $R_h = 0$ (courant maximal) ou $R_h = R_{h,\max}$ (courant minimal).
\item \textbf{Tracer ou lire $U = f(I)$.} Pour le rhéostat seul : $U = R_h I$ (droite de pente $R_h$ si $R_h$ est fixé). Pour l'association générateur--récepteur : $U = E - R_0 I$, droite décroissante dont l'intersection avec la caractéristique du récepteur donne le point de fonctionnement.
\item \textbf{Vérifier la puissance admissible} : $P = R_h I^2 \leq P_{\max}$, condition souvent oubliée et pourtant essentielle pour ne pas détériorer le rhéostat.
\end{enumerate}
\end{methode}

\begin{exoresolu}[Rhéostat de limitation dans un circuit de laboratoire]
Un rhéostat de résistance maximale $R_h = \SI{100}{\ohm}$ et de courant admissible $I_{\max} = \SI{2}{A}$ est monté en série avec une lampe de résistance $R_L = \SI{20}{\ohm}$ (supposée constante) et un générateur de f.é.m. $E = \SI{24}{V}$ de résistance interne négligeable.
\begin{enumerate}
\item Quelle valeur minimale $R_{h,\min}$ doit-on donner à la résistance du rhéostat pour ne pas dépasser $I_{\max}$ ?
\item Pour $R_h = \SI{40}{\ohm}$, calculer $I$, la tension $U_L$ aux bornes de la lampe et la puissance dissipée dans le rhéostat.
\item Tracer l'allure de $U_L = f(R_h)$ et conclure sur le rôle du rhéostat.
\end{enumerate}
\tcblower
\textbf{1. Valeur minimale du rhéostat.}

La loi des mailles s'écrit : $E = (R_L + R_h) I$, d'où $I = \dfrac{E}{R_L + R_h}$.

La condition $I \leq I_{\max}$ impose :
\[
R_L + R_h \geq \frac{E}{I_{\max}} \quad\Longrightarrow\quad R_h \geq \frac{E}{I_{\max}} - R_L = \frac{24}{2} - 20 = 12 - 20 = \SI{-8}{\ohm}.
\]
Le résultat étant négatif, la contrainte est automatiquement satisfaite : même à $R_h = 0$, le courant vaut $I_0 = \dfrac{24}{20} = \SI{1,2}{A} < \SI{2}{A}$. Donc $R_{h,\min} = \SI{0}{\ohm}$ : le rhéostat peut être utilisé sur toute sa plage sans danger.

\textbf{2. Point de fonctionnement pour $R_h = \SI{40}{\ohm}$.}
\begin{align*}
I &= \frac{E}{R_L + R_h} = \frac{24}{20 + 40} = \frac{24}{60} = \SI{0,40}{A},\\
U_L &= R_L I = 20 \times 0{,}40 = \SI{8,0}{V},\\
P_h &= R_h I^2 = 40 \times (0{,}40)^2 = 40 \times 0{,}16 = \SI{6,4}{W}.
\end{align*}
Vérification par la loi des mailles : $U_h = R_h I = \SI{16}{V}$ et $U_L + U_h = 8 + 16 = \SI{24}{V} = E$. \checkmark

\textbf{3. Allure de $U_L = f(R_h)$.}

$U_L = \dfrac{R_L E}{R_L + R_h} = \dfrac{480}{20 + R_h}$ : fonction décroissante de $R_h$, de $U_L = \SI{24}{V}$ (pour $R_h = 0$) à $U_L = \dfrac{480}{120} = \SI{4,0}{V}$ (pour $R_h = \SI{100}{\ohm}$). La courbe est une branche d'hyperbole : la tension aux bornes de la lampe n'est jamais nulle.

\textbf{Conclusion.} En montage série, le rhéostat permet de faire varier la tension de la lampe entre $\SI{4}{V}$ et $\SI{24}{V}$ seulement : c'est un \cle{limiteur de courant}. Pour obtenir une variation de $0$ à $E$, il faut le monter en \cle{potentiomètre} (3 bornes).
\end{exoresolu}

\courssub{Méthode 2 : Interpréter la caractéristique $R = f(\theta)$ d'une thermistance}

\begin{methode}[Exploiter une thermistance CTN]
\begin{enumerate}
\item \textbf{Reconnaître le type.} CTN (coefficient de température négatif) : $R$ \emph{décroît} quand $\theta$ croît ; la caractéristique est fortement \emph{non linéaire} (allure exponentielle décroissante). La KTY (au silicium) est quasi linéaire et croissante.
\item \textbf{Lire graphiquement} la résistance à la température demandée, ou inversement la température correspondant à une résistance mesurée (c'est le principe du thermomètre à CTN).
\item \textbf{Monter un pont diviseur} pour convertir la variation de résistance en variation de tension : avec $R$ fixe et $R_{CTN}$ en série sous $E$, $U = \dfrac{R}{R + R_{CTN}} E$.
\item \textbf{Exploiter la sensibilité} : $s = \left|\dfrac{\Delta R}{\Delta \theta}\right|$ autour du point de fonctionnement ; plus la courbe est pentue, plus le capteur est sensible.
\item \textbf{Conclure sur la linéarité} : un capteur est linéaire si sa caractéristique est une droite ; la CTN ne l'est pas, la KTY l'est approximativement.
\end{enumerate}
\end{methode}

\begin{exoresolu}[Thermomètre à thermistance CTN]
Une thermistance CTN obéit, entre $\SI{20}{\celsius}$ et $\SI{80}{\celsius}$, à la loi approchée $R(\theta) = R_0\, e^{-b(\theta - \theta_0)}$ avec $R_0 = \SI{1000}{\ohm}$ à $\theta_0 = \SI{25}{\celsius}$ et $b = \SI{0,045}{\celsius^{-1}}$.
\begin{enumerate}
\item Calculer $R$ à $\theta_1 = \SI{25}{\celsius}$, $\theta_2 = \SI{45}{\celsius}$ et $\theta_3 = \SI{65}{\celsius}$.
\item La CTN est montée en série avec un résistor $R_1 = \SI{500}{\ohm}$ sous $E = \SI{6,0}{V}$. Calculer la tension $U$ aux bornes de $R_1$ à $\theta_2 = \SI{45}{\celsius}$.
\item La tension mesurée vaut $U' = \SI{3,9}{V}$. Déterminer la résistance de la CTN puis la température correspondante.
\end{enumerate}
\tcblower
\textbf{1. Résistances aux trois températures.}
\begin{align*}
R(25) &= 1000 \times e^{0} = \SI{1000}{\ohm},\\
R(45) &= 1000 \times e^{-0{,}045 \times 20} = 1000 \times e^{-0{,}90} \approx 1000 \times 0{,}407 \approx \SI{407}{\ohm},\\
R(65) &= 1000 \times e^{-0{,}045 \times 40} = 1000 \times e^{-1{,}80} \approx 1000 \times 0{,}165 \approx \SI{165}{\ohm}.
\end{align*}
On constate que $R$ décroît fortement quand $\theta$ croît : c'est bien une \cle{CTN}, et la décroissance n'est pas régulière ($-593\,\Omega$ puis $-242\,\Omega$ par pas de $\SI{20}{\celsius}$) : le capteur est \cle{non linéaire}.

\textbf{2. Tension aux bornes de $R_1$ à $\SI{45}{\celsius}$.}

Pont diviseur de tension :
\[
U = \frac{R_1}{R_1 + R_{CTN}} E = \frac{500}{500 + 407} \times 6{,}0 = \frac{500}{907} \times 6{,}0 \approx \SI{3,3}{V}.
\]

\textbf{3. Remontée à la température.}

De $U' = \dfrac{R_1 E}{R_1 + R_{CTN}}$, on tire :
\[
R_{CTN} = R_1\left(\frac{E}{U'} - 1\right) = 500 \times \left(\frac{6{,}0}{3{,}9} - 1\right) = 500 \times 0{,}538 \approx \SI{269}{\ohm}.
\]
On inverse la loi $R = R_0 e^{-b(\theta - \theta_0)}$ :
\[
\theta = \theta_0 - \frac{1}{b}\ln\!\left(\frac{R}{R_0}\right) = 25 - \frac{1}{0{,}045}\ln(0{,}269) = 25 + \frac{1{,}313}{0{,}045} \approx 25 + 29{,}2 \approx \SI{54}{\celsius}.
\]
\textbf{Conclusion.} La mesure d'une simple tension ($\SI{3,9}{V}$) permet, après étalonnage, de remonter à la température ($\approx \SI{54}{\celsius}$) : c'est le principe des thermomètres électroniques utilisés dans les hôpitaux.
\end{exoresolu}

\courssub{Méthode 3 : Interpréter la caractéristique d'une photorésistance (LDR)}

\begin{methode}[Exploiter une photorésistance]
\begin{enumerate}
\item \textbf{Retenir le comportement qualitatif} : la résistance d'une LDR \emph{décroît} quand l'éclairement $\Phi$ augmente (de quelques M$\Omega$ dans l'obscurité à quelques centaines d'$\Omega$ en pleine lumière). La caractéristique est non linéaire.
\item \textbf{Choisir la position dans le pont diviseur} selon l'effet voulu :
\begin{itemize}
\item LDR en haut, $R$ en bas : $U_R$ \emph{augmente} avec la lumière ;
\item LDR en bas : $U_{LDR}$ \emph{diminue} avec la lumière.
\end{itemize}
\item \textbf{Appliquer le pont diviseur} : $U = \dfrac{R}{R + R_{LDR}} E$, puis comparer la tension obtenue au seuil de basculement de l'étage de commande (transistor).
\item \textbf{Prévoir les deux états} (obscurité / lumière) et conclure sur l'état de l'actionneur (lampe, sirène, relais).
\end{enumerate}
\end{methode}

\begin{exoresolu}[Éclairage automatique d'une cour à Yaoundé]
Pour allumer automatiquement une lampe à la tombée de la nuit, on utilise une LDR dont la résistance vaut $R_J = \SI{1,0}{k\ohm}$ en plein jour et $R_N = \SI{100}{k\ohm}$ dans l'obscurité. Elle est placée en série avec un résistor $R = \SI{10}{k\ohm}$ sous $E = \SI{12}{V}$, la LDR étant reliée au pôle positif. La tension $U_R$ aux bornes de $R$ attaque la base d'un transistor qui sature dès que $U_R \geq \SI{1,4}{V}$ (la lampe, placée dans le collecteur, s'allume alors).
\begin{enumerate}
\item Calculer $U_R$ en plein jour puis dans l'obscurité.
\item Conclure sur le fonctionnement du dispositif.
\end{enumerate}
\tcblower
\textbf{1. Calcul des tensions.}

Pont diviseur : $U_R = \dfrac{R}{R + R_{LDR}} E$.

En plein jour :
\[
U_R = \frac{10}{10 + 1{,}0} \times 12 = \frac{120}{11} \approx \SI{10,9}{V}.
\]
Dans l'obscurité :
\[
U_R = \frac{10}{10 + 100} \times 12 = \frac{120}{110} \approx \SI{1,1}{V}.
\]

\textbf{2. Fonctionnement.}

En plein jour, $U_R \approx \SI{10,9}{V} > \SI{1,4}{V}$ : le transistor est \emph{saturé} et la lampe serait allumée... ce qui est l'inverse de l'effet souhaité ! Il faut donc \emph{inverser les positions} de la LDR et de $R$ dans le pont (ou insérer un étage inverseur). Avec la LDR en bas :
\[
\text{Jour : } U_{LDR} = \frac{1{,}0}{11} \times 12 \approx \SI{1,1}{V} < \SI{1,4}{V} \Rightarrow \text{transistor bloqué, lampe éteinte} ;
\]
\[
\text{Nuit : } U_{LDR} = \frac{100}{110} \times 12 \approx \SI{10,9}{V} > \SI{1,4}{V} \Rightarrow \text{transistor saturé, lampe allumée}.
\]
\textbf{Conclusion.} Le dispositif n'allume la lampe qu'à la tombée de la nuit, à condition de placer la LDR \emph{en bas} du pont diviseur. La position du capteur dans le pont détermine la logique (directe ou inversée) de la commande : c'est un point fréquemment exigé au Bac.
\end{exoresolu}

\courssub{Méthode 4 : Distinguer un capteur linéaire d'un capteur non linéaire}

\begin{methode}[Test de linéarité d'un capteur]
\begin{enumerate}
\item \textbf{Critère graphique} : un capteur est linéaire si sa caractéristique (grandeur électrique en fonction de la grandeur physique mesurée) est une \emph{droite} sur la plage d'utilisation.
\item \textbf{Critère numérique} : calculer le rapport $\dfrac{\Delta R}{\Delta \theta}$ (ou $\dfrac{\Delta U}{\Delta \Phi}$) sur plusieurs intervalles ; s'il est constant, le capteur est linéaire.
\item \textbf{Conséquences pratiques} : un capteur linéaire s'exploite par une simple proportionnalité (étalonnage à 2 points) ; un capteur non linéaire exige une courbe d'étalonnage point par point ou une formule d'inversion.
\item \textbf{Exemples à connaître} : KTY et sonde Pt100 $\approx$ linéaires ; CTN, LDR, thermocouple : non linéaires.
\end{enumerate}
\end{methode}

\begin{exoresolu}[Linéaire ou non linéaire ?]
On relève les caractéristiques de deux capteurs de température :

\begin{center}
\begin{tabular}{|c|c|c|c|c|}
\hline
\rowcolor{popblueL} $\theta$ (\si{\celsius}) & 0 & 20 & 40 & 60 \\
\hline
Capteur A : $R_A$ (\si{\ohm}) & 100 & 108 & 116 & 124 \\
\hline
Capteur B : $R_B$ (\si{\ohm}) & 1000 & 650 & 420 & 270 \\
\hline
\end{tabular}
\end{center}

\begin{enumerate}
\item Déterminer, en le justifiant, lequel des deux capteurs est linéaire.
\item Donner l'équation $R_A = f(\theta)$ du capteur linéaire.
\item Quelle température indique le capteur A lorsqu'on mesure $R_A = \SI{112}{\ohm}$ ?
\end{enumerate}
\tcblower
\textbf{1. Test de linéarité.}

Calculons les variations par pas de $\SI{20}{\celsius}$ :

Capteur A : $\Delta R_A = +8\,\Omega$, $+8\,\Omega$, $+8\,\Omega$ : variation \emph{constante}, donc $\dfrac{\Delta R_A}{\Delta \theta} = \dfrac{8}{20} = \SI{0,40}{\ohm\per\celsius}$ partout. Le capteur A est \cle{linéaire}.

Capteur B : $\Delta R_B = -350\,\Omega$, $-230\,\Omega$, $-150\,\Omega$ : variations non constantes. Le capteur B est \cle{non linéaire} (comportement typique d'une CTN).

\textbf{2. Équation du capteur A.}

Droite de pente $a = \SI{0,40}{\ohm\per\celsius}$ et d'ordonnée à l'origine $R_A(0) = \SI{100}{\ohm}$ :
\[
R_A(\theta) = 100 + 0{,}40\,\theta \quad (R_A \text{ en } \Omega,\ \theta \text{ en } \si{\celsius}).
\]
Vérification dimensionnelle : le produit d'une pente en $\Omega \cdot \si{\celsius}^{-1}$ par une température en $\si{\celsius}$ donne bien des ohms. \checkmark

\textbf{3. Mesure de température.}

On inverse la relation :
\[
\theta = \frac{R_A - 100}{0{,}40} = \frac{112 - 100}{0{,}40} = \frac{12}{0{,}40} = \SI{30}{\celsius}.
\]
\textbf{Conclusion.} Le capteur A (type sonde platine) est linéaire : deux points d'étalonnage suffisent et la mesure $R_A = \SI{112}{\ohm}$ correspond à $\theta = \SI{30}{\celsius}$. Le capteur B, non linéaire, nécessiterait une courbe d'étalonnage complète.
\end{exoresolu}

\courssub{Méthode 5 : Expliquer une chaîne électronique complète (capteur — traitement — actionneur)}

\begin{methode}[Analyser ou concevoir une chaîne électronique]
\begin{enumerate}
\item \textbf{Découper la chaîne en trois blocs fonctionnels} :
\begin{itemize}
\item \cle{Capteur} : convertit la grandeur physique (lumière, température, niveau d'eau, onde sonore) en signal électrique ;
\item \cle{Unité de traitement} : amplifie, compare à un seuil, met en forme (transistor en commutation, amplificateur) ;
\item \cle{Actionneur} : convertit le signal traité en action (sirène, lampe, moteur, écouteur, relais).
\end{itemize}
\item \textbf{Schématiser} la chaîne par des blocs reliés par des flèches, en précisant la nature du signal entre les blocs.
\item \textbf{Justifier chaque étage} : rôle du pont diviseur (conversion $R \to U$), du transistor (commutation tout-ou-rien : bloqué si $U_{BE} < \SI{0,7}{V}$, saturé si $I_B$ est suffisant), du relais (isolement entre circuit de commande basse tension et circuit de puissance).
\item \textbf{Décrire le fonctionnement chronologiquement} : état 1 (grandeur physique faible) puis état 2 (grandeur physique forte), en précisant l'état du transistor et de l'actionneur dans chaque cas.
\item \textbf{Vérifier la cohérence} : la logique (action quand la grandeur augmente ou diminue) doit correspondre au cahier des charges.
\end{enumerate}
\end{methode}

\begin{exoresolu}[Indicateur de niveau d'eau d'un forage]
Pour surveiller le niveau d'eau d'un réservoir alimentant un quartier de Bafoussam, on plonge deux électrodes en cuivre dans la cuve. Quand l'eau atteint les électrodes, sa conductivité (eau non pure) ferme le circuit de base d'un transistor NPN alimenté sous $E = \SI{12}{V}$ ; le transistor sature et alimente une sirène. Un résistor $R_B = \SI{47}{k\ohm}$ protège la base. On donne $U_{BE} = \SI{0,7}{V}$ en conduction et $U_{CE,\text{sat}} \approx \SI{0}{V}$. La sirène a une résistance $R_S = \SI{120}{\ohm}$.
\begin{enumerate}
\item Schématiser la chaîne électronique en trois blocs.
\item Expliquer le fonctionnement quand le niveau d'eau est bas puis quand il atteint les électrodes.
\item Calculer le courant de base $I_B$ et le courant dans la sirène $I_C$ quand l'eau touche les électrodes (résistance de l'eau entre électrodes : $R_{eau} \approx \SI{10}{k\ohm}$).
\end{enumerate}
\tcblower
\textbf{1. Schéma de la chaîne.}

\begin{center}
\begin{tikzpicture}[
 bloc/.style={rectangle, draw=popblue, fill=popblueL, rounded corners, minimum height=1cm, minimum width=2.6cm, align=center, font=\small},
 fleche/.style={-{Stealth[length=3mm]}, thick, popdark}]
\node[bloc, align=center] (cap) {Capteur\\ électrodes + eau};
\node[bloc, right=1.2cm of cap, align=center] (trai) {Traitement\\ transistor NPN};
\node[bloc, right=1.2cm of trai, align=center] (act) {Actionneur\\ sirène};
\draw[fleche] (cap) -- (trai) node[midway, above, font=\small]{$I_B$};
\draw[fleche] (trai) -- (act) node[midway, above, font=\small]{$I_C$};
\end{tikzpicture}
\legende{Chaîne électronique de l'indicateur de niveau d'eau.}
\end{center}

\textbf{2. Fonctionnement.}

\emph{Niveau bas :} l'eau n'atteint pas les électrodes ; le circuit de base est ouvert (l'air est isolant), donc $I_B = 0$ : le transistor est \emph{bloqué}, $I_C = 0$, la sirène est silencieuse.

\emph{Niveau haut :} l'eau (rendue conductrice par les sels minéraux dissous) établit le contact entre les électrodes ; un courant de base $I_B$ circule à travers $R_B$ et l'eau. Le transistor devient passant et se \emph{sature} : $I_C$ circule dans la sirène qui retentit et alerte le fontainier.

\textbf{3. Calcul des courants.}

Maille de base : $E = R_B I_B + R_{eau} I_B + U_{BE}$, d'où :
\[
I_B = \frac{E - U_{BE}}{R_B + R_{eau}} = \frac{12 - 0{,}7}{47\,000 + 10\,000} = \frac{11{,}3}{57\,000} \approx \SI{2,0e-4}{A} = \SI{0,20}{mA}.
\]
Maille de collecteur (transistor saturé, $U_{CE,\text{sat}} \approx 0$) :
\[
I_C = \frac{E}{R_S} = \frac{12}{120} = \SI{0,10}{A} = \SI{100}{mA}.
\]
\textbf{Conclusion.} Un courant de base de seulement $\SI{0,20}{mA}$ commande un courant de $\SI{100}{mA}$ dans la sirène : le transistor joue bien son rôle d'\cle{interrupteur commandé} et d'amplificateur de courant ($I_C / I_B \approx 500$). L'eau, grâce à ses ions, sert ici de capteur de niveau.
\end{exoresolu}

\courssub{Méthode 6 : Expliquer le fonctionnement d'une antenne, d'un microphone et d'un écouteur}

\begin{methode}[Décrire les transducteurs d'ondes]
\begin{enumerate}
\item \textbf{Identifier le sens de la conversion} :
\begin{itemize}
\item \cle{Antenne réceptrice} : onde électromagnétique $\to$ tension alternative induite (phénomène d'induction électromagnétique) ;
\item \cle{Microphone} : onde sonore (pression) $\to$ tension électrique (membrane + bobine mobile dans un aimant, ou capsule capacitive) ;
\item \cle{Écouteur / haut-parleur} : tension électrique $\to$ onde sonore (force de Laplace sur la bobine parcourue par le courant, solidaire de la membrane).
\end{itemize}
\item \textbf{Citer la loi physique mise en jeu} : induction électromagnétique ($e = -\dfrac{d\Phi}{dt}$) pour l'antenne et le microphone électrodynamique ; forces de Laplace ($\vec{F} = I\vec{\ell} \wedge \vec{B}$) pour l'écouteur.
\item \textbf{Préciser la fidélité} : la tension délivrée reproduit les variations (fréquence, amplitude) de l'onde captée, et réciproquement pour l'écouteur.
\item \textbf{Structurer la réponse au Bac} : nommer la grandeur d'entrée, le phénomène physique, la grandeur de sortie, et le rôle dans la chaîne (capteur ou actionneur).
\end{enumerate}
\end{methode}

\begin{exoresolu}[La chaîne radio : de l'émetteur de CRTV à l'auditeur]
Un poste radio reçoit les émissions de la CRTV. On considère la chaîne : antenne réceptrice, amplificateur, écouteur.
\begin{enumerate}
\item Expliquer comment l'antenne transforme l'onde électromagnétique en signal électrique.
\item Expliquer le fonctionnement de l'écouteur (ou haut-parleur) du poste.
\item Situer l'antenne et l'écouteur dans les trois blocs d'une chaîne électronique.
\item Le champ magnétique de l'aimant de l'écouteur vaut $B = \SI{0,40}{T}$ ; la bobine comporte $N = 50$ spires de rayon $r = \SI{1,0}{cm}$ et est parcourue par $I = \SI{20}{mA}$. Calculer la force de Laplace totale exercée sur la bobine (conducteurs perpendiculaires à $\vec{B}$).
\end{enumerate}
\tcblower
\textbf{1. Antenne réceptrice.}

L'onde électromagnétique émise par l'émetteur est constituée de champs $\vec{E}$ et $\vec{B}$ variables. Le champ électrique variable met en mouvement les électrons libres du conducteur de l'antenne ; de plus, la variation du flux magnétique à travers les spires induit une f.é.m. $e = -\dfrac{d\Phi}{dt}$ (loi de Faraday). Il naît ainsi aux bornes de l'antenne une \emph{tension alternative de même fréquence} que l'onde captée : c'est le signal électrique, très faible, qui sera amplifié.

\textbf{2. Écouteur (haut-parleur électrodynamique).}

Le signal amplifié, tension alternative, parcourt une bobine mobile placée dans l'entrefer d'un aimant permanent. La bobine subit des \emph{forces de Laplace} dont le sens et l'intensité varient au rythme du courant. Solidaire d'une membrane, elle la fait vibrer : la membrane comprime l'air et reproduit l'onde sonore. Conversion : énergie électrique $\to$ énergie mécanique (sonore).

\textbf{3. Place dans la chaîne.}

L'antenne est le \cle{capteur} (onde $\to$ signal électrique) ; l'amplificateur et l'étage de démodulation constituent l'\cle{unité de traitement} ; l'écouteur est l'\cle{actionneur} (signal électrique $\to$ son).

\textbf{4. Force de Laplace sur la bobine.}

Longueur totale de conducteur : $\ell = N \times 2\pi r = 50 \times 2\pi \times 0{,}010 = \pi \approx \SI{3,14}{m}$.

Conducteur perpendiculaire à $\vec{B}$ : $F = B I \ell$ :
\[
F = 0{,}40 \times 0{,}020 \times 3{,}14 \approx \SI{2,5e-2}{N}.
\]
Vérification dimensionnelle : $[B I \ell] = T \cdot A \cdot m = \dfrac{N}{A \cdot m} \cdot A \cdot m = N$. \checkmark

\textbf{Conclusion.} Une force d'environ $\SI{25}{mN}$, variable au rythme du courant, suffit à faire vibrer la membrane légère de l'écouteur et à restituer le son. Antenne et écouteur sont deux \emph{transducteurs} réciproques : l'un capte l'onde, l'autre produit le son.
\end{exoresolu}

\begin{attention}[Les 5 erreurs qui coûtent des points au Bac]
\begin{enumerate}
\item \textbf{Inverser le comportement des capteurs.} CTN et LDR : la résistance \emph{diminue} quand la température ou l'éclairement \emph{augmente}. Écrire le contraire fausse toute l'étude du pont diviseur. Vérifie toujours le sens de variation avant de calculer.
\item \textbf{Mal placer le capteur dans le pont diviseur.} La tension prélevée aux bornes de l'élément du bas augmente quand la résistance du haut diminue... et inversement. Avant de conclure sur la logique (allumer la nuit ou le jour), calcule explicitement les \emph{deux} tensions (jour/nuit, chaud/froid) et compare-les au seuil.
\item \textbf{Oublier $U_{BE} = \SI{0,7}{V}$ dans le calcul de $I_B$.} La maille de base s'écrit $E = R_B I_B + U_{BE}$, donc $I_B = \dfrac{E - U_{BE}}{R_B}$, et non $\dfrac{E}{R_B}$. De même, au collecteur saturé : $I_C = \dfrac{E - U_{CE,\text{sat}}}{R_C}$ avec $U_{CE,\text{sat}} \approx 0$.
\item \textbf{Déclarer « linéaire » sans justification.} Le critère exigé est la constance de $\dfrac{\Delta R}{\Delta \theta}$ (ou une caractéristique en droite). Montre au moins deux calculs de pente sur des intervalles différents avant de conclure.
\item \textbf{Confondre capteur et actionneur dans la chaîne.} L'antenne et le microphone sont des \emph{capteurs} (grandeur physique $\to$ signal électrique) ; l'écouteur, la sirène, la lampe et le relais sont des \emph{actionneurs}. Une chaîne électronique se rédige toujours dans l'ordre : capteur $\to$ traitement $\to$ actionneur, en nommant le phénomène physique de chaque conversion (induction, forces de Laplace, effet Joule).
\end{enumerate}
\end{attention}

\newpage

\coursec{Exercices}

\courssub{Je vérifie mes connaissances}

\begin{exercice}[QCM : notions fondamentales]
Pour chaque question, une seule réponse est exacte. Recopie la lettre correspondante.

\textbf{1.} Un dipôle commandé est un dipôle dont la caractéristique électrique dépend :
\begin{itemize}
\item[a)] uniquement de la température ambiante ;
\item[b)] d'une grandeur de commande (mécanique, électrique, lumineuse, thermique...) ;
\item[c)] uniquement de la tension à ses bornes.
\end{itemize}

\textbf{2.} Dans une chaîne électronique, le capteur est l'élément qui :
\begin{itemize}
\item[a)] traite l'information électrique ;
\item[b)] convertit une grandeur physique en signal électrique ;
\item[c)] réalise l'action finale sur le milieu extérieur.
\end{itemize}

\textbf{3.} Une thermistance CTN voit sa résistance :
\begin{itemize}
\item[a)] augmenter quand la température augmente ;
\item[b)] diminuer quand la température augmente ;
\item[c)] rester constante quelle que soit la température.
\end{itemize}

\textbf{4.} Un transistor utilisé en commutation fonctionne :
\begin{itemize}
\item[a)] uniquement en régime linéaire ;
\item[b)] alternativement à l'état bloqué et à l'état saturé ;
\item[c)] toujours à l'état saturé.
\end{itemize}

\textbf{5.} Le relais électromagnétique permet de :
\begin{itemize}
\item[a)] amplifier un courant continu ;
\item[b)] commander un circuit de puissance par un circuit basse tension tout en les isolant électriquement ;
\item[c)] convertir la lumière en courant électrique.
\end{itemize}
\end{exercice}

\begin{exercice}[Vrai ou faux, en justifiant]
Réponds par vrai ou faux, puis justifie chaque réponse en une ou deux phrases.

\textbf{1.} Un capteur est dit linéaire lorsque la courbe donnant le signal électrique en fonction de la grandeur mesurée est une droite passant par l'origine.

\textbf{2.} La résistance d'une photorésistance (LDR) augmente lorsque l'éclairement augmente.

\textbf{3.} Un rhéostat monté en série dans un circuit permet de faire varier l'intensité du courant.

\textbf{4.} Un microphone est un actionneur car il produit un signal électrique.

\textbf{5.} Dans un transistor NPN, un courant de base nul impose un courant de collecteur nul : le transistor est bloqué.
\end{exercice}

\begin{exercice}[Définitions et vocabulaire]
\textbf{1.} Définis les termes suivants : \cle{dipôle commandé} ; \cle{capteur} ; \cle{actionneur}.

\textbf{2.} Explique en quelques phrases le principe général du captage d'une grandeur physique.

\textbf{3.} Pour chacun des dipôles suivants, indique la grandeur de commande et précise s'il s'agit d'une commande manuelle ou électrique : le rhéostat ; la thermistance CTN ; la photorésistance ; le relais électromagnétique ; le transistor.

\textbf{4.} Nomme les trois blocs fonctionnels d'une chaîne électronique et donne le rôle de chacun.
\end{exercice}

\begin{exercice}[Calculs directs]
\textbf{1.} Un rhéostat est monté en série avec un conducteur ohmique de résistance $R = 20~\Omega$ sous une tension $U = 12~\text{V}$. Calcule l'intensité du courant lorsque la résistance du rhéostat vaut $R_h = 10~\Omega$.

\textbf{2.} La résistance d'une photorésistance vaut $R_1 = 2{,}0~\text{k}\Omega$ en plein jour et $R_2 = 500~\text{k}\Omega$ dans l'obscurité. Elle est parcourue par un courant d'intensité $I = 1{,}0~\text{mA}$. Calcule la tension à ses bornes dans chaque cas et commente l'ordre de grandeur obtenu dans l'obscurité.

\textbf{3.} Un transistor de gain en courant $\beta = 100$ est parcouru par un courant de base $I_B = 0{,}20~\text{mA}$ en régime linéaire. Calcule le courant de collecteur $I_C$.

\textbf{4.} Une thermistance présente une résistance $R = 4{,}7~\text{k}\Omega$ à $\theta = 25~^\circ\text{C}$. Elle est associée en série avec un résistor $R_0 = 4{,}7~\text{k}\Omega$ aux bornes d'un générateur de f.é.m. $E = 9~\text{V}$. Calcule la tension aux bornes de la thermistance.
\end{exercice}

\courssub{Je m'entraîne}

\begin{exercice}[Rhéostat et lampe]
Une lampe portant les indications $(6~\text{V}\,;\,0{,}5~\text{A})$ est branchée en série avec un rhéostat aux bornes d'un générateur de tension continue $E = 12~\text{V}$ supposée constante.

\textbf{1.} Fais le schéma du montage.

\textbf{2.} Calcule la valeur de la résistance $R_h$ à régler sur le rhéostat pour que la lampe fonctionne dans ses conditions nominales.

\textbf{3.} Calcule la puissance électrique dissipée par effet Joule dans le rhéostat dans ces conditions.

\textbf{4.} Calcule le rendement énergétique du montage (rapport de la puissance utile de la lampe à la puissance fournie par le générateur).

\textbf{5.} Propose un montage (potentiométrique) permettant de faire varier la tension aux bornes de la lampe de $0$ à $12~\text{V}$ et explique son intérêt par rapport au montage série.
\end{exercice}

\begin{exercice}[Exploitation de la caractéristique d'une CTN]
La courbe $R = f(\theta)$ d'une thermistance CTN donne les valeurs suivantes :

\begin{center}
\begin{tabular}{|l|c|c|c|c|c|c|}
\hline
\rowcolor{popblueL} $\theta$ ($^\circ$C) & 0 & 20 & 40 & 60 & 80 & 100 \\
\hline
$R$ (k$\Omega$) & 27{,}2 & 12{,}1 & 5{,}9 & 3{,}0 & 1{,}7 & 1{,}0 \\
\hline
\end{tabular}
\end{center}

\textbf{1.} Trace la courbe $R = f(\theta)$ sur papier millimétré (échelles : $1~\text{cm}$ pour $10~^\circ\text{C}$ ; $1~\text{cm}$ pour $2~\text{k}\Omega$).

\textbf{2.} Détermine graphiquement la résistance de la CTN à $\theta = 40~^\circ\text{C}$ et à $\theta = 50~^\circ\text{C}$.

\textbf{3.} Montre, à l'aide de deux intervalles de température de même amplitude, que ce capteur n'est pas linéaire.

\textbf{4.} La CTN est placée en série avec un résistor $R_0 = 5{,}9~\text{k}\Omega$ sous une tension $E = 9~\text{V}$. Calcule la tension aux bornes de la CTN lorsque $\theta = 40~^\circ\text{C}$.
\end{exercice}

\begin{exercice}[Photorésistance et éclairement]
On relève, pour une photorésistance parcourue par un courant d'intensité $I = 2{,}0~\text{mA}$, les tensions suivantes en fonction de l'éclairement $\Phi$ :

\begin{center}
\begin{tabular}{|l|c|c|c|c|}
\hline
\rowcolor{popblueL} $\Phi$ (lux) & 10 & 100 & 500 & 1000 \\
\hline
$U$ (V) & 60 & 8{,}0 & 2{,}0 & 1{,}0 \\
\hline
\end{tabular}
\end{center}

\textbf{1.} Calcule la résistance $R$ de la photorésistance pour chaque éclairement.

\textbf{2.} Conclus sur le sens de variation de $R$ lorsque l'éclairement augmente. Ce capteur est-il linéaire ? Justifie.

\textbf{3.} Cette photorésistance commande la base d'un transistor qui sature lorsque la tension aux bornes de la LDR devient supérieure à $3~\text{V}$. Prévois l'état du transistor (bloqué ou saturé) de jour ($\Phi = 1000$ lux) et de nuit ($\Phi = 10$ lux).

\textbf{4.} Propose une application concrète de ce montage dans une ville camerounaise.
\end{exercice}

\begin{exercice}[Chaîne de réception radio]
Un poste de radio utilisé dans un village de la région de l'Adamaoua comporte les éléments suivants : antenne, circuit d'accord, amplificateur, écouteur.

\textbf{1.} Associe chaque élément à l'un des trois blocs fonctionnels d'une chaîne électronique (capteur, traitement, actionneur). Justifie chaque association.

\textbf{2.} Explique le rôle de l'antenne : quelle grandeur physique capte-t-elle et quel signal électrique fournit-elle ?

\textbf{3.} Explique les conversions d'énergie mises en jeu :
\begin{itemize}
\item[a)] au niveau du microphone d'un studio de la CRTV ;
\item[b)] au niveau de l'écouteur du poste récepteur.
\end{itemize}

\textbf{4.} Schématise la chaîne électronique complète sous forme de blocs reliés par des flèches.
\end{exercice}

\courssub{Je me prépare au Bac}

\begin{exercice}[Type Bac : alarme thermique d'une maternité à Garoua]
Pour surveiller la température d'une couveuse dans un hôpital de Garoua, un technicien utilise une thermistance CTN dont la caractéristique $R = f(\theta)$ est donnée ci-dessous :

\begin{center}
\begin{tabular}{|l|c|c|c|c|c|c|}
\hline
\rowcolor{popblueL} $\theta$ ($^\circ$C) & 20 & 30 & 37 & 40 & 50 & 60 \\
\hline
$R$ (k$\Omega$) & 12{,}0 & 8{,}0 & 6{,}0 & 5{,}0 & 3{,}4 & 2{,}4 \\
\hline
\end{tabular}
\end{center}

La CTN est associée en série avec un résistor $R_0 = 6{,}0~\text{k}\Omega$ ; l'ensemble est alimenté par une source de tension continue $E = 9{,}0~\text{V}$. La tension $u$ aux bornes de la CTN est appliquée à l'entrée d'un circuit de traitement qui déclenche une alarme sonore dès que $u < 3{,}0~\text{V}$.

\textbf{1.} Définis les termes \cle{capteur} et \cle{dipôle commandé}. Pourquoi la CTN est-elle un dipôle commandé ? Quelle est sa grandeur de commande ?

\textbf{2.} Trace la courbe $R = f(\theta)$ (échelles : $1~\text{cm}$ pour $5~^\circ\text{C}$ ; $1~\text{cm}$ pour $1~\text{k}\Omega$) et détermine graphiquement la résistance de la CTN à $\theta = 37~^\circ\text{C}$ (température normale de la couveuse).

\textbf{3.} Montre, par un argument quantitatif, que ce capteur de température n'est pas linéaire.

\textbf{4.} Calcule la tension $u$ aux bornes de la CTN lorsque $\theta = 37~^\circ\text{C}$. L'alarme se déclenche-t-elle ?

\textbf{5.} Détermine la valeur maximale de la résistance de la CTN pour laquelle l'alarme se déclenche, puis la température de déclenchement de l'alarme à l'aide de la courbe.

\textbf{6.} Identifie les trois blocs fonctionnels de cette chaîne électronique et schématise-la.
\end{exercice}

\begin{exercice}[Type Bac : éclairage automatique d'un lycée à Douala]
Le proviseur d'un lycée de Douala souhaite que les lampes de la cour s'allument automatiquement à la tombée de la nuit. On utilise une photorésistance (LDR) dont la résistance vaut $R_J = 1{,}0~\text{k}\Omega$ en plein jour et $R_N = 100~\text{k}\Omega$ la nuit. La LDR est placée en série avec un résistor $R_0 = 10~\text{k}\Omega$ sous une tension $E = 12~\text{V}$ (le résistor $R_0$ est relié au pôle $+$, la LDR à la masse) ; la tension aux bornes de la LDR attaque, à travers un résistor de base $R_B = 10~\text{k}\Omega$, la base d'un transistor NPN de gain $\beta = 200$, dont le collecteur alimente la bobine d'un relais électromagnétique. Le relais ferme le circuit de puissance ($220~\text{V}$) des lampes. On donne : tension base-émetteur à l'état conducteur $U_{BE} = 0{,}7~\text{V}$ ; courant nécessaire au collage du relais $I_{C(\text{sat})} = 50~\text{mA}$ ; tension de saturation $U_{CE(\text{sat})} \approx 0~\text{V}$.

\textbf{1.} Définis une photorésistance et précise sa grandeur de commande.

\textbf{2.} Calcule la tension $u$ aux bornes de la LDR de jour, puis la nuit.

\textbf{3.} Calcule le courant de base minimal $I_{B(\text{min})}$ nécessaire pour saturer le transistor et faire coller le relais.

\textbf{4.} Calcule le courant de base réellement fourni de jour, puis la nuit. En déduire l'état du transistor (bloqué, en régime linéaire ou saturé) dans chaque cas, puis l'état des lampes. Le cahier des charges du proviseur est-il respecté ?

\textbf{5.} Explique le rôle du relais électromagnétique dans cette chaîne et justifie son intérêt en termes d'isolement entre le circuit de commande (basse tension) et le circuit de puissance ($220~\text{V}$).

\textbf{6.} Schématise la chaîne électronique complète en identifiant le capteur, l'unité de traitement et l'actionneur.
\end{exercice}

\begin{exercice}[Type Bac : indicateur de niveau d'eau d'un forage à Maroua]
Une coopérative agricole de Maroua équipe un château d'eau d'un indicateur de niveau. Deux électrodes métalliques $A$ (au fond) et $B$ (à hauteur du niveau maximal) plongent dans la cuve. L'électrode $A$ est reliée au pôle $+$ d'un générateur de f.é.m. $E = 6{,}0~\text{V}$ ; l'électrode $B$ est reliée à la base d'un transistor NPN ($\beta = 150$, $U_{BE} = 0{,}7~\text{V}$) à travers un résistor de protection $R_B = 10~\text{k}\Omega$. Le collecteur alimente une lampe témoin $(6~\text{V}\,;\,0{,}30~\text{W})$ ; l'émetteur est à la masse. L'eau de la cuve est légèrement conductrice : lorsque le niveau atteint l'électrode $B$, la résistance équivalente de l'eau entre $A$ et $B$ vaut $R_{\text{eau}} = 20~\text{k}\Omega$.

\textbf{1.} Explique le rôle des électrodes dans ce dispositif : pourquoi peut-on les considérer comme un dipôle commandé ? Quelle est la grandeur de commande ?

\textbf{2.} Décris le fonctionnement du montage :
\begin{itemize}
\item[a)] lorsque le niveau d'eau est inférieur à l'électrode $B$ ;
\item[b)] lorsque le niveau monte et atteint l'électrode $B$.
\end{itemize}

\textbf{3.} Calcule l'intensité du courant de base $I_B$ lorsque l'eau atteint l'électrode $B$.

\textbf{4.} Calcule le courant de collecteur en régime linéaire $I_C = \beta I_B$, puis le courant de saturation $I_{C(\text{sat})}$ correspondant au fonctionnement nominal de la lampe. Le transistor est-il saturé ? La lampe brille-t-elle normalement ?

\textbf{5.} Identifie les trois blocs fonctionnels de cette chaîne électronique et propose une modification du montage pour déclencher, au lieu d'une lampe, une pompe électrique de puissance alimentée sous $220~\text{V}$.
\end{exercice}



\newpage

\coursec{Corrigés des exercices}

\coursec{Leçon P13 : Dipôles commandés, capteurs et chaînes électroniques}

\courssub{Je vérifie mes connaissances}

\begin{corrige}[Exercice 1 — QCM : notions fondamentales]
\textbf{1. b)} : par définition, un \cle{dipôle commandé} possède une caractéristique électrique qui dépend d'une grandeur de commande extérieure.

\textbf{2. b)} : le \cle{capteur} est le transducteur d'entrée ; le traitement est assuré par l'unité de traitement, l'action finale par l'\cle{actionneur}.

\textbf{3. b)} : CTN signifie \og coefficient de température négatif \fg{} : la résistance $R$ diminue quand la température $\theta$ augmente.

\textbf{4. b)} : en régime de \cle{commutation}, le transistor ne connaît que deux états : bloqué ($I_B = 0$, $I_C = 0$) ou saturé ($I_C = I_{C(\text{sat})}$, $U_{CE} \approx 0$).

\textbf{5. b)} : le \cle{relais électromagnétique} assure l'isolement galvanique entre le circuit de commande (basse tension) et le circuit de puissance.
\end{corrige}

\begin{attention}[Pièges classiques du QCM]
\begin{itemize}
\item Ne pas confondre \emph{capteur} (entrée : physique $\rightarrow$ électrique) et \emph{actionneur} (sortie : électrique $\rightarrow$ physique).
\item Un capteur \emph{linéaire} au sens strict du Bac : caractéristique \textbf{droite passant par l'origine} (proportionnalité). Une droite ne passant pas par l'origine est \emph{affine}, souvent qualifiée de « linéaire » dans l'industrie mais pas au programme.
\item Relais : l'isolement galvanique est son atout majeur pour la sécurité (commande 12 V / puissance 220 V).
\end{itemize}
\end{attention}

\begin{corrige}[Exercice 2 — Vrai ou faux, en justifiant]
\textbf{1. Vrai} (dans le sens strict du programme) : la linéarité signifie la proportionnalité entre le signal de sortie et la grandeur mesurée, donc une droite passant par l'origine. On admet parfois, dans un langage plus large, qu'un capteur est linéaire si sa caractéristique est une droite (relation affine) sur la plage d'utilisation.

\textbf{2. Faux} : la résistance d'une \cle{photorésistance (LDR)} \emph{diminue} lorsque l'éclairement augmente (grande valeur dans l'obscurité, faible valeur en pleine lumière).

\textbf{3. Vrai} : en série, le rhéostat ajoute une résistance variable $R_h$ ; d'après la loi d'Ohm $I = \dfrac{U}{R + R_h}$, faire varier $R_h$ fait varier $I$.

\textbf{4. Faux} : le microphone convertit une grandeur physique (onde sonore, pression acoustique) en signal électrique : c'est un \cle{capteur}. Un actionneur (haut-parleur, moteur, lampe, sirène) réalise l'action finale sur le milieu.

\textbf{5. Vrai} : si $I_B = 0$, alors $I_C = \beta I_B = 0$ (modèle idéal) : le transistor est bloqué et se comporte comme un interrupteur ouvert entre collecteur et émetteur.
\end{corrige}

\begin{corrige}[Exercice 3 — Définitions et vocabulaire]
\textbf{1.} Un \cle{dipôle commandé} est un dipôle dont la caractéristique électrique (résistance, tension, courant) dépend d'une grandeur extérieure appelée \cle{grandeur de commande}. Un \cle{capteur} est un dispositif qui convertit une grandeur physique (température, lumière, pression, onde\ldots) en un signal électrique mesurable. Un \cle{actionneur} est un dispositif qui, à partir d'un signal électrique, réalise une action sur le milieu extérieur (lampe, moteur, sirène, écouteur, relais\ldots).

\textbf{2.} Le captage repose sur un phénomène physique (effet thermique pour la CTN, photoélectrique pour la LDR, piézoélectrique pour le microphone, induction électromagnétique pour l'antenne\ldots) par lequel la grandeur à mesurer modifie une propriété électrique d'un composant (résistance, tension, charge). La variation de cette propriété est ensuite convertie en signal électrique exploitable (tension ou courant).

\textbf{3.} 
\begin{center}
\begin{tabularx}{\linewidth}{|l|X|X|}\hline
\rowcolor{popblueL} \textbf{Dipôle commandé} & \textbf{Grandeur de commande} & \textbf{Type de commande} \\ \hline
Rhéostat & Position du curseur & Manuelle \\ \hline
Thermistance CTN & Température $\theta$ & Électrique (automatique) \\ \hline
Photorésistance (LDR) & Éclairement $\Phi$ & Électrique (automatique) \\ \hline
Relais électromagnétique & Courant dans la bobine $i_b$ & Électrique \\ \hline
Transistor (BJT) & Courant de base $I_B$ & Électrique \\ \hline
\end{tabularx}
\end{center}

\textbf{4.} Les trois maillons d'une \cle{chaîne électronique} : 
\begin{enumerate}
\item le \cle{capteur} (acquisition de la grandeur physique $\rightarrow$ signal électrique) ;
\item l'\cle{unité de traitement} (amplification, filtrage, comparaison, mise en forme, commutation) ;
\item l'\cle{actionneur} (action finale sur le milieu extérieur).
\end{enumerate}
\end{corrige}

\begin{savaistu}[Le relais, un héritage industriel camerounais]
Les relais électromagnétiques équipent encore de nombreux tableaux de commande d'ENEO (postes de transformation, démarrage moteurs de pompage) et les automatismes des barrages (Lom Pangar, Nachtigal). Leur robustesse et leur isolement galvanique parfait en font le standard pour commander du 220 V / 380 V depuis une électronique 12 V / 24 V, malgré l'essor des transistors de puissance (MOSFET, IGBT).
\end{savaistu}

\begin{corrige}[Exercice 4 — Calculs directs]
\textbf{1.} Les résistances s'ajoutent en série : $I = \dfrac{U}{R + R_h} = \dfrac{\SI{12}{\volt}}{\SI{20}{\ohm} + \SI{10}{\ohm}} = \SI{0,40}{\ampere}$.

\textbf{2.} Loi d'Ohm : $U_1 = R_1 I = \SI{2,0}{\kilo\ohm} \times \SI{1,0}{\milli\ampere} = \SI{2,0}{\volt}$ ; $U_2 = R_2 I = \SI{500}{\kilo\ohm} \times \SI{1,0}{\milli\ampere} = \SI{500}{\volt}$. Cette valeur très élevée montre qu'en pratique on ne peut pas imposer $\SI{1,0}{\milli\ampere}$ à une LDR dans l'obscurité avec une alimentation ordinaire : on utilise plutôt un montage \emph{diviseur de tension}.

\textbf{3.} $I_C = \beta I_B = 100 \times \SI{0,20}{\milli\ampere} = \SI{20}{\milli\ampere}$.

\textbf{4.} Diviseur de tension : $u = E \times \dfrac{R}{R + R_0} = \SI{9}{\volt} \times \dfrac{\SI{4,7}{\kilo\ohm}}{\SI{4,7}{\kilo\ohm} + \SI{4,7}{\kilo\ohm}} = \SI{4,5}{\volt}$.
\end{corrige}

\courssub{Je m'entraîne}

\begin{corrige}[Exercice 5 — Rhéostat et lampe]
\textbf{1.} Schéma : générateur de tension $E = \SI{12}{\volt}$, rhéostat $R_h$ et lampe $L$ (\SI{6}{\volt} -- \SI{0,5}{\ampere}) en série, formant une boucle fermée.

\textbf{2.} En régime nominal, la lampe est traversée par $I = \SI{0,5}{\ampere}$ et sa tension est $U_L = \SI{6}{\volt}$. Le rhéostat absorbe la différence : $U_h = E - U_L = \SI{6}{\volt}$. D'où $R_h = \dfrac{U_h}{I} = \dfrac{\SI{6}{\volt}}{\SI{0,5}{\ampere}} = \SI{12}{\ohm}$.

\textbf{3.} Puissance dissipée par effet Joule dans le rhéostat : $P_J = R_h I^2 = \SI{12}{\ohm} \times (\SI{0,5}{\ampere})^2 = \SI{3,0}{\watt}$ (ou $P_J = U_h I = \SI{6}{\volt} \times \SI{0,5}{\ampere} = \SI{3,0}{\watt}$).

\textbf{4.} Puissance utile (lampe) : $P_u = U_L I = \SI{6}{\volt} \times \SI{0,5}{\ampere} = \SI{3,0}{\watt}$. Puissance fournie par le générateur : $P_g = E I = \SI{12}{\volt} \times \SI{0,5}{\ampere} = \SI{6,0}{\watt}$. Rendement : $\eta = \dfrac{P_u}{P_g} = 0,50$, soit \SI{50}{\percent}.

\textbf{5.} \textbf{Montage potentiométrique} : les deux extrémités du rhéostat sont branchées aux bornes du générateur et la lampe est branchée entre le curseur et une extrémité. 
\begin{itemize}
\item La tension aux bornes de la lampe varie alors de \SI{0}{\volt} à \SI{12}{\volt} (impossible en montage série, où la tension minimale reste $>0$).
\item Le réglage est plus fin et la lampe peut être complètement éteinte.
\item Inconvénient : le courant total fourni par le générateur varie (maximum quand le curseur est au milieu), ce qui fatigue la source.
\end{itemize}
\end{corrige}

\begin{methode}[Choisir le montage du rhéostat]
\begin{enumerate}
\item \textbf{Série (rhéostat)} : pour limiter le courant dans une charge, simple, mais tension résiduelle non nulle $\rightarrow$ lampe ne s'éteint pas totalement. Rendement médiocre (pertes dans $R_h$).
\item \textbf{Potentiométrique} : pour faire varier la tension aux bornes d'une charge de \SI{0}{\volt} à $E$. Meilleur contrôle, extinction possible. Courant source variable.
\item \textbf{Électronique (PWM, transistor)} : rendement optimal, commande numérique, standard moderne (variateurs de vitesse, gradateurs LED).
\end{enumerate}
\end{methode}

\begin{corrige}[Exercice 6 — Exploitation de la caractéristique d'une CTN]
\textbf{1.} La courbe $R = f(\theta)$ est \emph{décroissante} (CTN), fortement pentue aux basses températures puis s'aplatissant : ce n'est pas une droite, le capteur est \cle{non linéaire}.

\textbf{2.} Lecture graphique (interpolation) : à $\theta = \SI{40}{\celsius}$, $R \approx \SI{5,9}{\kilo\ohm}$ ; à $\theta = \SI{50}{\celsius}$, $R \approx \SI{4,2}{\kilo\ohm}$ (interpolation entre $\SI{5,9}{\kilo\ohm}$ à \SI{40}{\celsius} et $\SI{3,0}{\kilo\ohm}$ à \SI{60}{\celsius}).

\textbf{3.} Preuve de non-linéarité : 
\begin{itemize}
\item Entre \SI{0}{\celsius} et \SI{20}{\celsius} ($\Delta\theta = \SI{20}{\celsius}$) : $\Delta R = \SI{12,1}{\kilo\ohm} - \SI{27,2}{\kilo\ohm} = \SI{-15,1}{\kilo\ohm}$.
\item Entre \SI{80}{\celsius} et \SI{100}{\celsius} (même $\Delta\theta = \SI{20}{\celsius}$) : $\Delta R = \SI{1,0}{\kilo\ohm} - \SI{1,7}{\kilo\ohm} = \SI{-0,7}{\kilo\ohm}$.
\end{itemize}
Des variations de température égales ne produisent \textbf{pas} des variations de résistance égales : le capteur est \cle{non linéaire}.

\textbf{4.} Montage diviseur de tension avec $R_0 = \SI{5,9}{\kilo\ohm}$ (égale à $R$ à \SI{40}{\celsius}) : $u = E \times \dfrac{R}{R + R_0} = \SI{9}{\volt} \times \dfrac{\SI{5,9}{\kilo\ohm}}{\SI{5,9}{\kilo\ohm} + \SI{5,9}{\kilo\ohm}} = \SI{4,5}{\volt}$.
\end{corrige}

\begin{attention}[Linéarisation d'un capteur non linéaire]
Au Bac, on peut vous demander de \emph{linéariser} la réponse : 
\begin{itemize}
\item Choix judicieux de $R_0$ (souvent $R_0 = R(\theta_{\text{milieu de plage}})$) pour obtenir une tension $u(\theta)$ \emph{quasi-linéaire} sur une plage réduite.
\item Utilisation d'un montage pont de Wheatstone + amplificateur d'instrumentation.
\item Traitement logiciel (table de correspondance, polynôme) dans un microcontrôleur (Arduino, STM32).
\end{itemize}
\end{attention}

\begin{corrige}[Exercice 7 — Photorésistance et éclairement]
\textbf{1.} La caractéristique donne $U = f(I, \Phi)$. Pour un courant imposé $I = \SI{2,0}{\milli\ampere}$, on lit les tensions et on calcule $R = \dfrac{U}{I}$ :
\begin{align*}
\Phi = \SI{10}{\lux} &: U = \SI{60}{\volt} \Rightarrow R = \frac{\SI{60}{\volt}}{\SI{2,0}{\milli\ampere}} = \SI{30}{\kilo\ohm} ;\\
\Phi = \SI{100}{\lux} &: U = \SI{8,0}{\volt} \Rightarrow R = \SI{4,0}{\kilo\ohm} ;\\
\Phi = \SI{500}{\lux} &: U = \SI{2,0}{\volt} \Rightarrow R = \SI{1,0}{\kilo\ohm} ;\\
\Phi = \SI{1000}{\lux} &: U = \SI{1,0}{\volt} \Rightarrow R = \SI{0,50}{\kilo\ohm}.
\end{align*}

\textbf{2.} $R$ diminue quand l'éclairement augmente. Entre \SI{10}{\lux} et \SI{100}{\lux}, $R$ chute de \SI{26}{\kilo\ohm} ; entre \SI{500}{\lux} et \SI{1000}{\lux}, de seulement \SI{0,5}{\kilo\ohm} : la variation n'est pas proportionnelle, le capteur est \cle{non linéaire}.

\textbf{3.} On utilise la tension $U$ aux bornes de la LDR (parcourue par $I = \SI{2}{\milli\ampere}$) pour piloter la base d'un transistor NPN via une résistance $R_B$.
\begin{itemize}
\item \textbf{De jour} ($\Phi = \SI{1000}{\lux}$) : $R_{\text{LDR}}$ est faible et $U = \SI{1,0}{\volt} < \SI{3}{\volt}$ : le transistor est \textbf{bloqué}.
\item \textbf{De nuit} ($\Phi = \SI{10}{\lux}$) : $R_{\text{LDR}}$ est grande et $U = \SI{60}{\volt} > \SI{3}{\volt}$ : le transistor est \textbf{saturé}.
\end{itemize}
\emph{Remarque :} ce montage allume la charge la nuit (transistor saturé $\rightarrow$ relais collé $\rightarrow$ lampe allumée), ce qui correspond à un éclairage public automatique.

\textbf{4.} Applications camerounaises : éclairage automatique des rues de Yaoundé, Douala, Bafoussam à la tombée de la nuit ; extinction automatique des lampes des lycées et campus universitaires au lever du jour pour économiser l'énergie (facture ENEO).
\end{corrige}

\begin{popfigure}
\begin{tikzpicture}[scale=0.9, every node/.style={font=\small}]
\draw (0,5) node[left] {$+E = \SI{12}{V}$} -- (7,5);
\draw (0,5) to[R, l_={$R_0$}] (0,2.5) to[photoresistor, l_={$R_{\text{LDR}}$}] (0,0);
\node[npn] (Q) at (3.5,2.5) {};
\draw (0,2.5) to[R, l={$R_B$}] (Q.B);
\draw (Q.C) to[L, l_=bobine du relais] (Q.C |- 0,5);
\draw (Q.E) -- (Q.E |- 0,0);
\draw (0,0) -- (4.5,0);
\draw (3.5,0) node[ground] {};
\draw (6,4) to[nos, l=contact] (6,2) to[lamp, l_=lampes] (6,0.6);
\draw (6,4) -- (7.2,4) node[right] {\SI{220}{V}};
\draw (6,0.6) -- (7.2,0.6);
\end{tikzpicture}
\legende{Schéma principe : commande d'éclairage public par LDR, transistor et relais (Ex. 7 \& 10).}
\end{popfigure}

\begin{corrige}[Exercice 8 — Chaîne de réception radio]
\textbf{1.} 
\begin{itemize}
\item \textbf{Antenne} : \cle{capteur} (elle convertit l'onde électromagnétique en signal électrique alternatif de même fréquence, de très faible amplitude).
\item \textbf{Circuit d'accord + amplificateur} : \cle{unité de traitement} (sélection de la station par résonance, amplification du signal utile, démodulation).
\item \textbf{Écouteur} : \cle{actionneur} (il convertit le signal électrique en ondes sonores).
\end{itemize}

\textbf{2.} L'antenne capte les ondes hertziennes émises par la station de radio ; le champ électromagnétique variable induit dans le conducteur de l'antenne une tension alternative de même fréquence (phénomène d'induction / loi de Faraday) : c'est le signal électrique utile, de très faible amplitude (µV à mV).

\textbf{3. a)} \textbf{Microphone} (ex. : à électret ou dynamique) : l'énergie mécanique (acoustique) de l'onde sonore fait vibrer une membrane (ou un diaphragme) ; cette vibration est convertie en signal électrique (variation de capacité ou induction) : conversion \textbf{mécanique $\rightarrow$ électrique}.

\textbf{b)} \textbf{Écouteur} (casque, oreillette) : le courant électrique variable traverse une bobine placée dans un champ magnétique (aimant permanent) ; la force de Laplace fait vibrer une membrane solidaire : conversion \textbf{électrique $\rightarrow$ mécanique (acoustique)}. Microphone et écouteur réalisent des conversions \emph{inverses} (réciprocité).

\textbf{4.} Chaîne complète :
\[
\text{Onde hertzienne} \;\rightarrow\; \fbox{Antenne (capteur)} \;\rightarrow\; \fbox{Circuit d'accord + Amplificateur + Démodulateur (traitement)} \;\rightarrow\; \fbox{Écouteur (actionneur)} \;\rightarrow\; \text{Son}
\]
\end{corrige}

\courssub{Je me prépare au Bac}

\begin{corrige}[Exercice 9 — Type Bac : alarme thermique d'une maternité à Garoua]
\textbf{1.} Un \cle{capteur} convertit une grandeur physique en signal électrique. Un \cle{dipôle commandé} est un dipôle dont la caractéristique électrique dépend d'une grandeur de commande extérieure. La résistance de la CTN dépend de la température $\theta$ : c'est un dipôle commandé dont la grandeur de commande est la température.

\textbf{2.} La courbe $R = f(\theta)$ est décroissante, non rectiligne (non linéaire). Lecture à $\theta = \SI{37}{\celsius}$ : $R \approx \SI{6,0}{\kilo\ohm}$.

\textbf{3.} Preuve de non-linéarité : 
\begin{itemize}
\item Entre \SI{20}{\celsius} et \SI{30}{\celsius} ($\Delta\theta = \SI{10}{\celsius}$) : $\Delta R = \SI{8,0}{\kilo\ohm} - \SI{12,0}{\kilo\ohm} = \SI{-4,0}{\kilo\ohm}$.
\item Entre \SI{50}{\celsius} et \SI{60}{\celsius} (même $\Delta\theta = \SI{10}{\celsius}$) : $\Delta R = \SI{2,4}{\kilo\ohm} - \SI{3,4}{\kilo\ohm} = \SI{-1,0}{\kilo\ohm}$.
\end{itemize}
À variation de température égale, les variations de résistance diffèrent : le capteur est \cle{non linéaire}.

\textbf{4.} Diviseur de tension ($R_0 = \SI{6,0}{\kilo\ohm}$, $E = \SI{9,0}{\volt}$) : 
\[ u = E \times \frac{R}{R + R_0} = \SI{9,0}{\volt} \times \frac{\SI{6,0}{\kilo\ohm}}{\SI{6,0}{\kilo\ohm} + \SI{6,0}{\kilo\ohm}} = \SI{4,5}{\volt}. \]
Comme $u = \SI{4,5}{\volt} > \SI{3,0}{\volt}$ (seuil de déclenchement), l'alarme \textbf{ne se déclenche pas} : la température est normale.

\textbf{5.} L'alarme se déclenche si $u < \SI{3,0}{\volt}$, c'est-à-dire :
\[ E \frac{R}{R + R_0} < 3,0 \;\Longleftrightarrow\; \SI{9,0}{\volt}\,R < \SI{3,0}{\volt}\,(R + \SI{6,0}{\kilo\ohm}) \;\Longleftrightarrow\; \SI{6,0}{\volt}\,R < \SI{18}{\volt\cdot\kilo\ohm} \;\Longleftrightarrow\; R < \SI{3,0}{\kilo\ohm}. \]
La valeur maximale est donc $R_{\max} = \SI{3,0}{\kilo\ohm}$. Sur la courbe, $R = \SI{3,0}{\kilo\ohm}$ correspond (interpolation linéaire entre $\SI{3,4}{\kilo\ohm}$ à \SI{50}{\celsius} et $\SI{2,4}{\kilo\ohm}$ à \SI{60}{\celsius}) à $\theta \approx \SI{54}{\celsius}$ : l'alarme se déclenche dès que la température dépasse environ \SI{54}{\celsius}, signe d'une surchauffe dangereuse de la couveuse.

\textbf{6.} Chaîne électronique :
\begin{itemize}
\item \textbf{Capteur} : la CTN (température $\rightarrow$ résistance).
\item \textbf{Traitement} : le diviseur de tension ($R_0$ + CTN) qui convertit $R$ en tension $u$, suivi d'un comparateur (ou transistor en commutation) qui surveille $u$ par rapport au seuil \SI{3,0}{\volt}.
\item \textbf{Actionneur} : l'alarme sonore (buzzer, sirène).
\end{itemize}
Schéma fonctionnel : 
\[
\text{Température} \;\rightarrow\; \fbox{CTN} \;\rightarrow\; \fbox{Diviseur + Comparateur} \;\rightarrow\; \fbox{Alarme sonore}
\]
\end{corrige}

\begin{corrige}[Exercice 10 — Type Bac : éclairage automatique d'un lycée à Douala]
\textbf{1.} Une \cle{photorésistance (LDR)} est un dipôle commandé dont la résistance diminue lorsque l'éclairement augmente ; sa grandeur de commande est l'éclairement lumineux $\Phi$.

\textbf{2.} Montage diviseur de tension, la tension $u$ étant prise aux bornes de la LDR ($R_0 = \SI{10}{\kilo\ohm}$ fixe en série, $E = \SI{12}{\volt}$) :
\begin{align*}
\text{Jour (} \Phi \text{ fort, } R_J = \SI{1,0}{\kilo\ohm} \text{)} &: u_J = E \frac{R_J}{R_0 + R_J} = \SI{12}{\volt} \times \frac{\SI{1,0}{\kilo\ohm}}{\SI{10}{\kilo\ohm} + \SI{1,0}{\kilo\ohm}} \approx \SI{1,1}{\volt} ;\\
\text{Nuit (} \Phi \text{ faible, } R_N = \SI{100}{\kilo\ohm} \text{)} &: u_N = E \frac{R_N}{R_0 + R_N} = \SI{12}{\volt} \times \frac{\SI{100}{\kilo\ohm}}{\SI{10}{\kilo\ohm} + \SI{100}{\kilo\ohm}} \approx \SI{10,9}{\volt}.
\end{align*}

\textbf{3.} Pour saturer le transistor, il faut $I_C \geq I_{C(\text{sat})} = \SI{50}{\milli\ampere}$, donc courant de base minimal :
\[ I_{B(\text{min})} = \frac{I_{C(\text{sat})}}{\beta} = \frac{\SI{50}{\milli\ampere}}{200} = \SI{0,25}{\milli\ampere}. \]

\textbf{4.} Le courant de base est $I_B = \dfrac{u - U_{BE}}{R_B}$ (si $u > U_{BE} \approx \SI{0,7}{\volt}$), avec $R_B = \SI{10}{\kilo\ohm}$.

\begin{itemize}
\item \textbf{Jour} : $u_J = \SI{1,1}{\volt}$. $I_B = \dfrac{\SI{1,1}{\volt} - \SI{0,7}{\volt}}{\SI{10}{\kilo\ohm}} = \SI{0,040}{\milli\ampere} = \SI{40}{\micro\ampere}$. Alors $I_C = \beta I_B = 200 \times \SI{0,040}{\milli\ampere} = \SI{8,0}{\milli\ampere} < \SI{50}{\milli\ampere}$ : le transistor n'est \textbf{pas saturé} (régime linéaire, courant insuffisant), le relais ne colle pas : les lampes restent \textbf{éteintes}.
\item \textbf{Nuit} : $u_N = \SI{10,9}{\volt}$. $I_B = \dfrac{\SI{10,9}{\volt} - \SI{0,7}{\volt}}{\SI{10}{\kilo\ohm}} = \SI{1,02}{\milli\ampere} > I_{B(\text{min})}$ : le transistor est \textbf{saturé}, le relais colle et les lampes s'\textbf{allument}.
\end{itemize}
Le cahier des charges est respecté : lampes éteintes le jour, allumées la nuit.

\textbf{5.} Le \cle{relais électromagnétique} est un interrupteur commandé électriquement : le faible courant de collecteur (\SI{50}{\milli\ampere} sous \SI{12}{\volt}) traverse la bobine, crée un champ magnétique qui attire le mobile et ferme les contacts du circuit de puissance (\SI{220}{\volt}). Il assure l'\cle{isolement électrique (galvanique)} entre le circuit de commande basse tension (sans danger, électronique) et le circuit de puissance \SI{220}{\volt} : aucune liaison conductrice directe ne relie les deux circuits, ce qui protège l'électronique de commande et l'utilisateur contre les surtensions et les chocs électriques.

\textbf{6.} Chaîne électronique :
\[
\text{Éclairement} \;\rightarrow\; \fbox{LDR (capteur)} \;\rightarrow\; \fbox{Diviseur + Transistor (traitement / commutation)} \;\rightarrow\; \fbox{Relais + Lampes (actionneur)} \;\rightarrow\; \text{Éclairage de la cour}
\]
\end{corrige}

\begin{savaistu}[Éclairage public intelligent au Cameroun]
À Douala et Yaoundé, les nouveaux candélabres intègrent souvent une \textbf{photodiode} (plus rapide et linéaire qu'une LDR) ou une \textbf{sonde de présence} (radar/PIR) couplée à un \textbf{module radio (LoRa/Sigfox)} pour la télégestion centralisée (détection de panne, variation d'intensité selon l'heure). Le principe de base (capteur $\rightarrow$ traitement $\rightarrow$ actionneur) reste identique.
\end{savaistu}

\begin{corrige}[Exercice 11 — Type Bac : indicateur de niveau d'eau d'un forage à Maroua]
\textbf{1.} Les deux électrodes $A$ et $B$ et l'eau comprise entre elles forment un dipôle dont la résistance dépend du niveau d'eau : circuit ouvert (résistance quasi infinie) si l'eau n'atteint pas $B$, résistance finie $R_{\text{eau}} = \SI{20}{\kilo\ohm}$ si l'eau relie $A$ et $B$. C'est donc un dipôle commandé dont la grandeur de commande est le \cle{niveau d'eau} ; il joue le rôle de \cle{capteur de niveau}.

\textbf{2. a)} Niveau inférieur à $B$ : le circuit de base est ouvert, $I_B = 0$, donc $I_C = 0$ : le transistor est \textbf{bloqué} et la lampe est \textbf{éteinte}.

\textbf{b)} L'eau atteint $B$ : l'eau conductrice ferme le circuit de base ; un courant $I_B$ circule, le transistor conduit et la lampe s'allume : le niveau maximal est atteint (alarme ou arrêt pompe).

\textbf{3.} Maille de base (sens du courant $I_B$) : $E = R_{\text{eau}} I_B + R_B I_B + U_{BE}$, d'où :
\[ I_B = \frac{E - U_{BE}}{R_{\text{eau}} + R_B} = \frac{\SI{6,0}{\volt} - \SI{0,7}{\volt}}{(\SI{20}{\kilo\ohm} + \SI{10}{\kilo\ohm})} = \frac{\SI{5,3}{\volt}}{\SI{30}{\kilo\ohm}} \approx \SI{0,18}{\milli\ampere}. \]

\textbf{4.} Régime linéaire (hypothèse) : $I_C = \beta I_B = 150 \times \SI{0,18}{\milli\ampere} = \SI{27}{\milli\ampere}$.

Courant nominal de la lampe (\SI{6,0}{\volt} -- \SI{0,30}{\watt}) : $I_{C(\text{sat})} = \dfrac{P}{U} = \dfrac{\SI{0,30}{\watt}}{\SI{6,0}{\volt}} = \SI{50}{\milli\ampere}$.

Comme $\beta I_B = \SI{27}{\milli\ampere} < I_{C(\text{sat})} = \SI{50}{\milli\ampere}$, le transistor n'est \textbf{pas saturé} : il fonctionne en \textbf{régime linéaire} (actif). La lampe n'est traversée que par \SI{27}{\milli\ampere} au lieu de \SI{50}{\milli\ampere} : elle s'allume mais brille faiblement (elle reste visible comme témoin, ce qui suffit pour un indicateur visuel).

\textbf{5.} Chaîne électronique :
\begin{itemize}
\item \textbf{Capteur} : les électrodes $A$ et $B$ (niveau d'eau $\rightarrow$ résistance).
\item \textbf{Traitement} : le transistor et son réseau de résistances $R_B$, $R_{\text{eau}}$ (commutation / amplification).
\item \textbf{Actionneur} : la lampe témoin (indicateur visuel).
\end{itemize}
Pour commander une \textbf{pompe sous \SI{220}{\volt}} : on remplace la lampe par la bobine d'un \cle{relais électromagnétique} (commande \SI{12}{\volt} ou \SI{24}{\volt}) dont les contacts de puissance ferment le circuit \SI{220}{\volt} de la pompe. Le relais assure l'isolement entre la commande basse tension et le \SI{220}{\volt}. On peut aussi augmenter le courant de base (diminuer $R_B$ ou utiliser un montage \emph{Darlington} avec deux transistors) pour garantir la saturation du transistor et une lampe à pleine luminosité.
\end{corrige}

\begin{methode}[Analyse d'une chaîne électronique (type Bac)]
\begin{enumerate}
\item \textbf{Identifier les 3 blocs} : Capteur (grandeur physique $\rightarrow$ élec.), Traitement (signal $\rightarrow$ décision/puissance), Actionneur (élec. $\rightarrow$ action physique).
\item \textbf{Analyser le capteur} : Nature (CTN, LDR, électrodes...), grandeur de commande, caractéristique ($R=f(\theta)$, $R=f(\Phi)$...), linéarité.
\item \textbf{Étudier le traitement} : Montage (diviseur, pont, amplificateur, transistor commutation). Calculer tensions/courants seuils. Vérifier saturation ($I_B > I_{C(sat)}/\beta$).
\item \textbf{Vérifier l'actionneur} : Compatibilité tension/courant avec la sortie traitement. Nécessité d'un relais (isolement, puissance) ?
\item \textbf{Rédiger la chaîne fonctionnelle} : Flèches entre boîtes, grandeurs d'entrée/sortie de chaque boîte.
\end{enumerate}
\end{methode}

\newpage

\coursec{Fiche bilan}

\begin{popfigure}
\begin{tikzpicture}[mindmap, grow cyclic, every node/.style={concept, text=white, font=\small}, concept color=popblue,
level 1/.append style={level distance=3.4cm, sibling angle=51},
level 2/.append style={level distance=2.1cm, sibling angle=45, font=\scriptsize}]
\node[concept, font=\small\bfseries] {Dipôles commandés, capteurs et chaînes électroniques}
  child[concept color=poporange] { node {Dipôle commandé}
    child { node {grandeur de commande} }
    child { node {manuel ou électrique} }
  }
  child[concept color=popgreen] { node {Rhéostat}
    child { node {$U = RI$} }
    child { node {potentiomètre} }
  }
  child[concept color=poppurple] { node {Capteurs}
    child { node {CTN, KTY : température} }
    child { node {LDR : lumière} }
    child { node {antenne, micro : ondes} }
  }
  child[concept color=popgold] { node {Transistor}
    child { node {bloqué / saturé} }
    child { node {$I_C = \beta I_B$} }
  }
  child[concept color=poppink] { node {Relais et VDR}
    child { node {isolement des circuits} }
  }
  child[concept color=popblue!70] { node {Chaîne électronique}
    child { node {acquérir} }
    child { node {traiter} }
    child { node {commander} }
  }
  child[concept color=popdark] { node {Linéarité}
    child { node {$R = f(\theta)$ droite ?} }
  };
\end{tikzpicture}
\legende{Carte mentale de la leçon : les sept idées forces à retenir.}
\end{popfigure}

\begin{aretenir}[L'essentiel de la leçon]
\textbf{Définitions clés}
\begin{itemize}
  \item \cle{Dipôle commandé} : dipôle dont la grandeur électrique de sortie (résistance, tension, courant) dépend d'une grandeur de commande (mécanique, électrique ou physique : température, lumière, pression, onde).
  \item \cle{Capteur} : dipôle commandé dont la grandeur de commande est une grandeur \emph{physique non électrique} ; il la convertit en signal électrique : c'est le principe du captage.
  \item \cle{Chaîne électronique} : succession de trois fonctions : \textbf{acquérir} (capteur) $\rightarrow$ \textbf{traiter} (montage électronique : transistor, amplificateur) $\rightarrow$ \textbf{commander} (actionneur : lampe, sirène, moteur, écouteur).
\end{itemize}

\textbf{Formules et résultats à connaître par cœur}
\begin{center}
\begin{tabularx}{\linewidth}{|l|X|X|}
\hline
\rowcolor{popblueL} \textbf{Dipôle} & \textbf{Loi / caractéristique} & \textbf{Comportement} \\
\hline
Rhéostat & $U = R\,I$ avec $R$ réglable ; montage potentiométrique : $U_s = \dfrac{R_2}{R_1+R_2}\,E$ & $U=f(I,R)$ : faisceau de droites passant par l'origine, de pente $R$ \\
\hline
Thermistance CTN & $R$ diminue quand $\theta$ augmente & capteur \textbf{non linéaire} (courbe décroissante) \\
\hline
KTY & $R$ augmente quasi linéairement avec $\theta$ & capteur quasi \textbf{linéaire} : $R = a\theta + b$ \\
\hline
Photorésistance (LDR) & $R$ diminue quand l'éclairement $\Phi$ augmente & non linéaire ; caractéristique $U=f(I)$ dépend de $\Phi$ \\
\hline
Diode & $I \approx 0$ si $U < U_s$ (seuil $\approx \num{0,6}$ V) ; passante sinon & conductrice dans un seul sens \\
\hline
Transistor (commutation) & $I_C = \beta\, I_B$ (régime linéaire) ; $V_{CE} \approx 0$ en saturation & bloqué ($I_B=0$ : interrupteur ouvert) ou saturé (interrupteur fermé) \\
\hline
Relais électromagnétique & bobine alimentée $\Rightarrow$ contact bascule & commande électrique avec isolement des circuits \\
\hline
\end{tabularx}
\end{center}

\textbf{Méthodes express}
\begin{itemize}
  \item \textbf{Capteur linéaire ou non ?} Tracer $R=f(\theta)$ ou $U=f(I)$ : si la courbe est une droite, le capteur est linéaire (réponse proportionnelle ou affine) ; sinon il est non linéaire.
  \item \textbf{Analyser une chaîne} : identifier les trois blocs (capteur $\rightarrow$ traitement $\rightarrow$ actionneur), puis préciser la grandeur physique d'entrée et l'action de sortie.
  \item \textbf{Transistor en commutation} : comparer $I_B$ au courant de saturation $I_{B,\mathrm{sat}} = \dfrac{E}{\beta R_C}$ ; si $I_B \geq I_{B,\mathrm{sat}}$, le transistor est saturé et l'actionneur est alimenté.
  \item \textbf{Pont diviseur avec capteur} : la tension aux bornes du capteur $U = \dfrac{R_{\text{capteur}}}{R_{\text{capteur}}+R}\,E$ varie avec la grandeur physique : c'est le signal traité par l'étage suivant.
\end{itemize}

\textbf{Exemples types à savoir redécrire}
\begin{itemize}
  \item Sirène commandée par photorésistance et transistor (alarme crépusculaire) : la nuit, $R_{\text{LDR}}$ augmente, $I_B$ devient suffisant, le transistor se sature et la sirène sonne.
  \item Indicateur de niveau d'eau par électrodes et transistor : l'eau, conductrice, referme le circuit de base entre les électrodes et déclenche l'allumage d'une lampe.
  \item Poste radio : antenne (capteur d'ondes électromagnétiques) $\rightarrow$ amplification $\rightarrow$ écouteur (actionneur : membrane mise en vibration par une bobine mobile).
\end{itemize}
\end{aretenir}

\begin{attention}[Pièges fréquents au Bac]
\begin{itemize}
  \item Ne pas confondre \textbf{CTN} ($R$ diminue quand $\theta$ augmente) et \textbf{KTY} ($R$ augmente avec $\theta$) : les deux sont des capteurs de température, mais de comportements opposés.
  \item La formule du pont diviseur $U_s = \dfrac{R_2}{R_1+R_2}E$ suppose que la sortie ne débite \textbf{aucun courant} (sortie à vide ou sur une entrée de très grande résistance).
  \item En saturation, la relation $I_C = \beta I_B$ \textbf{ne s'applique plus} : le courant collecteur est fixé par le circuit extérieur, $I_{C,\mathrm{sat}} \approx \dfrac{E}{R_C}$.
  \item Une caractéristique $U=f(I)$ qui n'est pas une droite ne signifie pas que le dipôle est défectueux : c'est la signature d'un dipôle \textbf{non linéaire} (diode, VDR, thermistance).
  \item Toujours préciser le \textbf{sens de branchement} d'une diode : elle ne conduit que si sa tension anode--cathode dépasse le seuil $U_s \approx \SI{0,6}{V}$ dans le sens passant.
\end{itemize}
\end{attention}

\begin{aretenir}[Checklist avant le Bac]
\begin{itemize}
  \item $\square$ Je sais définir un dipôle commandé et un capteur, et citer la grandeur de commande de chacun.
  \item $\square$ Je sais tracer et exploiter la caractéristique $U = f(I,R)$ d'un rhéostat (faisceau de droites de pente $R$).
  \item $\square$ Je connais le montage potentiométrique et la formule $U_s = \dfrac{R_2}{R_1+R_2}E$.
  \item $\square$ Je sais que la résistance d'une CTN diminue avec la température et que celle d'une LDR diminue avec l'éclairement.
  \item $\square$ Je sais distinguer un capteur linéaire d'un capteur non linéaire à partir de sa courbe caractéristique.
  \item $\square$ Je connais le rôle de la diode (seuil $U_s$, un seul sens de conduction) et de la VDR (résistance qui diminue quand la tension augmente).
  \item $\square$ Je sais expliquer le fonctionnement du relais électromagnétique (bobine, contact, isolement entre circuit de commande et circuit de puissance).
  \item $\square$ Je maîtrise le transistor en commutation : bloqué si $I_B=0$, saturé si $I_B \geq I_{B,\mathrm{sat}}$, avec $I_C = \beta I_B$ en régime linéaire.
  \item $\square$ Je sais décrire les trois blocs d'une chaîne électronique : acquérir, traiter, commander.
  \item $\square$ Je sais expliquer la commande d'une sirène par photorésistance et transistor.
  \item $\square$ Je sais expliquer l'indicateur de niveau d'eau à électrodes et le rôle de l'antenne et de l'écouteur dans un poste radio.
  \item $\square$ Je vérifie toujours les unités (résistances en $\Omega$, températures en \textdegree C, tensions en V, courants en A) et je justifie chaque étape par une loi (Ohm, mailles, nœuds).
\end{itemize}
\end{aretenir}

\findecours

\end{document}
